% GENERATED FILE -- DO NOT EDIT.
%
% Produced by harness/assemble_contract.py from module.yaml and claims/*.tex.
% Edit the claim files and regenerate; `--check` fails if this file is stale.
\documentclass[11pt,a4paper]{article}
\pdftrailerid{888abf6d16980b033c67c25a98bb8569}
%% --- 1. Mathematics (order matters: physics after braket wins \ket, \tr) ---
\usepackage{amsmath}
\usepackage{amssymb}
\usepackage{amsfonts}
\usepackage{amsthm}
\usepackage{mathtools}
\usepackage{dsfont}
\usepackage{braket}
\usepackage{physics}
%% --- 2. Page furniture ------------------------------------------------------
\usepackage{xcolor}
\usepackage{array}    % >{...} column prefixes in the wide tables of some claims
\usepackage{enumitem}
\usepackage{graphicx}
\usepackage{booktabs}
\usepackage{longtable}
\usepackage{xltabular}              % front- and back-matter tables across pages
\usepackage{needspace}
\usepackage{xurl}
\usepackage{tcolorbox}
\tcbuselibrary{skins,breakable}
%% --- 3. Links and navigation (hyperref late; bookmark, refcount after it) ---
\usepackage{xr-hyper}   % before hyperref: the contract's labels
\usepackage{hyperref}
\usepackage{bookmark}
\usepackage{refcount}
%% --- 4. Colours: every distinction also has a shape (rule, caps, size, weight)
\definecolor{nnGrey}{gray}{0.40}      % non-normative text
\definecolor{metaGrey}{gray}{0.30}    % the generated meta line
\definecolor{techGrey}{gray}{0.45}    % heading of a technical claim
\definecolor{mainRule}{RGB}{0,110,140}
\definecolor{reqAccent}{RGB}{165,70,0}
\definecolor{linkBlue}{RGB}{0,62,150}
\definecolor{extLink}{RGB}{120,40,140}  % links into another PDF
\hypersetup{colorlinks=true, linkcolor=linkBlue, citecolor=linkBlue,
  urlcolor=linkBlue, filecolor=extLink, bookmarksnumbered=true,
  bookmarksopen=true, bookmarksopenlevel=1, pdfdisplaydoctitle=true}
% A link into another PDF (\href to a .pdf, \hyperref{file}..., xr-hyper) is a
% relative URI, <file>#nameddest=<destination>, in the file-link colour:
% browser viewers (Edge, Chrome) ignore hyperref's default GoToR action.
\makeatletter
\renewcommand*{\hyper@linkfile}[3]{% anchor text, file name, destination
  \begingroup
    \ifx\\#3\\\def\CQ@target{#2}\else\def\CQ@target{#2\hyper@hash nameddest=#3}\fi
    \let\@urlcolor\@filecolor
    \hyper@linkurl{#1}{\CQ@target}%
  \endgroup}
\makeatother
%% --- 5. The mapped Unicode set (harness/contract_navigation.py, UNICODE_TEX)
\DeclareUnicodeCharacter{03B1}{\ensuremath{\alpha}}
\DeclareUnicodeCharacter{03B2}{\ensuremath{\beta}}
\DeclareUnicodeCharacter{03B3}{\ensuremath{\gamma}}
\DeclareUnicodeCharacter{03B4}{\ensuremath{\delta}}
\DeclareUnicodeCharacter{03B5}{\ensuremath{\varepsilon}}
\DeclareUnicodeCharacter{03F5}{\ensuremath{\epsilon}}
\DeclareUnicodeCharacter{03B6}{\ensuremath{\zeta}}
\DeclareUnicodeCharacter{03B7}{\ensuremath{\eta}}
\DeclareUnicodeCharacter{03B8}{\ensuremath{\theta}}
\DeclareUnicodeCharacter{03D1}{\ensuremath{\vartheta}}
\DeclareUnicodeCharacter{03B9}{\ensuremath{\iota}}
\DeclareUnicodeCharacter{03BA}{\ensuremath{\kappa}}
\DeclareUnicodeCharacter{03BB}{\ensuremath{\lambda}}
\DeclareUnicodeCharacter{03BC}{\ensuremath{\mu}}
\DeclareUnicodeCharacter{03BD}{\ensuremath{\nu}}
\DeclareUnicodeCharacter{03BE}{\ensuremath{\xi}}
\DeclareUnicodeCharacter{03C0}{\ensuremath{\pi}}
\DeclareUnicodeCharacter{03C1}{\ensuremath{\rho}}
\DeclareUnicodeCharacter{03F1}{\ensuremath{\varrho}}
\DeclareUnicodeCharacter{03C3}{\ensuremath{\sigma}}
\DeclareUnicodeCharacter{03C2}{\ensuremath{\varsigma}}
\DeclareUnicodeCharacter{03C4}{\ensuremath{\tau}}
\DeclareUnicodeCharacter{03C5}{\ensuremath{\upsilon}}
\DeclareUnicodeCharacter{03C6}{\ensuremath{\varphi}}
\DeclareUnicodeCharacter{03D5}{\ensuremath{\phi}}
\DeclareUnicodeCharacter{03C7}{\ensuremath{\chi}}
\DeclareUnicodeCharacter{03C8}{\ensuremath{\psi}}
\DeclareUnicodeCharacter{03C9}{\ensuremath{\omega}}
\DeclareUnicodeCharacter{0393}{\ensuremath{\Gamma}}
\DeclareUnicodeCharacter{0394}{\ensuremath{\Delta}}
\DeclareUnicodeCharacter{0398}{\ensuremath{\Theta}}
\DeclareUnicodeCharacter{039B}{\ensuremath{\Lambda}}
\DeclareUnicodeCharacter{039E}{\ensuremath{\Xi}}
\DeclareUnicodeCharacter{03A0}{\ensuremath{\Pi}}
\DeclareUnicodeCharacter{03A3}{\ensuremath{\Sigma}}
\DeclareUnicodeCharacter{03A5}{\ensuremath{\Upsilon}}
\DeclareUnicodeCharacter{03A6}{\ensuremath{\Phi}}
\DeclareUnicodeCharacter{03A8}{\ensuremath{\Psi}}
\DeclareUnicodeCharacter{03A9}{\ensuremath{\Omega}}
\DeclareUnicodeCharacter{00E1}{\'a}
\DeclareUnicodeCharacter{00E0}{\`a}
\DeclareUnicodeCharacter{00E2}{\^a}
\DeclareUnicodeCharacter{00E4}{\"a}
\DeclareUnicodeCharacter{00E5}{\aa{}}
\DeclareUnicodeCharacter{00E9}{\'e}
\DeclareUnicodeCharacter{00E8}{\`e}
\DeclareUnicodeCharacter{00EA}{\^e}
\DeclareUnicodeCharacter{00EB}{\"e}
\DeclareUnicodeCharacter{00ED}{\'{\i}}
\DeclareUnicodeCharacter{00EF}{\"{\i}}
\DeclareUnicodeCharacter{00F3}{\'o}
\DeclareUnicodeCharacter{00F2}{\`o}
\DeclareUnicodeCharacter{00F4}{\^o}
\DeclareUnicodeCharacter{00F6}{\"o}
\DeclareUnicodeCharacter{00F8}{\o{}}
\DeclareUnicodeCharacter{00FA}{\'u}
\DeclareUnicodeCharacter{00F9}{\`u}
\DeclareUnicodeCharacter{00FB}{\^u}
\DeclareUnicodeCharacter{00FC}{\"u}
\DeclareUnicodeCharacter{00F1}{\~n}
\DeclareUnicodeCharacter{00E7}{\c{c}}
\DeclareUnicodeCharacter{00DF}{\ss{}}
\DeclareUnicodeCharacter{00C4}{\"A}
\DeclareUnicodeCharacter{00C5}{\AA{}}
\DeclareUnicodeCharacter{00C9}{\'E}
\DeclareUnicodeCharacter{00D6}{\"O}
\DeclareUnicodeCharacter{00DC}{\"U}
\DeclareUnicodeCharacter{2013}{\textendash{}}
\DeclareUnicodeCharacter{2014}{\textemdash{}}
\DeclareUnicodeCharacter{2018}{`}
\DeclareUnicodeCharacter{2019}{'}
\DeclareUnicodeCharacter{201C}{``}
\DeclareUnicodeCharacter{201D}{''}
\DeclareUnicodeCharacter{2026}{\ldots{}}
\DeclareUnicodeCharacter{00A0}{~}
\DeclareUnicodeCharacter{00A7}{\S{}}
\DeclareUnicodeCharacter{2264}{\ensuremath{\le}}
\DeclareUnicodeCharacter{2265}{\ensuremath{\ge}}
\DeclareUnicodeCharacter{2260}{\ensuremath{\neq}}
\DeclareUnicodeCharacter{2248}{\ensuremath{\approx}}
\DeclareUnicodeCharacter{2261}{\ensuremath{\equiv}}
\DeclareUnicodeCharacter{223C}{\ensuremath{\sim}}
\DeclareUnicodeCharacter{2243}{\ensuremath{\simeq}}
\DeclareUnicodeCharacter{2245}{\ensuremath{\cong}}
\DeclareUnicodeCharacter{221D}{\ensuremath{\propto}}
\DeclareUnicodeCharacter{226A}{\ensuremath{\ll}}
\DeclareUnicodeCharacter{226B}{\ensuremath{\gg}}
\DeclareUnicodeCharacter{2192}{\ensuremath{\to}}
\DeclareUnicodeCharacter{2190}{\ensuremath{\leftarrow}}
\DeclareUnicodeCharacter{2194}{\ensuremath{\leftrightarrow}}
\DeclareUnicodeCharacter{21D2}{\ensuremath{\Rightarrow}}
\DeclareUnicodeCharacter{21D0}{\ensuremath{\Leftarrow}}
\DeclareUnicodeCharacter{21D4}{\ensuremath{\Leftrightarrow}}
\DeclareUnicodeCharacter{21A6}{\ensuremath{\mapsto}}
\DeclareUnicodeCharacter{2208}{\ensuremath{\in}}
\DeclareUnicodeCharacter{2209}{\ensuremath{\notin}}
\DeclareUnicodeCharacter{2286}{\ensuremath{\subseteq}}
\DeclareUnicodeCharacter{2282}{\ensuremath{\subset}}
\DeclareUnicodeCharacter{2287}{\ensuremath{\supseteq}}
\DeclareUnicodeCharacter{2283}{\ensuremath{\supset}}
\DeclareUnicodeCharacter{2229}{\ensuremath{\cap}}
\DeclareUnicodeCharacter{222A}{\ensuremath{\cup}}
\DeclareUnicodeCharacter{2216}{\ensuremath{\setminus}}
\DeclareUnicodeCharacter{2205}{\ensuremath{\emptyset}}
\DeclareUnicodeCharacter{2200}{\ensuremath{\forall}}
\DeclareUnicodeCharacter{2203}{\ensuremath{\exists}}
\DeclareUnicodeCharacter{2115}{\ensuremath{\mathbb{N}}}
\DeclareUnicodeCharacter{2124}{\ensuremath{\mathbb{Z}}}
\DeclareUnicodeCharacter{211A}{\ensuremath{\mathbb{Q}}}
\DeclareUnicodeCharacter{211D}{\ensuremath{\mathbb{R}}}
\DeclareUnicodeCharacter{2102}{\ensuremath{\mathbb{C}}}
\DeclareUnicodeCharacter{00B7}{\ensuremath{\cdot}}
\DeclareUnicodeCharacter{00D7}{\ensuremath{\times}}
\DeclareUnicodeCharacter{2212}{\ensuremath{-}}
\DeclareUnicodeCharacter{00B1}{\ensuremath{\pm}}
\DeclareUnicodeCharacter{2213}{\ensuremath{\mp}}
\DeclareUnicodeCharacter{221A}{\ensuremath{\surd}}
\DeclareUnicodeCharacter{2211}{\ensuremath{\sum}}
\DeclareUnicodeCharacter{220F}{\ensuremath{\prod}}
\DeclareUnicodeCharacter{222B}{\ensuremath{\int}}
\DeclareUnicodeCharacter{221E}{\ensuremath{\infty}}
\DeclareUnicodeCharacter{2202}{\ensuremath{\partial}}
\DeclareUnicodeCharacter{2207}{\ensuremath{\nabla}}
\DeclareUnicodeCharacter{2297}{\ensuremath{\otimes}}
\DeclareUnicodeCharacter{2295}{\ensuremath{\oplus}}
\DeclareUnicodeCharacter{2218}{\ensuremath{\circ}}
\DeclareUnicodeCharacter{2308}{\ensuremath{\lceil}}
\DeclareUnicodeCharacter{2309}{\ensuremath{\rceil}}
\DeclareUnicodeCharacter{230A}{\ensuremath{\lfloor}}
\DeclareUnicodeCharacter{230B}{\ensuremath{\rfloor}}
\DeclareUnicodeCharacter{27E8}{\ensuremath{\langle}}
\DeclareUnicodeCharacter{27E9}{\ensuremath{\rangle}}
\DeclareUnicodeCharacter{2016}{\ensuremath{\|}}
\DeclareUnicodeCharacter{2113}{\ensuremath{\ell}}
\DeclareUnicodeCharacter{00B0}{\ensuremath{^{\circ}}}
\DeclareUnicodeCharacter{00B2}{\ensuremath{^{2}}}
\DeclareUnicodeCharacter{00B3}{\ensuremath{^{3}}}
\DeclareUnicodeCharacter{00B9}{\ensuremath{^{1}}}
\DeclareUnicodeCharacter{2080}{\ensuremath{_{0}}}
\DeclareUnicodeCharacter{2081}{\ensuremath{_{1}}}
\DeclareUnicodeCharacter{2082}{\ensuremath{_{2}}}
\DeclareUnicodeCharacter{207F}{\ensuremath{^{n}}}
\DeclareUnicodeCharacter{2096}{\ensuremath{_{k}}}
%% --- 6. Numbering: one counter per section -----------------------
\numberwithin{equation}{section}
\setcounter{tocdepth}{1}
%% --- 7. The project macro set, then the claim environments -----------------
% spec/macros.tex, read relative to this module's directory (latexmk -cd).
% =============================================================================
%  spec/macros.tex -- the one macro set every generated contract inputs.
%
%  Two parts: part A, the macros spec/notation.md catalogues, and part B,
%  the macros local to one module, grouped by module.
%  harness/assemble_contract.py emits % =============================================================================
%  spec/macros.tex -- the one macro set every generated contract inputs.
%
%  Two parts: part A, the macros spec/notation.md catalogues, and part B,
%  the macros local to one module, grouped by module.
%  harness/assemble_contract.py emits % =============================================================================
%  spec/macros.tex -- the one macro set every generated contract inputs.
%
%  Two parts: part A, the macros spec/notation.md catalogues, and part B,
%  the macros local to one module, grouped by module.
%  harness/assemble_contract.py emits \input{../../macros} for it, so the
%  file is read from spec/modules/<id>/.
%
%  RULES
%    * A macro that spec/notation.md catalogues is a verbatim copy of it
%      (part A). Changing one is a notation change: bump spec/notation.md.
%    * Part B holds the module-local macros. A macro defined by two modules
%      with the same expansion appears once, under the first module (in the
%      order below) that defines it. Two macros with the same name and
%      DIFFERENT expansions are both kept under a module prefix, and the
%      claims write the prefixed names:
%        \Chain  -> \redChain (reductions, \mathfrak{C}: the reduction chain)
%                   \pcChain  (protocol-compiler, \mathfrak{R}: the two-segment
%                             chain; \Chsrc and \Chrec are built on it)
%        \Nmsm   -> \redNmsm  (reductions, \mathcal{N}: a POVM measurement
%                             channel)
%                   \pcNmsm   (\mathcal{N}_{\mathsf A}: Alice's measurement
%                             channel, written by
%                             transmitter/thm:certificate-sufficiency,
%                             transmitter/rem:xreg and
%                             reductions/lem:source-transfer)
%    * Deliberately omitted: a \provenance{tag}{details} macro. Provenance
%      is a header field of every claim file (spec/claim-format.md section
%      3); a claim body that calls such a macro fails to compile, which is
%      the intended signal.
%
%  REQUIRED PACKAGES, loaded by the assembler's preamble before this file:
%    amsmath, amssymb, amsfonts, mathtools, dsfont, braket, physics, xcolor
%    (braket before physics, so that physics wins for \ket, \bra, \tr, ...).
%
%  This file contains only definitions: no \usepackage, no document body.
% =============================================================================


% =============================================================================
%  PART A -- notation macros, the ones spec/notation.md catalogues, each a
%  verbatim copy of its catalogue row; grouped under numbered headings.
% =============================================================================

% -----------------------------------------------------------------------------
% 1. State spaces and operator sets (a "system label" -> spaces of operators)
% -----------------------------------------------------------------------------
\newcommand{\Hs}[1]{\mathcal{H}(#1)}        % Hilbert space of a system label
\newcommand{\Lin}[1]{\mathcal{L}(#1)}       % bounded linear operators on H(#1)
\newcommand{\Linn}[2]{\mathcal{L}(#1,#2)}   % linear maps H(#1) -> H(#2)
\newcommand{\Pos}[1]{\mathrm{Pos}(#1)}      % positive semidefinite operators on H(#1)
\newcommand{\Density}{\mathrm{D}}           % density operators: \Density(A) = D(A)
\newcommand{\Seq}[1]{S_{=}(#1)}             % normalized states (trace = 1)
\newcommand{\Sleq}[1]{S_{\leq}(#1)}         % subnormalized states (0 <= trace <= 1)
\newcommand{\id}{\mathds{1}}                % identity operator
\newcommand{\idon}[1]{\mathds{1}_{#1}}      % identity on a named system, e.g. \idon{C}

% DEPRECATED: applies \mathcal to both arguments itself, so it takes bare
% letters, not system labels. Prefer \Linn; do not use in new contracts.
\providecommand{\Lop}[2]{\mathcal{L}(\mathcal{#1}, \mathcal{#2})}


% -----------------------------------------------------------------------------
% 2. Composite systems
% -----------------------------------------------------------------------------
\newcommand{\dsum}{\oplus}                  % direct sum of system labels
% Tensor product: use \otimes directly.
% n-fold direct sum of an operator space: \Lin{A}^{\oplus n}


% -----------------------------------------------------------------------------
% 3. Events, conditioning and classical registers (for STATES)
%    Convention:  \wedge  -> subnormalized "and-event" component, e.g. \rho_{A\wedge x}
%                 |        -> normalized conditional state,        e.g. \rho_{A|x}
% (Write these directly in subscripts; no extra macros needed.)
% -----------------------------------------------------------------------------


% -----------------------------------------------------------------------------
% 4. Photon-number subspaces / direct-sum blocks
%    Convention:  operators carry the restriction "on top", in brackets;
%                 states carry it as a \wedge subscript (see Sec. 3).
% -----------------------------------------------------------------------------
\newcommand{\Pnum}[1]{P_{#1}}               % photon-number projector (system from context)
\newcommand{\Pnumon}[2]{P_{#1}^{#2}}        % with explicit system superscript, e.g. P_n^C
\newcommand{\blk}[2]{#1^{(#2)}}             % operator block, e.g. \blk{M}{\leqN} = M^{(<=N)}
\newcommand{\leqN}{\leq N}                  % use as \blk{M}{\leqN}
\newcommand{\gtN}{> N}                      % use as \blk{M}{\gtN}
\newcommand{\eqn}{= n}                      % use as \blk{M}{\eqn}


% -----------------------------------------------------------------------------
% 5. System labels used in the protocol
% -----------------------------------------------------------------------------
\newcommand{\Chan}{C}                       % quantum-channel (Eve) input to Bob
\newcommand{\Prb}{P}                        % probe input
\newcommand{\Flag}{{\mathrm{Flg}}}          % flag system (EXISTS ONLY AFTER squashing)
\newcommand{\Loss}{{\mathrm{Lss}}}          % loss system
\newcommand{\Vc}{{\mathrm{Vcm}}}            % vacuum system
\newcommand{\Nout}{N_{\mathrm{out}}}
\newcommand{\Nin}{N_{\mathrm{in}}}


% -----------------------------------------------------------------------------
% 6. Alphabets
% -----------------------------------------------------------------------------
\newcommand{\Yset}{\mathcal{Y}}             % Bob's (coarse-grained) detection alphabet
\newcommand{\Yraw}{\mathcal{Y}_{\mathrm{raw}}}% raw click patterns {0,1}^4
\newcommand{\Yscset}{\mathcal{Y}_{\singlec}}% single-click outcomes
\newcommand{\Ymcset}{\mathcal{Y}_{\mc}}     % multi-click outcomes
\newcommand{\Tset}{\mathcal{T}_{\mathrm{set}}}% private setting set (the probe settings in the optical instance)
\newcommand{\Xset}{\mathcal{X}}             % Alice's alphabet
\newcommand{\Dset}{\{H,V,+,-\}}             % polarization / detector labels
\newcommand{\pass}{\top}                    % pass sifting


% -----------------------------------------------------------------------------
% 7. Mode / outcome abbreviations
% -----------------------------------------------------------------------------
\newcommand{\on}{{\mathrm{on}}}
\newcommand{\off}{{\mathrm{off}}}
\newcommand{\dc}{{\mathrm{dc}}}             % dark count
\newcommand{\sq}{\mathrm{sq}}               % squashed
\newcommand{\mc}{\mathrm{mc}}               % multi-click
\newcommand{\nc}{\mathrm{nc}}               % no-click
\newcommand{\singlec}{\mathrm{sc}}          % single-click
\newcommand{\vac}{\mathrm{vac}}             % vacuum
\newcommand{\kout}{{k_\mathrm{out}}}
\newcommand{\kin}{{k_\mathrm{in}}}


% -----------------------------------------------------------------------------
% 8. Maps, channels and miscellaneous operators
% -----------------------------------------------------------------------------
\newcommand{\Prob}{\mathrm{Prob}}           % probability (use \Prob[\cdot] or \Prob(\cdot))
\newcommand{\lmin}{\lambda_{\min}}          % minimum eigenvalue
\newcommand{\Wflag}{W_{\Flag}}              % flag-space weight
\newcommand{\transpose}{\mathsf{T}}
\newcommand{\innerp}[2]{\left\langle #1 \vert #2 \right\rangle}
\renewcommand{\vec}[1]{\boldsymbol{#1}}     % NB: overrides the LaTeX arrow accent


% -----------------------------------------------------------------------------
% 9. Key-rate numerics
% -----------------------------------------------------------------------------
\newcommand{\cG}{\mathcal{G}}               % protocol map
\newcommand{\cZ}{\mathcal{Z}}               % pinching map
\newcommand{\Amap}{\mathcal{A}_{\mathrm{ann}}} % announcement + measurement map
\newcommand{\Amapnb}{\mathcal{A}_{\mathrm{ann}|\mathrm{nb}}}
\newcommand{\simpl}{}                       % empty suffix for simplified block maps
\newcommand{\acSet}{\mathcal{S}_{\mathrm{acc}}} % acceptance set
\newcommand{\testE}{\mathcal{E}_{\mathrm{test}}} % test-statistics map
\newcommand{\ecLeak}{\mathrm{leak}^{\mathrm{EC}}_{\mathrm{obs}}}
\newcommand{\pgen}{p_{\mathrm{gen}}}
\newcommand{\askeyr}{R_{\infty}}             % asymptotic key rate
\newcommand{\ppass}{p_{\mathrm{pass}}}       % sifting pass probability


% -----------------------------------------------------------------------------
% 10. Finite-size analysis / Renyi entropies
% -----------------------------------------------------------------------------

% --- Round types: test vs. key generation (register T, alphabet \Fset) --------
\newcommand{\tst}{\mathrm{test}}            % test round
\newcommand{\gen}{\mathrm{gen}}             % key-generation round
\newcommand{\Fset}{\mathcal{T}}             % round-type alphabet {test, gen} (not \Tset)
\newcommand{\Chat}{\widehat{C}}             % acceptance-test register
\newcommand{\Chatset}{\widehat{\mathcal{C}}}% its alphabet

% --- Basis / key extraction functions ----------------------------------------
\newcommand{\bA}{b_{A}}                     % Alice's basis-extraction function
\newcommand{\bB}{b_{B}}                     % Bob's basis-extraction function
\newcommand{\kB}{k_{B}}                     % Bob's key-extraction function

% --- Maps --------------------------------------------------------------------
\newcommand{\Mprot}{\mathcal{M}}            % single-round protocol map
\newcommand{\Phimap}{\boldsymbol{\Phi}}     % test-statistics map
\newcommand{\Pdeph}{\mathcal{P}}            % dephasing (pinching) map on \Blreg

% --- Isometries --------------------------------------------------------------
\newcommand{\Vmsm}{V_{\mathrm{msm}}}        % measurement isometry
\newcommand{\Vann}{V_{\mathrm{ann}}}        % announcement isometry
\newcommand{\VI}{V_{I}}                     % Eve's copy of the announcements
\newcommand{\VChat}{V_{\Chat}}              % acceptance-data isometry
\newcommand{\VR}{V_{R}}                     % sifting + key map -> raw key R
\newcommand{\VZ}{V_{Z}}                     % raw key -> key register Z
\newcommand{\Wprot}{W}                      % full protocol isometry

% --- Renyi entropies ---------------------------------------------------------
% SECURITY-CRITICAL: see spec/notation.md §4. Never change an arrow to
% make an inequality work out; never convert silently between the sandwiched
% (tilde) and Petz (overline) families.
\newcommand{\Hsandup}{\widetilde{H}^{\uparrow}_{\alpha}}   % sandwiched, up
\newcommand{\Hsanddown}{\widetilde{H}^{\downarrow}_{\alpha}}% sandwiched, down
\newcommand{\Hpetzup}{\overline{H}^{\uparrow}_{\alpha}}    % Petz, up
\newcommand{\Hpetzdown}{\overline{H}^{\downarrow}_{\alpha}}% Petz, down
\newcommand{\hup}{h^{\uparrow}_{\alpha}}                   % single-round rate
\newcommand{\hqkd}{h^{\uparrow,\mathrm{QKD}}_{\alpha}}     % single-round QKD rate
\newcommand{\hlow}{\underline{h}_{\alpha}}                 % lower bound on \hqkd
% --- spec/notation.md 3.5 and 4.5 -------------------------------------------
% \hlb  is the first-order linearization lower bound on \hlow: the headline
%       key rate (the `first-order` tier of foundations/def:tiers). NOT \hlow
%       itself.
% \tnum is the instance's one declared numerical tolerance. Its
%       near-collision with the declared marginal \tau_A is disclosed in
%       spec/notation.md 3.5; bare \tau is reserved and denotes neither.
\newcommand{\hlb}{\underline{h}^{\mathrm{FO}}_{\alpha}}    % first-order lower bound on \hlow
\newcommand{\tnum}{\tau_{\mathrm{num}}}                    % declared numerical tolerance


% =============================================================================
%  PART B -- module-local macros, grouped by module. Deprecated aliases,
%  defined so that a text written with them compiles, used in no new contract: \Lop,
%  \testE, \Pdeph, \pass, \ecLeak (spec/notation.md section 3.9).
% =============================================================================

% --- meat-core -----------------------------------------------------------------
\newcommand{\SecSet}{\mathbb{S}}            % sector index set of Q
\newcommand{\Qsys}{Q}                       % declared sector-carrying measured system
\newcommand{\Vset}{\mathcal{V}}             % declared joint outcome alphabet of the Q-POVM
\newcommand{\CAset}{\mathcal{C}_{A}}        % Alice announcement alphabet
\newcommand{\CBset}{\mathcal{C}_{B}}        % Bob announcement alphabet
\newcommand{\Party}{\mathfrak{P}}           % declared party partition
\newcommand{\Oset}[1]{\mathcal{O}_{#1}}     % readout alphabet of party #1
\newcommand{\rd}[1]{\pi_{#1}}               % readout map of party #1
\newcommand{\Sadm}{\mathcal{S}_{A}}         % admissible Alice subspace
\newcommand{\nomarg}{\mathsf{none}}         % "no trusted marginal declared"
\newcommand{\Zset}{\mathcal{Z}_{\perp}}     % raw-key alphabet incl. \perp
\newcommand{\Zval}{\mathcal{Z}}             % key alphabet without \perp
\newcommand{\Etil}{\widetilde{E}}           % single-round purifying register
\newcommand{\Ehat}{\widehat{E}}             % Eve's n-round register
\newcommand{\OmAT}{\Omega_{\mathrm{AT}}}    % acceptance event
\newcommand{\epa}{\varepsilon_{\mathrm{PA}}}
\newcommand{\eev}{\varepsilon_{\mathrm{EV}}}
\newcommand{\lEC}{\lambda_{\mathrm{EC}}}
\newcommand{\Decl}{\mathfrak{D}}            % protocol declaration
\newcommand{\CList}{\mathfrak{L}}           % declared constraint list
\newcommand{\PinS}{\mathcal{P}_{\SecSet}}   % sector pinching on Q
\newcommand{\Feas}{\mathcal{W}}             % feasible set of block states
\newcommand{\Chkey}{\Chatset_{\mathrm{key}}}% key-live acceptance values
\newcommand{\Kann}{\mathcal{K}}             % announcement-class alphabet
\newcommand{\hold}{\mathsf{h}}              % declared key holder (A or B)
\newcommand{\sfD}{D}                        % classical KL divergence (base 2)
\newcommand{\Dtwo}{D}                       % quantum relative entropy (base 2)
\newcommand{\Dsand}{\widetilde{D}_{\alpha}} % sandwiched Renyi divergence
\newcommand{\idmap}{\mathrm{id}}            % identity channel
\newcommand{\Aev}{\mathcal{A}}              % Eve's per-round attack channel
\newcommand{\AttSet}{\mathfrak{A}}          % licensed attack-channel set
\newcommand{\Chacc}{\Chatset_{\mathrm{acc}}}% acceptance support of \acSet
\newcommand{\epsstar}{\varepsilon^{\star}}      % distance to an exact POVM
\newcommand{\dfH}{\Delta_{\mathrm{herm}}}       % Hermiticity defect
\newcommand{\dfP}{\Delta_{\mathrm{pos}}}        % negative part
\newcommand{\dfC}{\Delta_{\mathrm{cpl}}}        % completeness defect
\newcommand{\lmoq}{g}                       % acceptance-side gradient block
\newcommand{\lmoJ}{G}                       % feasible-side gradient block
\newcommand{\bacc}{\beta_{\mathrm{acc}}}    % certified acceptance-side minimum
\newcommand{\Itset}{\mathfrak{I}}           % finite index set of iterates (C22)

% --- reductions ----------------------------------------------------------------
\newcommand{\Bsys}{B}                        % physical (pre-squashing) system
\newcommand{\Tcl}[1]{\mathcal{T}(#1)}        % trace-class operators on H(#1)
\newcommand{\Ins}{\mathcal{I}}               % generalized-instrument branch map
\newcommand{\redNmsm}{\mathcal{N}}              % POVM measurement channel
\newcommand{\Eprot}{\mathcal{E}}             % n-round protocol channel
\newcommand{\Grnd}{\mathsf{G}}               % one round map of a protocol
\newcommand{\Kfin}{\mathsf{K}}               % end-of-protocol processing
\newcommand{\Crit}{\mathsf{sec}}             % value of a security criterion
\newcommand{\Pflag}{\Pi_{\Flag}}             % flag-sector projector on Q
\newcommand{\Pleq}{\Pi_{\leqN}}              % preserved-block projector
\newcommand{\Pgt}{\Pi_{\gtN}}                % above-cutoff projector
\newcommand{\Sph}{S^{\mathrm{ph}}}           % family of physical inputs
\newcommand{\Srch}{S^{\mathrm{rch}}}         % reachable squashed states
\newcommand{\Ssq}{S^{\sq}}                   % squashed admissible set
\newcommand{\Ffss}{F^{\mathrm{FSS}}}         % flag-state-squasher target
\newcommand{\Fwp}{F^{\mathrm{WP}}}           % multi-bound WPFSS target
\newcommand{\Mwp}{M^{\mathrm{WP}}}           % renormalized above-cutoff POVM
\newcommand{\Mdec}{{M'}^{\mathrm{ph}}}       % declared physical family on B
\newcommand{\Ppvm}{P^{\Flag}}                % flag PVM
\newcommand{\Lmp}{{\Lambda_{\mathrm{mp}}}}   % measure-and-prepare channel
\newcommand{\Lwp}{{\Lambda^{\mathrm{WP}}}}   % multi-bound WPFSS channel
\newcommand{\Lfss}{{\Lambda^{\mathrm{FSS}}}} % flag-state-squasher channel
\newcommand{\lmax}{\lambda_{\max}}           % maximum eigenvalue
\newcommand{\Lnoise}{\widetilde{\Lambda}}    % flag-noise channel
\newcommand{\cg}{\varpi}                     % declared coarse-graining map
\newcommand{\MCK}{\mathsf{MC}}               % multiclick raw-pattern set
\newcommand{\redChain}{\mathfrak{C}}            % reduction chain
\newcommand{\Lnk}{\ell}                      % a single typed link
\newcommand{\dsq}{\preceq}                      % the squashing relation
\newcommand{\opn}[1]{\left\lVert #1\right\rVert}           % operator (spectral) norm
\newcommand{\mxn}[1]{\left\lVert #1\right\rVert_{\max}}    % entrywise max modulus
\newcommand{\Fld}{\mathbb{F}}                   % exact ordered substrate field
\newcommand{\Fldi}{\mathbb{F}_{i}}              % its Gaussian extension F + iF
\newcommand{\lat}[1]{\big\lfloor #1\big\rfloor_{\kappa}} % round DOWN onto the lattice

% --- protocol-compiler ---------------------------------------------------------
\newcommand{\Spec}{\mathfrak{S}}             % protocol specification (compiler input)
\newcommand{\pcChain}{\mathfrak{R}}            % reduction chain (two typed segments)
\newcommand{\Chsrc}{\pcChain_{\mathrm{src}}}   % its Alice-side source-map segment
\newcommand{\Chrec}{\pcChain_{\mathrm{rec}}}   % its receiver-side squashing segment
\newcommand{\rvirt}{\rho^{\mathrm{v}}}       % virtual (PRE-source-map) signal state
\newcommand{\Fph}{F^{\mathrm{ph}}}           % physical (pre-reduction) POVM element
\newcommand{\Vph}{\mathcal{V}^{\mathrm{ph}}} % physical outcome alphabet
\newcommand{\Tpub}{\mathcal{T}_{\mathrm{pub}}}% public-choice alphabet
\newcommand{\Uset}{\mathcal{U}}              % residual outcome alphabet
\newcommand{\idxm}{\mathsf{idx}}             % label-to-index map
\newcommand{\Asig}{A'}                       % transmitted (signal) system -- the REAL one
\newcommand{\Avirt}{A''}                     % virtual transmitted system (source descriptor of Xi_1)
\newcommand{\Abar}{\bar{A}}                  % Alice's measured register
\newcommand{\Ahat}{\hat{A}}                  % Alice's shield register
\newcommand{\pcNmsm}{\mathcal{N}_{\mathsf A}}  % Alice's measurement channel (see the RULES block)
\newcommand{\Phip}{\Phi^{+}}                 % maximally entangled vector label
\newcommand{\Smap}{\mathfrak{s}}             % the emitted sector map (fraktur: a compile-time object)

% --- receiver ------------------------------------------------------------------
\newcommand{\Modes}{\mathbb{M}}              % declared ordered mode set
\newcommand{\Inp}{\mathbb{I}}                % declared input modes
\newcommand{\InpC}{\mathbb{I}_{\Chan}}       % channel input modes
\newcommand{\InpS}{\mathbb{I}_{\Prb}}        % probe (source) input modes
\newcommand{\Env}{\mathbb{E}}                % dilation (environment) modes
\newcommand{\Port}[1]{\mathbb{P}_{#1}}       % mode set of detector port #1
\newcommand{\Emd}[1]{E_{#1}}                 % coordinate projector on modes #1
\newcommand{\Gmat}{G}                        % instrument (transfer) matrix
\newcommand{\Fock}{\mathfrak{F}}             % Fock / second-quantization functor
\newcommand{\Symn}[1]{\mathrm{Sym}^{#1}}     % symmetric tensor power
\newcommand{\normord}[1]{{:}\mspace{2mu}#1\mspace{2mu}{:}}   % normal ordering
\newcommand{\rk}{\operatorname{rank}}        % rank (physics.sty owns \rank)
\newcommand{\SCK}{\mathsf{SC}}               % single-click pattern set
\newcommand{\Recv}{\mathfrak{R}}             % receiver instrument tuple
\newcommand{\Cat}{\mathfrak{C}}              % component catalogue
\newcommand{\Grph}{\mathfrak{E}}             % declared element graph
\newcommand{\defic}{\mathbb{L}}              % loss deficiency operator
\newcommand{\Pvac}[1]{\Pi_{\vac}(#1)}        % joint no-click element on ports #1
\newcommand{\pmc}{p_{\mc}}                   % uniform high-photon mc weight
\newcommand{\Bnd}{\beta}                     % the scalar bound function beta(n)
\newcommand{\Rot}{R}                         % planar rotation matrix
\newcommand{\Wplate}{W}                      % retarder (wave-plate) matrix
\newcommand{\Bsp}{B}                         % beamsplitter core matrix

% --- statistics ----------------------------------------------------------------
\newcommand{\honest}{p^{\mathrm{hon}}}          % honest per-round distribution
\newcommand{\ecomp}{\varepsilon_{\mathrm{comp}}}% completeness budget
\newcommand{\simplex}[1]{\Delta(#1)}            % probability simplex
\newcommand{\Chatcon}{\Chatset_{\mathrm{con}}}  % constrained outcomes
\newcommand{\Ireg}[3]{\mathcal{I}_{#1}(#2,#3)}  % regularized incomplete beta
\newcommand{\Ncount}[1]{\mathsf{N}_{#1}}        % per-outcome count
\newcommand{\qlo}{\mathbf{q}^{-}}               % box lower bounds
\newcommand{\qhi}{\mathbf{q}^{+}}               % box upper bounds
 for it, so the
%  file is read from spec/modules/<id>/.
%
%  RULES
%    * A macro that spec/notation.md catalogues is a verbatim copy of it
%      (part A). Changing one is a notation change: bump spec/notation.md.
%    * Part B holds the module-local macros. A macro defined by two modules
%      with the same expansion appears once, under the first module (in the
%      order below) that defines it. Two macros with the same name and
%      DIFFERENT expansions are both kept under a module prefix, and the
%      claims write the prefixed names:
%        \Chain  -> \redChain (reductions, \mathfrak{C}: the reduction chain)
%                   \pcChain  (protocol-compiler, \mathfrak{R}: the two-segment
%                             chain; \Chsrc and \Chrec are built on it)
%        \Nmsm   -> \redNmsm  (reductions, \mathcal{N}: a POVM measurement
%                             channel)
%                   \pcNmsm   (\mathcal{N}_{\mathsf A}: Alice's measurement
%                             channel, written by
%                             transmitter/thm:certificate-sufficiency,
%                             transmitter/rem:xreg and
%                             reductions/lem:source-transfer)
%    * Deliberately omitted: a \provenance{tag}{details} macro. Provenance
%      is a header field of every claim file (spec/claim-format.md section
%      3); a claim body that calls such a macro fails to compile, which is
%      the intended signal.
%
%  REQUIRED PACKAGES, loaded by the assembler's preamble before this file:
%    amsmath, amssymb, amsfonts, mathtools, dsfont, braket, physics, xcolor
%    (braket before physics, so that physics wins for \ket, \bra, \tr, ...).
%
%  This file contains only definitions: no \usepackage, no document body.
% =============================================================================


% =============================================================================
%  PART A -- notation macros, the ones spec/notation.md catalogues, each a
%  verbatim copy of its catalogue row; grouped under numbered headings.
% =============================================================================

% -----------------------------------------------------------------------------
% 1. State spaces and operator sets (a "system label" -> spaces of operators)
% -----------------------------------------------------------------------------
\newcommand{\Hs}[1]{\mathcal{H}(#1)}        % Hilbert space of a system label
\newcommand{\Lin}[1]{\mathcal{L}(#1)}       % bounded linear operators on H(#1)
\newcommand{\Linn}[2]{\mathcal{L}(#1,#2)}   % linear maps H(#1) -> H(#2)
\newcommand{\Pos}[1]{\mathrm{Pos}(#1)}      % positive semidefinite operators on H(#1)
\newcommand{\Density}{\mathrm{D}}           % density operators: \Density(A) = D(A)
\newcommand{\Seq}[1]{S_{=}(#1)}             % normalized states (trace = 1)
\newcommand{\Sleq}[1]{S_{\leq}(#1)}         % subnormalized states (0 <= trace <= 1)
\newcommand{\id}{\mathds{1}}                % identity operator
\newcommand{\idon}[1]{\mathds{1}_{#1}}      % identity on a named system, e.g. \idon{C}

% DEPRECATED: applies \mathcal to both arguments itself, so it takes bare
% letters, not system labels. Prefer \Linn; do not use in new contracts.
\providecommand{\Lop}[2]{\mathcal{L}(\mathcal{#1}, \mathcal{#2})}


% -----------------------------------------------------------------------------
% 2. Composite systems
% -----------------------------------------------------------------------------
\newcommand{\dsum}{\oplus}                  % direct sum of system labels
% Tensor product: use \otimes directly.
% n-fold direct sum of an operator space: \Lin{A}^{\oplus n}


% -----------------------------------------------------------------------------
% 3. Events, conditioning and classical registers (for STATES)
%    Convention:  \wedge  -> subnormalized "and-event" component, e.g. \rho_{A\wedge x}
%                 |        -> normalized conditional state,        e.g. \rho_{A|x}
% (Write these directly in subscripts; no extra macros needed.)
% -----------------------------------------------------------------------------


% -----------------------------------------------------------------------------
% 4. Photon-number subspaces / direct-sum blocks
%    Convention:  operators carry the restriction "on top", in brackets;
%                 states carry it as a \wedge subscript (see Sec. 3).
% -----------------------------------------------------------------------------
\newcommand{\Pnum}[1]{P_{#1}}               % photon-number projector (system from context)
\newcommand{\Pnumon}[2]{P_{#1}^{#2}}        % with explicit system superscript, e.g. P_n^C
\newcommand{\blk}[2]{#1^{(#2)}}             % operator block, e.g. \blk{M}{\leqN} = M^{(<=N)}
\newcommand{\leqN}{\leq N}                  % use as \blk{M}{\leqN}
\newcommand{\gtN}{> N}                      % use as \blk{M}{\gtN}
\newcommand{\eqn}{= n}                      % use as \blk{M}{\eqn}


% -----------------------------------------------------------------------------
% 5. System labels used in the protocol
% -----------------------------------------------------------------------------
\newcommand{\Chan}{C}                       % quantum-channel (Eve) input to Bob
\newcommand{\Prb}{P}                        % probe input
\newcommand{\Flag}{{\mathrm{Flg}}}          % flag system (EXISTS ONLY AFTER squashing)
\newcommand{\Loss}{{\mathrm{Lss}}}          % loss system
\newcommand{\Vc}{{\mathrm{Vcm}}}            % vacuum system
\newcommand{\Nout}{N_{\mathrm{out}}}
\newcommand{\Nin}{N_{\mathrm{in}}}


% -----------------------------------------------------------------------------
% 6. Alphabets
% -----------------------------------------------------------------------------
\newcommand{\Yset}{\mathcal{Y}}             % Bob's (coarse-grained) detection alphabet
\newcommand{\Yraw}{\mathcal{Y}_{\mathrm{raw}}}% raw click patterns {0,1}^4
\newcommand{\Yscset}{\mathcal{Y}_{\singlec}}% single-click outcomes
\newcommand{\Ymcset}{\mathcal{Y}_{\mc}}     % multi-click outcomes
\newcommand{\Tset}{\mathcal{T}_{\mathrm{set}}}% private setting set (the probe settings in the optical instance)
\newcommand{\Xset}{\mathcal{X}}             % Alice's alphabet
\newcommand{\Dset}{\{H,V,+,-\}}             % polarization / detector labels
\newcommand{\pass}{\top}                    % pass sifting


% -----------------------------------------------------------------------------
% 7. Mode / outcome abbreviations
% -----------------------------------------------------------------------------
\newcommand{\on}{{\mathrm{on}}}
\newcommand{\off}{{\mathrm{off}}}
\newcommand{\dc}{{\mathrm{dc}}}             % dark count
\newcommand{\sq}{\mathrm{sq}}               % squashed
\newcommand{\mc}{\mathrm{mc}}               % multi-click
\newcommand{\nc}{\mathrm{nc}}               % no-click
\newcommand{\singlec}{\mathrm{sc}}          % single-click
\newcommand{\vac}{\mathrm{vac}}             % vacuum
\newcommand{\kout}{{k_\mathrm{out}}}
\newcommand{\kin}{{k_\mathrm{in}}}


% -----------------------------------------------------------------------------
% 8. Maps, channels and miscellaneous operators
% -----------------------------------------------------------------------------
\newcommand{\Prob}{\mathrm{Prob}}           % probability (use \Prob[\cdot] or \Prob(\cdot))
\newcommand{\lmin}{\lambda_{\min}}          % minimum eigenvalue
\newcommand{\Wflag}{W_{\Flag}}              % flag-space weight
\newcommand{\transpose}{\mathsf{T}}
\newcommand{\innerp}[2]{\left\langle #1 \vert #2 \right\rangle}
\renewcommand{\vec}[1]{\boldsymbol{#1}}     % NB: overrides the LaTeX arrow accent


% -----------------------------------------------------------------------------
% 9. Key-rate numerics
% -----------------------------------------------------------------------------
\newcommand{\cG}{\mathcal{G}}               % protocol map
\newcommand{\cZ}{\mathcal{Z}}               % pinching map
\newcommand{\Amap}{\mathcal{A}_{\mathrm{ann}}} % announcement + measurement map
\newcommand{\Amapnb}{\mathcal{A}_{\mathrm{ann}|\mathrm{nb}}}
\newcommand{\simpl}{}                       % empty suffix for simplified block maps
\newcommand{\acSet}{\mathcal{S}_{\mathrm{acc}}} % acceptance set
\newcommand{\testE}{\mathcal{E}_{\mathrm{test}}} % test-statistics map
\newcommand{\ecLeak}{\mathrm{leak}^{\mathrm{EC}}_{\mathrm{obs}}}
\newcommand{\pgen}{p_{\mathrm{gen}}}
\newcommand{\askeyr}{R_{\infty}}             % asymptotic key rate
\newcommand{\ppass}{p_{\mathrm{pass}}}       % sifting pass probability


% -----------------------------------------------------------------------------
% 10. Finite-size analysis / Renyi entropies
% -----------------------------------------------------------------------------

% --- Round types: test vs. key generation (register T, alphabet \Fset) --------
\newcommand{\tst}{\mathrm{test}}            % test round
\newcommand{\gen}{\mathrm{gen}}             % key-generation round
\newcommand{\Fset}{\mathcal{T}}             % round-type alphabet {test, gen} (not \Tset)
\newcommand{\Chat}{\widehat{C}}             % acceptance-test register
\newcommand{\Chatset}{\widehat{\mathcal{C}}}% its alphabet

% --- Basis / key extraction functions ----------------------------------------
\newcommand{\bA}{b_{A}}                     % Alice's basis-extraction function
\newcommand{\bB}{b_{B}}                     % Bob's basis-extraction function
\newcommand{\kB}{k_{B}}                     % Bob's key-extraction function

% --- Maps --------------------------------------------------------------------
\newcommand{\Mprot}{\mathcal{M}}            % single-round protocol map
\newcommand{\Phimap}{\boldsymbol{\Phi}}     % test-statistics map
\newcommand{\Pdeph}{\mathcal{P}}            % dephasing (pinching) map on \Blreg

% --- Isometries --------------------------------------------------------------
\newcommand{\Vmsm}{V_{\mathrm{msm}}}        % measurement isometry
\newcommand{\Vann}{V_{\mathrm{ann}}}        % announcement isometry
\newcommand{\VI}{V_{I}}                     % Eve's copy of the announcements
\newcommand{\VChat}{V_{\Chat}}              % acceptance-data isometry
\newcommand{\VR}{V_{R}}                     % sifting + key map -> raw key R
\newcommand{\VZ}{V_{Z}}                     % raw key -> key register Z
\newcommand{\Wprot}{W}                      % full protocol isometry

% --- Renyi entropies ---------------------------------------------------------
% SECURITY-CRITICAL: see spec/notation.md §4. Never change an arrow to
% make an inequality work out; never convert silently between the sandwiched
% (tilde) and Petz (overline) families.
\newcommand{\Hsandup}{\widetilde{H}^{\uparrow}_{\alpha}}   % sandwiched, up
\newcommand{\Hsanddown}{\widetilde{H}^{\downarrow}_{\alpha}}% sandwiched, down
\newcommand{\Hpetzup}{\overline{H}^{\uparrow}_{\alpha}}    % Petz, up
\newcommand{\Hpetzdown}{\overline{H}^{\downarrow}_{\alpha}}% Petz, down
\newcommand{\hup}{h^{\uparrow}_{\alpha}}                   % single-round rate
\newcommand{\hqkd}{h^{\uparrow,\mathrm{QKD}}_{\alpha}}     % single-round QKD rate
\newcommand{\hlow}{\underline{h}_{\alpha}}                 % lower bound on \hqkd
% --- spec/notation.md 3.5 and 4.5 -------------------------------------------
% \hlb  is the first-order linearization lower bound on \hlow: the headline
%       key rate (the `first-order` tier of foundations/def:tiers). NOT \hlow
%       itself.
% \tnum is the instance's one declared numerical tolerance. Its
%       near-collision with the declared marginal \tau_A is disclosed in
%       spec/notation.md 3.5; bare \tau is reserved and denotes neither.
\newcommand{\hlb}{\underline{h}^{\mathrm{FO}}_{\alpha}}    % first-order lower bound on \hlow
\newcommand{\tnum}{\tau_{\mathrm{num}}}                    % declared numerical tolerance


% =============================================================================
%  PART B -- module-local macros, grouped by module. Deprecated aliases,
%  defined so that a text written with them compiles, used in no new contract: \Lop,
%  \testE, \Pdeph, \pass, \ecLeak (spec/notation.md section 3.9).
% =============================================================================

% --- meat-core -----------------------------------------------------------------
\newcommand{\SecSet}{\mathbb{S}}            % sector index set of Q
\newcommand{\Qsys}{Q}                       % declared sector-carrying measured system
\newcommand{\Vset}{\mathcal{V}}             % declared joint outcome alphabet of the Q-POVM
\newcommand{\CAset}{\mathcal{C}_{A}}        % Alice announcement alphabet
\newcommand{\CBset}{\mathcal{C}_{B}}        % Bob announcement alphabet
\newcommand{\Party}{\mathfrak{P}}           % declared party partition
\newcommand{\Oset}[1]{\mathcal{O}_{#1}}     % readout alphabet of party #1
\newcommand{\rd}[1]{\pi_{#1}}               % readout map of party #1
\newcommand{\Sadm}{\mathcal{S}_{A}}         % admissible Alice subspace
\newcommand{\nomarg}{\mathsf{none}}         % "no trusted marginal declared"
\newcommand{\Zset}{\mathcal{Z}_{\perp}}     % raw-key alphabet incl. \perp
\newcommand{\Zval}{\mathcal{Z}}             % key alphabet without \perp
\newcommand{\Etil}{\widetilde{E}}           % single-round purifying register
\newcommand{\Ehat}{\widehat{E}}             % Eve's n-round register
\newcommand{\OmAT}{\Omega_{\mathrm{AT}}}    % acceptance event
\newcommand{\epa}{\varepsilon_{\mathrm{PA}}}
\newcommand{\eev}{\varepsilon_{\mathrm{EV}}}
\newcommand{\lEC}{\lambda_{\mathrm{EC}}}
\newcommand{\Decl}{\mathfrak{D}}            % protocol declaration
\newcommand{\CList}{\mathfrak{L}}           % declared constraint list
\newcommand{\PinS}{\mathcal{P}_{\SecSet}}   % sector pinching on Q
\newcommand{\Feas}{\mathcal{W}}             % feasible set of block states
\newcommand{\Chkey}{\Chatset_{\mathrm{key}}}% key-live acceptance values
\newcommand{\Kann}{\mathcal{K}}             % announcement-class alphabet
\newcommand{\hold}{\mathsf{h}}              % declared key holder (A or B)
\newcommand{\sfD}{D}                        % classical KL divergence (base 2)
\newcommand{\Dtwo}{D}                       % quantum relative entropy (base 2)
\newcommand{\Dsand}{\widetilde{D}_{\alpha}} % sandwiched Renyi divergence
\newcommand{\idmap}{\mathrm{id}}            % identity channel
\newcommand{\Aev}{\mathcal{A}}              % Eve's per-round attack channel
\newcommand{\AttSet}{\mathfrak{A}}          % licensed attack-channel set
\newcommand{\Chacc}{\Chatset_{\mathrm{acc}}}% acceptance support of \acSet
\newcommand{\epsstar}{\varepsilon^{\star}}      % distance to an exact POVM
\newcommand{\dfH}{\Delta_{\mathrm{herm}}}       % Hermiticity defect
\newcommand{\dfP}{\Delta_{\mathrm{pos}}}        % negative part
\newcommand{\dfC}{\Delta_{\mathrm{cpl}}}        % completeness defect
\newcommand{\lmoq}{g}                       % acceptance-side gradient block
\newcommand{\lmoJ}{G}                       % feasible-side gradient block
\newcommand{\bacc}{\beta_{\mathrm{acc}}}    % certified acceptance-side minimum
\newcommand{\Itset}{\mathfrak{I}}           % finite index set of iterates (C22)

% --- reductions ----------------------------------------------------------------
\newcommand{\Bsys}{B}                        % physical (pre-squashing) system
\newcommand{\Tcl}[1]{\mathcal{T}(#1)}        % trace-class operators on H(#1)
\newcommand{\Ins}{\mathcal{I}}               % generalized-instrument branch map
\newcommand{\redNmsm}{\mathcal{N}}              % POVM measurement channel
\newcommand{\Eprot}{\mathcal{E}}             % n-round protocol channel
\newcommand{\Grnd}{\mathsf{G}}               % one round map of a protocol
\newcommand{\Kfin}{\mathsf{K}}               % end-of-protocol processing
\newcommand{\Crit}{\mathsf{sec}}             % value of a security criterion
\newcommand{\Pflag}{\Pi_{\Flag}}             % flag-sector projector on Q
\newcommand{\Pleq}{\Pi_{\leqN}}              % preserved-block projector
\newcommand{\Pgt}{\Pi_{\gtN}}                % above-cutoff projector
\newcommand{\Sph}{S^{\mathrm{ph}}}           % family of physical inputs
\newcommand{\Srch}{S^{\mathrm{rch}}}         % reachable squashed states
\newcommand{\Ssq}{S^{\sq}}                   % squashed admissible set
\newcommand{\Ffss}{F^{\mathrm{FSS}}}         % flag-state-squasher target
\newcommand{\Fwp}{F^{\mathrm{WP}}}           % multi-bound WPFSS target
\newcommand{\Mwp}{M^{\mathrm{WP}}}           % renormalized above-cutoff POVM
\newcommand{\Mdec}{{M'}^{\mathrm{ph}}}       % declared physical family on B
\newcommand{\Ppvm}{P^{\Flag}}                % flag PVM
\newcommand{\Lmp}{{\Lambda_{\mathrm{mp}}}}   % measure-and-prepare channel
\newcommand{\Lwp}{{\Lambda^{\mathrm{WP}}}}   % multi-bound WPFSS channel
\newcommand{\Lfss}{{\Lambda^{\mathrm{FSS}}}} % flag-state-squasher channel
\newcommand{\lmax}{\lambda_{\max}}           % maximum eigenvalue
\newcommand{\Lnoise}{\widetilde{\Lambda}}    % flag-noise channel
\newcommand{\cg}{\varpi}                     % declared coarse-graining map
\newcommand{\MCK}{\mathsf{MC}}               % multiclick raw-pattern set
\newcommand{\redChain}{\mathfrak{C}}            % reduction chain
\newcommand{\Lnk}{\ell}                      % a single typed link
\newcommand{\dsq}{\preceq}                      % the squashing relation
\newcommand{\opn}[1]{\left\lVert #1\right\rVert}           % operator (spectral) norm
\newcommand{\mxn}[1]{\left\lVert #1\right\rVert_{\max}}    % entrywise max modulus
\newcommand{\Fld}{\mathbb{F}}                   % exact ordered substrate field
\newcommand{\Fldi}{\mathbb{F}_{i}}              % its Gaussian extension F + iF
\newcommand{\lat}[1]{\big\lfloor #1\big\rfloor_{\kappa}} % round DOWN onto the lattice

% --- protocol-compiler ---------------------------------------------------------
\newcommand{\Spec}{\mathfrak{S}}             % protocol specification (compiler input)
\newcommand{\pcChain}{\mathfrak{R}}            % reduction chain (two typed segments)
\newcommand{\Chsrc}{\pcChain_{\mathrm{src}}}   % its Alice-side source-map segment
\newcommand{\Chrec}{\pcChain_{\mathrm{rec}}}   % its receiver-side squashing segment
\newcommand{\rvirt}{\rho^{\mathrm{v}}}       % virtual (PRE-source-map) signal state
\newcommand{\Fph}{F^{\mathrm{ph}}}           % physical (pre-reduction) POVM element
\newcommand{\Vph}{\mathcal{V}^{\mathrm{ph}}} % physical outcome alphabet
\newcommand{\Tpub}{\mathcal{T}_{\mathrm{pub}}}% public-choice alphabet
\newcommand{\Uset}{\mathcal{U}}              % residual outcome alphabet
\newcommand{\idxm}{\mathsf{idx}}             % label-to-index map
\newcommand{\Asig}{A'}                       % transmitted (signal) system -- the REAL one
\newcommand{\Avirt}{A''}                     % virtual transmitted system (source descriptor of Xi_1)
\newcommand{\Abar}{\bar{A}}                  % Alice's measured register
\newcommand{\Ahat}{\hat{A}}                  % Alice's shield register
\newcommand{\pcNmsm}{\mathcal{N}_{\mathsf A}}  % Alice's measurement channel (see the RULES block)
\newcommand{\Phip}{\Phi^{+}}                 % maximally entangled vector label
\newcommand{\Smap}{\mathfrak{s}}             % the emitted sector map (fraktur: a compile-time object)

% --- receiver ------------------------------------------------------------------
\newcommand{\Modes}{\mathbb{M}}              % declared ordered mode set
\newcommand{\Inp}{\mathbb{I}}                % declared input modes
\newcommand{\InpC}{\mathbb{I}_{\Chan}}       % channel input modes
\newcommand{\InpS}{\mathbb{I}_{\Prb}}        % probe (source) input modes
\newcommand{\Env}{\mathbb{E}}                % dilation (environment) modes
\newcommand{\Port}[1]{\mathbb{P}_{#1}}       % mode set of detector port #1
\newcommand{\Emd}[1]{E_{#1}}                 % coordinate projector on modes #1
\newcommand{\Gmat}{G}                        % instrument (transfer) matrix
\newcommand{\Fock}{\mathfrak{F}}             % Fock / second-quantization functor
\newcommand{\Symn}[1]{\mathrm{Sym}^{#1}}     % symmetric tensor power
\newcommand{\normord}[1]{{:}\mspace{2mu}#1\mspace{2mu}{:}}   % normal ordering
\newcommand{\rk}{\operatorname{rank}}        % rank (physics.sty owns \rank)
\newcommand{\SCK}{\mathsf{SC}}               % single-click pattern set
\newcommand{\Recv}{\mathfrak{R}}             % receiver instrument tuple
\newcommand{\Cat}{\mathfrak{C}}              % component catalogue
\newcommand{\Grph}{\mathfrak{E}}             % declared element graph
\newcommand{\defic}{\mathbb{L}}              % loss deficiency operator
\newcommand{\Pvac}[1]{\Pi_{\vac}(#1)}        % joint no-click element on ports #1
\newcommand{\pmc}{p_{\mc}}                   % uniform high-photon mc weight
\newcommand{\Bnd}{\beta}                     % the scalar bound function beta(n)
\newcommand{\Rot}{R}                         % planar rotation matrix
\newcommand{\Wplate}{W}                      % retarder (wave-plate) matrix
\newcommand{\Bsp}{B}                         % beamsplitter core matrix

% --- statistics ----------------------------------------------------------------
\newcommand{\honest}{p^{\mathrm{hon}}}          % honest per-round distribution
\newcommand{\ecomp}{\varepsilon_{\mathrm{comp}}}% completeness budget
\newcommand{\simplex}[1]{\Delta(#1)}            % probability simplex
\newcommand{\Chatcon}{\Chatset_{\mathrm{con}}}  % constrained outcomes
\newcommand{\Ireg}[3]{\mathcal{I}_{#1}(#2,#3)}  % regularized incomplete beta
\newcommand{\Ncount}[1]{\mathsf{N}_{#1}}        % per-outcome count
\newcommand{\qlo}{\mathbf{q}^{-}}               % box lower bounds
\newcommand{\qhi}{\mathbf{q}^{+}}               % box upper bounds
 for it, so the
%  file is read from spec/modules/<id>/.
%
%  RULES
%    * A macro that spec/notation.md catalogues is a verbatim copy of it
%      (part A). Changing one is a notation change: bump spec/notation.md.
%    * Part B holds the module-local macros. A macro defined by two modules
%      with the same expansion appears once, under the first module (in the
%      order below) that defines it. Two macros with the same name and
%      DIFFERENT expansions are both kept under a module prefix, and the
%      claims write the prefixed names:
%        \Chain  -> \redChain (reductions, \mathfrak{C}: the reduction chain)
%                   \pcChain  (protocol-compiler, \mathfrak{R}: the two-segment
%                             chain; \Chsrc and \Chrec are built on it)
%        \Nmsm   -> \redNmsm  (reductions, \mathcal{N}: a POVM measurement
%                             channel)
%                   \pcNmsm   (\mathcal{N}_{\mathsf A}: Alice's measurement
%                             channel, written by
%                             transmitter/thm:certificate-sufficiency,
%                             transmitter/rem:xreg and
%                             reductions/lem:source-transfer)
%    * Deliberately omitted: a \provenance{tag}{details} macro. Provenance
%      is a header field of every claim file (spec/claim-format.md section
%      3); a claim body that calls such a macro fails to compile, which is
%      the intended signal.
%
%  REQUIRED PACKAGES, loaded by the assembler's preamble before this file:
%    amsmath, amssymb, amsfonts, mathtools, dsfont, braket, physics, xcolor
%    (braket before physics, so that physics wins for \ket, \bra, \tr, ...).
%
%  This file contains only definitions: no \usepackage, no document body.
% =============================================================================


% =============================================================================
%  PART A -- notation macros, the ones spec/notation.md catalogues, each a
%  verbatim copy of its catalogue row; grouped under numbered headings.
% =============================================================================

% -----------------------------------------------------------------------------
% 1. State spaces and operator sets (a "system label" -> spaces of operators)
% -----------------------------------------------------------------------------
\newcommand{\Hs}[1]{\mathcal{H}(#1)}        % Hilbert space of a system label
\newcommand{\Lin}[1]{\mathcal{L}(#1)}       % bounded linear operators on H(#1)
\newcommand{\Linn}[2]{\mathcal{L}(#1,#2)}   % linear maps H(#1) -> H(#2)
\newcommand{\Pos}[1]{\mathrm{Pos}(#1)}      % positive semidefinite operators on H(#1)
\newcommand{\Density}{\mathrm{D}}           % density operators: \Density(A) = D(A)
\newcommand{\Seq}[1]{S_{=}(#1)}             % normalized states (trace = 1)
\newcommand{\Sleq}[1]{S_{\leq}(#1)}         % subnormalized states (0 <= trace <= 1)
\newcommand{\id}{\mathds{1}}                % identity operator
\newcommand{\idon}[1]{\mathds{1}_{#1}}      % identity on a named system, e.g. \idon{C}

% DEPRECATED: applies \mathcal to both arguments itself, so it takes bare
% letters, not system labels. Prefer \Linn; do not use in new contracts.
\providecommand{\Lop}[2]{\mathcal{L}(\mathcal{#1}, \mathcal{#2})}


% -----------------------------------------------------------------------------
% 2. Composite systems
% -----------------------------------------------------------------------------
\newcommand{\dsum}{\oplus}                  % direct sum of system labels
% Tensor product: use \otimes directly.
% n-fold direct sum of an operator space: \Lin{A}^{\oplus n}


% -----------------------------------------------------------------------------
% 3. Events, conditioning and classical registers (for STATES)
%    Convention:  \wedge  -> subnormalized "and-event" component, e.g. \rho_{A\wedge x}
%                 |        -> normalized conditional state,        e.g. \rho_{A|x}
% (Write these directly in subscripts; no extra macros needed.)
% -----------------------------------------------------------------------------


% -----------------------------------------------------------------------------
% 4. Photon-number subspaces / direct-sum blocks
%    Convention:  operators carry the restriction "on top", in brackets;
%                 states carry it as a \wedge subscript (see Sec. 3).
% -----------------------------------------------------------------------------
\newcommand{\Pnum}[1]{P_{#1}}               % photon-number projector (system from context)
\newcommand{\Pnumon}[2]{P_{#1}^{#2}}        % with explicit system superscript, e.g. P_n^C
\newcommand{\blk}[2]{#1^{(#2)}}             % operator block, e.g. \blk{M}{\leqN} = M^{(<=N)}
\newcommand{\leqN}{\leq N}                  % use as \blk{M}{\leqN}
\newcommand{\gtN}{> N}                      % use as \blk{M}{\gtN}
\newcommand{\eqn}{= n}                      % use as \blk{M}{\eqn}


% -----------------------------------------------------------------------------
% 5. System labels used in the protocol
% -----------------------------------------------------------------------------
\newcommand{\Chan}{C}                       % quantum-channel (Eve) input to Bob
\newcommand{\Prb}{P}                        % probe input
\newcommand{\Flag}{{\mathrm{Flg}}}          % flag system (EXISTS ONLY AFTER squashing)
\newcommand{\Loss}{{\mathrm{Lss}}}          % loss system
\newcommand{\Vc}{{\mathrm{Vcm}}}            % vacuum system
\newcommand{\Nout}{N_{\mathrm{out}}}
\newcommand{\Nin}{N_{\mathrm{in}}}


% -----------------------------------------------------------------------------
% 6. Alphabets
% -----------------------------------------------------------------------------
\newcommand{\Yset}{\mathcal{Y}}             % Bob's (coarse-grained) detection alphabet
\newcommand{\Yraw}{\mathcal{Y}_{\mathrm{raw}}}% raw click patterns {0,1}^4
\newcommand{\Yscset}{\mathcal{Y}_{\singlec}}% single-click outcomes
\newcommand{\Ymcset}{\mathcal{Y}_{\mc}}     % multi-click outcomes
\newcommand{\Tset}{\mathcal{T}_{\mathrm{set}}}% private setting set (the probe settings in the optical instance)
\newcommand{\Xset}{\mathcal{X}}             % Alice's alphabet
\newcommand{\Dset}{\{H,V,+,-\}}             % polarization / detector labels
\newcommand{\pass}{\top}                    % pass sifting


% -----------------------------------------------------------------------------
% 7. Mode / outcome abbreviations
% -----------------------------------------------------------------------------
\newcommand{\on}{{\mathrm{on}}}
\newcommand{\off}{{\mathrm{off}}}
\newcommand{\dc}{{\mathrm{dc}}}             % dark count
\newcommand{\sq}{\mathrm{sq}}               % squashed
\newcommand{\mc}{\mathrm{mc}}               % multi-click
\newcommand{\nc}{\mathrm{nc}}               % no-click
\newcommand{\singlec}{\mathrm{sc}}          % single-click
\newcommand{\vac}{\mathrm{vac}}             % vacuum
\newcommand{\kout}{{k_\mathrm{out}}}
\newcommand{\kin}{{k_\mathrm{in}}}


% -----------------------------------------------------------------------------
% 8. Maps, channels and miscellaneous operators
% -----------------------------------------------------------------------------
\newcommand{\Prob}{\mathrm{Prob}}           % probability (use \Prob[\cdot] or \Prob(\cdot))
\newcommand{\lmin}{\lambda_{\min}}          % minimum eigenvalue
\newcommand{\Wflag}{W_{\Flag}}              % flag-space weight
\newcommand{\transpose}{\mathsf{T}}
\newcommand{\innerp}[2]{\left\langle #1 \vert #2 \right\rangle}
\renewcommand{\vec}[1]{\boldsymbol{#1}}     % NB: overrides the LaTeX arrow accent


% -----------------------------------------------------------------------------
% 9. Key-rate numerics
% -----------------------------------------------------------------------------
\newcommand{\cG}{\mathcal{G}}               % protocol map
\newcommand{\cZ}{\mathcal{Z}}               % pinching map
\newcommand{\Amap}{\mathcal{A}_{\mathrm{ann}}} % announcement + measurement map
\newcommand{\Amapnb}{\mathcal{A}_{\mathrm{ann}|\mathrm{nb}}}
\newcommand{\simpl}{}                       % empty suffix for simplified block maps
\newcommand{\acSet}{\mathcal{S}_{\mathrm{acc}}} % acceptance set
\newcommand{\testE}{\mathcal{E}_{\mathrm{test}}} % test-statistics map
\newcommand{\ecLeak}{\mathrm{leak}^{\mathrm{EC}}_{\mathrm{obs}}}
\newcommand{\pgen}{p_{\mathrm{gen}}}
\newcommand{\askeyr}{R_{\infty}}             % asymptotic key rate
\newcommand{\ppass}{p_{\mathrm{pass}}}       % sifting pass probability


% -----------------------------------------------------------------------------
% 10. Finite-size analysis / Renyi entropies
% -----------------------------------------------------------------------------

% --- Round types: test vs. key generation (register T, alphabet \Fset) --------
\newcommand{\tst}{\mathrm{test}}            % test round
\newcommand{\gen}{\mathrm{gen}}             % key-generation round
\newcommand{\Fset}{\mathcal{T}}             % round-type alphabet {test, gen} (not \Tset)
\newcommand{\Chat}{\widehat{C}}             % acceptance-test register
\newcommand{\Chatset}{\widehat{\mathcal{C}}}% its alphabet

% --- Basis / key extraction functions ----------------------------------------
\newcommand{\bA}{b_{A}}                     % Alice's basis-extraction function
\newcommand{\bB}{b_{B}}                     % Bob's basis-extraction function
\newcommand{\kB}{k_{B}}                     % Bob's key-extraction function

% --- Maps --------------------------------------------------------------------
\newcommand{\Mprot}{\mathcal{M}}            % single-round protocol map
\newcommand{\Phimap}{\boldsymbol{\Phi}}     % test-statistics map
\newcommand{\Pdeph}{\mathcal{P}}            % dephasing (pinching) map on \Blreg

% --- Isometries --------------------------------------------------------------
\newcommand{\Vmsm}{V_{\mathrm{msm}}}        % measurement isometry
\newcommand{\Vann}{V_{\mathrm{ann}}}        % announcement isometry
\newcommand{\VI}{V_{I}}                     % Eve's copy of the announcements
\newcommand{\VChat}{V_{\Chat}}              % acceptance-data isometry
\newcommand{\VR}{V_{R}}                     % sifting + key map -> raw key R
\newcommand{\VZ}{V_{Z}}                     % raw key -> key register Z
\newcommand{\Wprot}{W}                      % full protocol isometry

% --- Renyi entropies ---------------------------------------------------------
% SECURITY-CRITICAL: see spec/notation.md §4. Never change an arrow to
% make an inequality work out; never convert silently between the sandwiched
% (tilde) and Petz (overline) families.
\newcommand{\Hsandup}{\widetilde{H}^{\uparrow}_{\alpha}}   % sandwiched, up
\newcommand{\Hsanddown}{\widetilde{H}^{\downarrow}_{\alpha}}% sandwiched, down
\newcommand{\Hpetzup}{\overline{H}^{\uparrow}_{\alpha}}    % Petz, up
\newcommand{\Hpetzdown}{\overline{H}^{\downarrow}_{\alpha}}% Petz, down
\newcommand{\hup}{h^{\uparrow}_{\alpha}}                   % single-round rate
\newcommand{\hqkd}{h^{\uparrow,\mathrm{QKD}}_{\alpha}}     % single-round QKD rate
\newcommand{\hlow}{\underline{h}_{\alpha}}                 % lower bound on \hqkd
% --- spec/notation.md 3.5 and 4.5 -------------------------------------------
% \hlb  is the first-order linearization lower bound on \hlow: the headline
%       key rate (the `first-order` tier of foundations/def:tiers). NOT \hlow
%       itself.
% \tnum is the instance's one declared numerical tolerance. Its
%       near-collision with the declared marginal \tau_A is disclosed in
%       spec/notation.md 3.5; bare \tau is reserved and denotes neither.
\newcommand{\hlb}{\underline{h}^{\mathrm{FO}}_{\alpha}}    % first-order lower bound on \hlow
\newcommand{\tnum}{\tau_{\mathrm{num}}}                    % declared numerical tolerance


% =============================================================================
%  PART B -- module-local macros, grouped by module. Deprecated aliases,
%  defined so that a text written with them compiles, used in no new contract: \Lop,
%  \testE, \Pdeph, \pass, \ecLeak (spec/notation.md section 3.9).
% =============================================================================

% --- meat-core -----------------------------------------------------------------
\newcommand{\SecSet}{\mathbb{S}}            % sector index set of Q
\newcommand{\Qsys}{Q}                       % declared sector-carrying measured system
\newcommand{\Vset}{\mathcal{V}}             % declared joint outcome alphabet of the Q-POVM
\newcommand{\CAset}{\mathcal{C}_{A}}        % Alice announcement alphabet
\newcommand{\CBset}{\mathcal{C}_{B}}        % Bob announcement alphabet
\newcommand{\Party}{\mathfrak{P}}           % declared party partition
\newcommand{\Oset}[1]{\mathcal{O}_{#1}}     % readout alphabet of party #1
\newcommand{\rd}[1]{\pi_{#1}}               % readout map of party #1
\newcommand{\Sadm}{\mathcal{S}_{A}}         % admissible Alice subspace
\newcommand{\nomarg}{\mathsf{none}}         % "no trusted marginal declared"
\newcommand{\Zset}{\mathcal{Z}_{\perp}}     % raw-key alphabet incl. \perp
\newcommand{\Zval}{\mathcal{Z}}             % key alphabet without \perp
\newcommand{\Etil}{\widetilde{E}}           % single-round purifying register
\newcommand{\Ehat}{\widehat{E}}             % Eve's n-round register
\newcommand{\OmAT}{\Omega_{\mathrm{AT}}}    % acceptance event
\newcommand{\epa}{\varepsilon_{\mathrm{PA}}}
\newcommand{\eev}{\varepsilon_{\mathrm{EV}}}
\newcommand{\lEC}{\lambda_{\mathrm{EC}}}
\newcommand{\Decl}{\mathfrak{D}}            % protocol declaration
\newcommand{\CList}{\mathfrak{L}}           % declared constraint list
\newcommand{\PinS}{\mathcal{P}_{\SecSet}}   % sector pinching on Q
\newcommand{\Feas}{\mathcal{W}}             % feasible set of block states
\newcommand{\Chkey}{\Chatset_{\mathrm{key}}}% key-live acceptance values
\newcommand{\Kann}{\mathcal{K}}             % announcement-class alphabet
\newcommand{\hold}{\mathsf{h}}              % declared key holder (A or B)
\newcommand{\sfD}{D}                        % classical KL divergence (base 2)
\newcommand{\Dtwo}{D}                       % quantum relative entropy (base 2)
\newcommand{\Dsand}{\widetilde{D}_{\alpha}} % sandwiched Renyi divergence
\newcommand{\idmap}{\mathrm{id}}            % identity channel
\newcommand{\Aev}{\mathcal{A}}              % Eve's per-round attack channel
\newcommand{\AttSet}{\mathfrak{A}}          % licensed attack-channel set
\newcommand{\Chacc}{\Chatset_{\mathrm{acc}}}% acceptance support of \acSet
\newcommand{\epsstar}{\varepsilon^{\star}}      % distance to an exact POVM
\newcommand{\dfH}{\Delta_{\mathrm{herm}}}       % Hermiticity defect
\newcommand{\dfP}{\Delta_{\mathrm{pos}}}        % negative part
\newcommand{\dfC}{\Delta_{\mathrm{cpl}}}        % completeness defect
\newcommand{\lmoq}{g}                       % acceptance-side gradient block
\newcommand{\lmoJ}{G}                       % feasible-side gradient block
\newcommand{\bacc}{\beta_{\mathrm{acc}}}    % certified acceptance-side minimum
\newcommand{\Itset}{\mathfrak{I}}           % finite index set of iterates (C22)

% --- reductions ----------------------------------------------------------------
\newcommand{\Bsys}{B}                        % physical (pre-squashing) system
\newcommand{\Tcl}[1]{\mathcal{T}(#1)}        % trace-class operators on H(#1)
\newcommand{\Ins}{\mathcal{I}}               % generalized-instrument branch map
\newcommand{\redNmsm}{\mathcal{N}}              % POVM measurement channel
\newcommand{\Eprot}{\mathcal{E}}             % n-round protocol channel
\newcommand{\Grnd}{\mathsf{G}}               % one round map of a protocol
\newcommand{\Kfin}{\mathsf{K}}               % end-of-protocol processing
\newcommand{\Crit}{\mathsf{sec}}             % value of a security criterion
\newcommand{\Pflag}{\Pi_{\Flag}}             % flag-sector projector on Q
\newcommand{\Pleq}{\Pi_{\leqN}}              % preserved-block projector
\newcommand{\Pgt}{\Pi_{\gtN}}                % above-cutoff projector
\newcommand{\Sph}{S^{\mathrm{ph}}}           % family of physical inputs
\newcommand{\Srch}{S^{\mathrm{rch}}}         % reachable squashed states
\newcommand{\Ssq}{S^{\sq}}                   % squashed admissible set
\newcommand{\Ffss}{F^{\mathrm{FSS}}}         % flag-state-squasher target
\newcommand{\Fwp}{F^{\mathrm{WP}}}           % multi-bound WPFSS target
\newcommand{\Mwp}{M^{\mathrm{WP}}}           % renormalized above-cutoff POVM
\newcommand{\Mdec}{{M'}^{\mathrm{ph}}}       % declared physical family on B
\newcommand{\Ppvm}{P^{\Flag}}                % flag PVM
\newcommand{\Lmp}{{\Lambda_{\mathrm{mp}}}}   % measure-and-prepare channel
\newcommand{\Lwp}{{\Lambda^{\mathrm{WP}}}}   % multi-bound WPFSS channel
\newcommand{\Lfss}{{\Lambda^{\mathrm{FSS}}}} % flag-state-squasher channel
\newcommand{\lmax}{\lambda_{\max}}           % maximum eigenvalue
\newcommand{\Lnoise}{\widetilde{\Lambda}}    % flag-noise channel
\newcommand{\cg}{\varpi}                     % declared coarse-graining map
\newcommand{\MCK}{\mathsf{MC}}               % multiclick raw-pattern set
\newcommand{\redChain}{\mathfrak{C}}            % reduction chain
\newcommand{\Lnk}{\ell}                      % a single typed link
\newcommand{\dsq}{\preceq}                      % the squashing relation
\newcommand{\opn}[1]{\left\lVert #1\right\rVert}           % operator (spectral) norm
\newcommand{\mxn}[1]{\left\lVert #1\right\rVert_{\max}}    % entrywise max modulus
\newcommand{\Fld}{\mathbb{F}}                   % exact ordered substrate field
\newcommand{\Fldi}{\mathbb{F}_{i}}              % its Gaussian extension F + iF
\newcommand{\lat}[1]{\big\lfloor #1\big\rfloor_{\kappa}} % round DOWN onto the lattice

% --- protocol-compiler ---------------------------------------------------------
\newcommand{\Spec}{\mathfrak{S}}             % protocol specification (compiler input)
\newcommand{\pcChain}{\mathfrak{R}}            % reduction chain (two typed segments)
\newcommand{\Chsrc}{\pcChain_{\mathrm{src}}}   % its Alice-side source-map segment
\newcommand{\Chrec}{\pcChain_{\mathrm{rec}}}   % its receiver-side squashing segment
\newcommand{\rvirt}{\rho^{\mathrm{v}}}       % virtual (PRE-source-map) signal state
\newcommand{\Fph}{F^{\mathrm{ph}}}           % physical (pre-reduction) POVM element
\newcommand{\Vph}{\mathcal{V}^{\mathrm{ph}}} % physical outcome alphabet
\newcommand{\Tpub}{\mathcal{T}_{\mathrm{pub}}}% public-choice alphabet
\newcommand{\Uset}{\mathcal{U}}              % residual outcome alphabet
\newcommand{\idxm}{\mathsf{idx}}             % label-to-index map
\newcommand{\Asig}{A'}                       % transmitted (signal) system -- the REAL one
\newcommand{\Avirt}{A''}                     % virtual transmitted system (source descriptor of Xi_1)
\newcommand{\Abar}{\bar{A}}                  % Alice's measured register
\newcommand{\Ahat}{\hat{A}}                  % Alice's shield register
\newcommand{\pcNmsm}{\mathcal{N}_{\mathsf A}}  % Alice's measurement channel (see the RULES block)
\newcommand{\Phip}{\Phi^{+}}                 % maximally entangled vector label
\newcommand{\Smap}{\mathfrak{s}}             % the emitted sector map (fraktur: a compile-time object)

% --- receiver ------------------------------------------------------------------
\newcommand{\Modes}{\mathbb{M}}              % declared ordered mode set
\newcommand{\Inp}{\mathbb{I}}                % declared input modes
\newcommand{\InpC}{\mathbb{I}_{\Chan}}       % channel input modes
\newcommand{\InpS}{\mathbb{I}_{\Prb}}        % probe (source) input modes
\newcommand{\Env}{\mathbb{E}}                % dilation (environment) modes
\newcommand{\Port}[1]{\mathbb{P}_{#1}}       % mode set of detector port #1
\newcommand{\Emd}[1]{E_{#1}}                 % coordinate projector on modes #1
\newcommand{\Gmat}{G}                        % instrument (transfer) matrix
\newcommand{\Fock}{\mathfrak{F}}             % Fock / second-quantization functor
\newcommand{\Symn}[1]{\mathrm{Sym}^{#1}}     % symmetric tensor power
\newcommand{\normord}[1]{{:}\mspace{2mu}#1\mspace{2mu}{:}}   % normal ordering
\newcommand{\rk}{\operatorname{rank}}        % rank (physics.sty owns \rank)
\newcommand{\SCK}{\mathsf{SC}}               % single-click pattern set
\newcommand{\Recv}{\mathfrak{R}}             % receiver instrument tuple
\newcommand{\Cat}{\mathfrak{C}}              % component catalogue
\newcommand{\Grph}{\mathfrak{E}}             % declared element graph
\newcommand{\defic}{\mathbb{L}}              % loss deficiency operator
\newcommand{\Pvac}[1]{\Pi_{\vac}(#1)}        % joint no-click element on ports #1
\newcommand{\pmc}{p_{\mc}}                   % uniform high-photon mc weight
\newcommand{\Bnd}{\beta}                     % the scalar bound function beta(n)
\newcommand{\Rot}{R}                         % planar rotation matrix
\newcommand{\Wplate}{W}                      % retarder (wave-plate) matrix
\newcommand{\Bsp}{B}                         % beamsplitter core matrix

% --- statistics ----------------------------------------------------------------
\newcommand{\honest}{p^{\mathrm{hon}}}          % honest per-round distribution
\newcommand{\ecomp}{\varepsilon_{\mathrm{comp}}}% completeness budget
\newcommand{\simplex}[1]{\Delta(#1)}            % probability simplex
\newcommand{\Chatcon}{\Chatset_{\mathrm{con}}}  % constrained outcomes
\newcommand{\Ireg}[3]{\mathcal{I}_{#1}(#2,#3)}  % regularized incomplete beta
\newcommand{\Ncount}[1]{\mathsf{N}_{#1}}        % per-outcome count
\newcommand{\qlo}{\mathbf{q}^{-}}               % box lower bounds
\newcommand{\qhi}{\mathbf{q}^{+}}               % box upper bounds

\makeatletter
\newcommand*{\ClaimHeadFont}{\bfseries}   % switched by prominence
\newcommand*{\ClaimNoteFont}{\mdseries}
\newcommand*{\ClaimShortId}{}
\newcommand*{\SetShortId}[1]{\renewcommand*{\ClaimShortId}{#1}}
\newcommand*{\RestatedName}{}
% Upright bodies throughout: italic is reserved for non-normative text.
% \@endparpenalty\@M keeps the meta line on the statement's page.
\newtheoremstyle{claim}{\topsep}{\topsep}{\normalfont\@endparpenalty\@M}{}%
  {\ClaimHeadFont}{.}{.5em}{\thmname{#1}\thmnumber{ #2}\thmnote{ {\ClaimNoteFont(#3)}}}
\newtheoremstyle{claimass}{\topsep}{\topsep}{\normalfont\@endparpenalty\@M}{}%
  {\ClaimHeadFont}{.}{.5em}{\thmname{#1}\thmnumber{ #2}%
  \ifx\ClaimShortId\@empty\else\ [\textsf{\ClaimShortId}]\fi%
  \thmnote{ {\ClaimNoteFont(#3)}}}
\newtheoremstyle{claimremark}{.5\topsep}{.5\topsep}{\small\@endparpenalty\@M}{}%
  {\itshape}{.}{.5em}{\thmname{#1}\thmnumber{ #2}\thmnote{ {\upshape(#3)}}}
\newtheoremstyle{claimreq}{\topsep}{\topsep}{\normalfont\@endparpenalty\@M}{}%
  {\color{reqAccent}\scshape}{.}{.5em}%
  {\thmname{#1}\thmnumber{ #2}\thmnote{ {\upshape(#3)}}}
\makeatother
\theoremstyle{claim}
\newtheorem{theorem}{Theorem}[section]
\newtheorem{lemma}[theorem]{Lemma}
\newtheorem{proposition}[theorem]{Proposition}
\newtheorem{corollary}[theorem]{Corollary}
\newtheorem{fact}[theorem]{Fact}
\newtheorem{definition}[theorem]{Definition}
\newtheorem*{restatedclaim}{\RestatedName}
\theoremstyle{claimass}
\newtheorem{assumption}[theorem]{Assumption}
\newtheorem*{restatedassumption}{\RestatedName}
\theoremstyle{claimremark}
\newtheorem{remark}[theorem]{Remark}
\newtheorem{example}[theorem]{Example}
\theoremstyle{claimreq}
\newtheorem{requirement}[theorem]{Requirement}
\newtheorem*{restatedrequirement}{\RestatedName}
%% --- 8. Prominence: main results a margin rule and a bold heading ----------
\newtcolorbox{mainclaim}{enhanced jigsaw, breakable, blanker,
  borderline west={1.5pt}{-8pt}{mainRule}, left=0pt, right=0pt, top=0pt, bottom=0pt,
  boxsep=0pt, before skip=\topsep, after skip=\topsep, parbox=false,
  before upper={\renewcommand*{\ClaimNoteFont}{\bfseries}}}
\newenvironment{technicalclaim}{\renewcommand*{\ClaimHeadFont}{\mdseries\color{techGrey}}%
  \renewcommand*{\ClaimNoteFont}{}}{}
%% --- 9. Non-normative text, the meta line, navigation ------------------------
\makeatletter
% Tie what follows to this paragraph: no page break, no list penalty before
% the next theorem head (the \@nobreak route a section heading uses).
\newcommand{\keepwithnext}{\par\nobreak\vspace{3pt}\@nobreaktrue
  \everypar{\if@nobreak\@nobreakfalse\clubpenalty\@M\else\clubpenalty\@clubpenalty\everypar{}\fi}}
\newcommand{\claimsummary}[2]{% #1 claim id (named destination), #2 summary
  \par\addvspace{\bigskipamount}\Needspace{6\baselineskip}%
  {\small\itshape\color{nnGrey}\noindent\hypertarget{claim-#1}{}#2\par}\keepwithnext}
\newcommand{\claimanchor}[1]{% a claim without a summary: the destination alone
  \par\addvspace{\bigskipamount}\Needspace{6\baselineskip}%
  \phantomsection\hypertarget{claim-#1}{}\keepwithnext}
\newcommand{\claimleadin}[1]{{\color{nnGrey}\noindent#1\par}\keepwithnext}
\newcommand{\claimmeta}[1]{\par\nobreak\vspace{2pt}%
  {\footnotesize\color{metaGrey}\raggedright\interlinepenalty\@M\noindent#1\par}%
  \addvspace{\medskipamount}}
% After an assumption: the sentence generated from its verification field; after
% any claim with one: its plain-text reading (user_statement). Both
% non-normative, both grey.
\newcommand{\claimverification}[1]{\par\nobreak\vspace{2pt}%
  {\footnotesize\color{metaGrey}\noindent#1\par}\nobreak}
\newcommand{\claimreading}[1]{\par\nobreak\vspace{2pt}%
  {\footnotesize\color{nnGrey}\noindent\emph{Reading:}\enspace#1\par}\nobreak}
\makeatother
\newcommand{\claimbookmark}[4]{% #1 id, #2 kind word, #3 label, #4 plain title
  \bookmark[dest=claim-#1,level=2]{#2 \getrefnumber{#3} #4}}
\newcommand{\metasep}{\unskip\enspace\textperiodcentered\enspace\ignorespaces}
\newcommand{\metalabel}[1]{\emph{#1}\enspace}
\newcommand{\sectionintro}[1]{{\color{nnGrey}\noindent#1\par}\smallskip}
\newcommand{\sectionimports}[1]{{\footnotesize\color{metaGrey}\raggedright\noindent\metalabel{Imports}#1\par}\smallskip}
\newenvironment{overview}{\par\color{nnGrey}\itshape\noindent\ignorespaces}{\par}
% References: the whole "Def. 2.1" is the link; ids are the fallback text of
% every link into another PDF (a relative URI to a named destination).
\newcommand{\mref}[2]{\hyperref[#2]{#1~\ref*{#2}}}
\newcommand{\aref}[2]{\hyperref[#2]{\textsf{#1}}}
\DeclareUrlCommand{\claimid}{\urlstyle{sf}\def\UrlBreaks{\do\/\do\:\do\-\do\.\do\_}%
  \def\UrlBigBreaks{}\def\UrlNoBreaks{}}
\newcommand{\extref}[2]{\hyperref{#1}{}{claim-#2}{\claimid{#2}}}
\newcommand{\extaref}[3]{\hyperref{#1}{}{claim-#2}{\textsf{#3}}}
\newcommand{\coderef}[2]{\href{#1}{\claimid{#2}}}
% Claim-body commands (spec/claim-format.md section 4): \xref{<full id>} prints
% the kind, title and id of a claim of another contract, linked to its stable
% destination in that contract's PDF (one \DeclareXref per cited claim below);
% \term[gloss]{phrase} marks a defined phrase for the index (the gloss feeds
% the Terms table); \tref[phrase]{text} is a use of a defined term. The
% assembler rewrites both in the body into \termdef (the definition's named
% destination term-<slug>) and \termref / \exttermref (a link to it, in this
% or another contract's PDF); the fallbacks below serve text it copies
% elsewhere, such as the digest. \iface is an interface key.
\makeatletter
\newcommand{\DeclareXref}[4]{% id, kind word, title, contract PDF
  \@namedef{xref@#1}{\hyperref{#4}{}{claim-#1}{#2 ``#3'' ({\small\claimid{#1}})}}}
\newcommand{\xref}[1]{\@ifundefined{xref@#1}%
  {\textbf{[unresolved \nolinkurl{#1}]}\@latex@warning{Unresolved xref #1}}%
  {\@nameuse{xref@#1}}}
\makeatother
\DeclareUrlCommand{\iface}{\urlstyle{tt}}
\newcommand{\term}[2][]{\emph{#2}\index{#2}}
\newcommand{\tref}[2][]{#2}
\newcommand{\termdef}[2]{\hypertarget{#1}{}\emph{#2}\index{#2}}
\newcommand{\termref}[2]{\hyperlink{#1}{#2}}
\newcommand{\exttermref}[3]{\hyperref{#1}{}{#2}{#3}}
\hypersetup{pdftitle={Transmitter source model (transmitter) -- statements digest 0.13.0}}

% The digest: the contract's labels under the prefix C-, its numbers checked.
\externaldocument[C-]{contract}
\theoremstyle{claim}\newtheorem*{digestclaim}{\RestatedName}
\theoremstyle{claimass}\newtheorem*{digestassumption}{\RestatedName}
\theoremstyle{claimreq}\newtheorem*{digestrequirement}{\RestatedName}
\makeatletter
% #1 label, #2 the number the assembler computed: a local label carrying it, so
% that a reference between two digest items stays in the digest; the build
% stops when the contract's own number differs.
\newcommand{\digestlabel}[2]{\phantomsection
  \edef\digest@theirs{\getrefnumber{C-#1}}\def\digest@ours{#2}%
  \ifx\digest@theirs\digest@ours\else
    \PackageError{digest}{#1 is #2 here but \digest@theirs\space in the contract}{}\fi
  \def\@currentlabel{#2}\label{#1}}
\makeatother
% A display's local label (amsmath records the \tag, i.e. the contract's number).
\newcommand{\digesteqlabel}[1]{\label{#1}}

% Claims of other contracts cited here (\xref, [[...]]).
\DeclareXref{meat-core/def:declaration}{Definition}{The protocol declaration as the MEAT core reads it}{../meat-core/contract.pdf}
\DeclareXref{meat-core/def:roundmodel}{Definition}{Round model and licensed attack sets}{../meat-core/contract.pdf}
\DeclareXref{meat-core/lem:product-marginal}{Lemma}{Product marginal of Alice's retained systems}{../meat-core/contract.pdf}
\DeclareXref{meat-core/rem:modes-da1}{Remark}{The two source modes on the program}{../meat-core/contract.pdf}
\DeclareXref{protocol-compiler/ass:link}{Assumption}{Soundness of every reduction link}{../protocol-compiler/contract.pdf}
\DeclareXref{protocol-compiler/def:chaincond}{Definition}{Interface conditions on a chain}{../protocol-compiler/contract.pdf}
\DeclareXref{protocol-compiler/def:declaration}{Definition}{Protocol declaration}{../protocol-compiler/contract.pdf}
\DeclareXref{protocol-compiler/def:eligible}{Definition}{Key-length eligibility}{../protocol-compiler/contract.pdf}
\DeclareXref{protocol-compiler/prop:conservative}{Proposition}{Dropping a trusted marginal enlarges the feasible set}{../protocol-compiler/contract.pdf}
\DeclareXref{protocol-compiler/thm:chain-consumption}{Theorem}{Decided identities and named obligations of a consumed chain}{../protocol-compiler/contract.pdf}
\DeclareXref{reductions/def:chain}{Definition}{Reduction chain}{../reductions/contract.pdf}
\DeclareXref{reductions/def:link}{Definition}{Typed link}{../reductions/contract.pdf}
\DeclareXref{reductions/lem:precomposition}{Lemma}{Licensed admissibility survives precomposition with a source channel}{../reductions/contract.pdf}
\DeclareXref{reductions/lem:source-transfer}{Lemma}{Exactness of the source map in trusted-source mode}{../reductions/contract.pdf}
\DeclareXref{reductions/thm:two-sided}{Theorem}{The two-sided reduction in trusted-source mode}{../reductions/contract.pdf}

\begin{document}
\begin{center}
{\LARGE\bfseries Transmitter source model (transmitter)}\\[4pt]
{\large Statements digest, version 0.13.0}
\end{center}
\noindent{\small Every definition, fact, assumption, result and requirement of the \href{contract.pdf}{contract}, with the contract's numbers; remarks, examples, technical steps and all proofs are left out. Each heading opens its statement in the contract.}

\part*{Mathematics}\pdfbookmark[0]{Mathematics}{digest-part-1}

\section*{1\quad Source specifications and source modes}\pdfbookmark[1]{1 Source specifications and source modes}{digest-group-transmitter/def:source-spec}

% claim transmitter/def:source-spec
\digestlabel{def:source-spec}{1.1}\pdfbookmark[2]{Definition 1.1 Source specification}{digest-transmitter/def:source-spec}
\renewcommand*{\RestatedName}{\hyperref{contract.pdf}{}{claim-transmitter/def:source-spec}{Definition~1.1}}%
\begin{digestclaim}[Source specification]
A \emph{source specification} $\Spec_{\mathrm{src}}$ consists of the
following data. A finite non-empty alphabet $\Xset$ with a declared order. A
finite-dimensional system $A$ with a declared computational basis,
$d_A\coloneqq\dim\Hs{A}$, and a declared \emph{sector structure} on it: an
ordered finite index set $\mathbb{T}$ and a decomposition
$\Hs{A}=\bigoplus_{t\in\mathbb{T}}\Hs{A_t}$ into subspaces each spanned by
vectors of the declared basis, $d_{A_t}\coloneqq\dim\Hs{A_t}\ge1$, with
projectors $\Pi^{A}_t$; the \emph{one-sector structure}, $\mathbb{T}$ a
single index and $A_t=A$, is the default. A POVM $\{M_x\}_{x\in\Xset}$ on
$A$ presented as a block list against that structure,
$M_x=\bigoplus_{t}M^{(t)}_x$ with $M^{(t)}_x\in\Lin{A_t}$, each block
Hermitian and positive semidefinite, and
$\sum_{x\in\Xset}M^{(t)}_x=\idon{A_t}$ for every $t$ with the sum taken in
the declared order. The \emph{specified signal ensemble}
$\{p_x,\rho_x\}_{x\in\Xset}$ on the \emph{transmitted system} $\Asig$, with
$p_x\ge0$, $\sum_xp_x=1$ and $\rho_x$ a state on $\Asig$. With one sector
the block list $M_x$ is the matrix $M_x$, and every condition on its blocks
is the same condition on that matrix.

\emph{How the specified ensemble is declared.} In one of two ways.
(a) \emph{As matrices}, for a finite-dimensional $\Asig$: by its
\emph{weighted members} $\sigma_x\in\Pos{\Asig}$, with
$p_x\coloneqq\tr\sigma_x$, $\sum_xp_x=1$ in the declared order, and
$\rho_x\coloneqq\sigma_x/p_x$ where $p_x>0$. A declaration by the pair
$(p_x,\rho_x)$, $\rho_x\in\Seq{\Asig}$, is the declaration of the weighted
member $\sigma_x=p_x\rho_x$ with the product formed exactly, and it requires
$\rho_x\in\Seq{\Asig}$ of \emph{every} letter, a letter with $p_x=0$
included. Every condition of this module that reads an admitted ensemble
reads it through its weighted members only.
(b) \emph{By a} \emph{source model}, for a transmitted system $\Asig$ whose
Hilbert space is separable, possibly infinite dimensional: by a named
definition of a submodule of this module that fixes the weights $p_x$ and the
states $\rho_x$ as mathematical objects and not as matrices, together with
the physical-source assumption which that submodule declares in the place of
Assumption~\ref{ass:ens}. A specification declared by a source model is
presented with a non-empty source segment
(Definition~\ref{def:source-segment}): what a consumer receives is never the
model but the virtual ensemble, which is finite dimensional.

\emph{The two modes.} The specification is declared in one of two modes. In
\emph{trusted-source mode} it additionally carries a marginal
$\tau_A\in\Seq{A}$, presented as a block list
$\tau_A=\bigoplus_t\tau^{(t)}_A$ with $\tau^{(t)}_A\in\Pos{A_t}$ and
$\sum_t\tr\tau^{(t)}_A=1$, and, at $d_A\ge2$, a source-replacement
certificate in the sense of Definition~\ref{def:srcert}, written against the
\emph{virtual} ensemble $\{p_x,\rvirt_x\}$ on the virtual transmitted system
$\Avirt$: the finite-dimensional ensemble the source segment starts from, whose
channels produce the specified ensemble from it
(Definition~\ref{def:source-segment}\,(S1), (S3)); when the source segment is
empty, $\Avirt=\Asig$ and $\rvirt_x=\rho_x$. The certificate takes one of the
three forms of Definition~\ref{def:srcert}: a rank-one projector (the pure
form), positive sector blocks (the block form), or an ensemble presentation,
which declares only the weighted virtual ensemble, in blocks along a matched
sector structure of $\Avirt$, and defines the certificate from it. A sector structure with more than one sector is declared only in
trusted-source mode and only together with a certificate in block form or an
ensemble presentation. \emph{With an ensemble presentation} the specification
declares neither $A$, its basis and its sector structure, nor $\{M_x\}$, nor
$\tau_A$: they are constituents of the specification that \eqref{eq:enspres}
defines from the declared virtual ensemble, the conditions this definition
places on them hold by Lemma~\ref{lem:ensemble-presentation}\,(i), and of
them only $\tau_A$ is formed as numbers, that is, computed entry by entry from
the declared blocks (Lemma~\ref{lem:ensemble-presentation}\,(iii));
$\{M_x\}$ and the certificate are defined and never computed. At $d_A=1$ no
certificate is demanded and $\tau_A=[1]$ is the only marginal, so that there
the two modes impose the same condition. In
\emph{untrusted-source mode} the specification carries no certificate and no
marginal, written
$\tau_A=\nomarg$; the specification with
$d_A=1$ and a one-letter $\Xset$ is its degenerate case, in which Alice
prepares nothing that this module describes.

\emph{What the sector structure is not.} It is a declaration about
$\{M_x\}$ and $\tau_A$, which are block diagonal by their presentation. It
is no statement about any operator on $A$ together with a further system,
and no consumer may read it as one: that a joint state is block diagonal
along $\mathbb{T}$ is either proved of the states a certificate reaches
(Lemma~\ref{lem:matched-pinching}) or licensed by a theorem of the consuming
module that restricts its own feasible set, under that module's own
preconditions, and is never inferred from the presentation of $\tau_A$ or
$\{M_x\}$.

\emph{The number model.} Every declared entry is a binary64 number, hence a
dyadic rational, and is admissible exactly as declared. A source model
(case (b)) declares no entries: its weights and states are mathematical
objects, and what is decided of such a specification is decided on the
entries it does declare, those of its virtual ensemble among them. Every
condition above is decided \emph{exactly} on the declared entries: a
quantity that a condition compares --- a sum in the declared order, a
product, a partial
trace --- is \emph{formed} in exact arithmetic on the declared entries as
well as compared exactly, so that the identity decided is the identity
written and rounding enters nowhere. That the operands are binary64 is not
what makes a decision exact; a product of two binary64 numbers need not be
one. Positive semidefiniteness is decided by an exact decision on the
declared entries and never by a floating-point eigenvalue.
\end{digestclaim}

\section*{2\quad Source links and the source segment}\pdfbookmark[1]{2 Source links and the source segment}{digest-group-transmitter/def:source-link}

% claim transmitter/def:source-link
\digestlabel{def:source-link}{2.1}\pdfbookmark[2]{Definition 2.1 Source link}{digest-transmitter/def:source-link}
\renewcommand*{\RestatedName}{\hyperref{contract.pdf}{}{claim-transmitter/def:source-link}{Definition~2.1}}%
\begin{digestclaim}[Source link]
A \emph{source link} $\Xi$ is a typed link of a reduction chain
(\xref{reductions/def:chain}) at the stage
\textsf{source}: a tuple
\begin{equation*}\tag{\getrefnumber{C-eq:source-link}}\digesteqlabel{eq:source-link}
  \Xi=\Big(\textsf{kind},\
  \big(A^{\mathrm{in}},\{\sigma^{\mathrm{in}}_x\}_{x\in\Xset}\big),\
  \big(A^{\mathrm{out}},\{\sigma^{\mathrm{out}}_x\}_{x\in\Xset}\big),\
  \Gamma_{\Xi},\ \mathcal{D}_{\Xi},\ \Smap_{\Xi},\
  \CList_{\Xi},\ \textsf{licence},\ \textsf{status}\Big)
\end{equation*}
whose slots are the following.

\emph{(a) The kind.} \textsf{kind} is a source-stage link kind. The
vocabulary is the kinds this module declares, of which there is exactly one,
\textsf{source-map}, together with the kinds a submodule of this module
declares in its own interface. A submodule that declares a kind supplies its
soundness claim, and a kind for which no claim of this repository proves the
link's identity is declared and refused
(Definition~\ref{def:source-segment}, ``Which kinds are admitted''), as
\textsf{source-map} is. An admitted kind has an entry in the soundness
registry of that definition, and a kind without one is not admitted whatever
a link declares.

\emph{(b) The two descriptors.} The \emph{source descriptor}
$(A^{\mathrm{in}},\{\sigma^{\mathrm{in}}_x\})$ is the transmitted system and
signal ensemble the link starts from and the \emph{target descriptor}
$(A^{\mathrm{out}},\{\sigma^{\mathrm{out}}_x\})$ the one it produces, both
indexed by the same alphabet $\Xset$ in the same declared order, and both
read through their weighted members $\sigma_x=p_x\rho_x$, as every ensemble
of this module is (Definition~\ref{def:source-spec}): $\sigma^{\mathrm{in}}_x$ and
$\sigma^{\mathrm{out}}_x$ are positive semidefinite with
$\tr\sigma^{\mathrm{in}}_x=\tr\sigma^{\mathrm{out}}_x=p_x$, which is what
carrying the weights untouched means.

\emph{(c) The channel witness, in one of two ways.} $\Gamma_{\Xi}$ is a
channel from $A^{\mathrm{in}}$ to $A^{\mathrm{out}}$, supplied either
\emph{as data} or \emph{as a theorem constructor}.
\begin{enumerate}[label=(c\arabic*),leftmargin=2.6em]
\item \emph{As data}: both systems are finite dimensional and the channel is
declared, so that the identity
$\Gamma_{\Xi}(\sigma^{\mathrm{in}}_x)=\sigma^{\mathrm{out}}_x$ for every
$x\in\Xset$ is an exact matrix identity on declared numbers and is decided as
such (Definition~\ref{def:source-segment}\,(S3)).
\item \emph{As a} \emph{theorem constructor}: the witness is supplied by a
named claim of a submodule of this module, which constructs one channel for
each member of a declared family of descriptors and proves the identity of
(c1) for it. The identity is then \emph{not} decided on numbers: the link
carries an \emph{obligation record} of four components, the triple
\begin{equation*}\tag{\getrefnumber{C-eq:obligation}}\digesteqlabel{eq:obligation}
  \big(\textsf{kind},\ \text{the id of the soundness claim},\
  \mathcal{A}(\textsf{kind})\big)
\end{equation*}
and, as the fourth, the member of the declared family the link's descriptors
are. The third component is the assumptions the obligation carries, which
Definition~\ref{def:source-segment}\,(A1) defines as the soundness claim's
assumption closure together with the physical-source assumption the
submodule declares for the source model the claim is stated over, an
assumption the soundness claim itself does \emph{not} assume. The obligation
record travels with the link and, where the link carries a licence record, in
it. \emph{Its first three components are decided}, against the soundness
registry of the artefact that carries the link
(Definition~\ref{def:source-segment}\,(A1), (A2)), by a comparison of
declared strings. The fourth is not: which member of the family the
descriptors are is asserted by the module that emits the link and re-proved
by no consumer, and no verdict on a number is carried, there being none to
decide.
\end{enumerate}
A channel supplied as a theorem constructor is how a link whose source or
target is not a matrix is expressed at all, and it is the only way an
obligation enters this module: no claim here discharges one.

\emph{(d) The declared descriptor family.} $\mathcal{D}_{\Xi}$, the
\emph{declared descriptor family}, is the family of source descriptors the
witness's identity is asserted for. It contains the declared source
descriptor; for a link whose channel is supplied as data it is that
descriptor alone; for a constructor link it is the family the named claim
proves the identity for. The link is \emph{sound} when
$\Gamma_{\Xi}(\sigma^{\mathrm{in}}_x)=\sigma^{\mathrm{out}}_x$ holds for
every $x\in\Xset$ at every member of $\mathcal{D}_{\Xi}$; for a data link
this is what (S3) decides, and for a constructor link it is what the named
claim concludes and the obligation record names. A link whose family does not
contain its own declared descriptor asserts nothing about the descriptor it
carries, and is refused. For a constructor link that membership is not a
decidable condition: the family is the extension of a hypothesis of a claim,
the emitting module asserts that its descriptors satisfy it, and what a
consumer decides instead is the triple \eqref{eq:obligation} against the
registry and the provenance of the artefact
(Definition~\ref{def:source-segment}\,(A1) to (A3)). The refusal of a family
that excludes its own descriptor therefore applies where membership is
decidable, to a data link, whose family is by definition its own descriptor,
so that it refuses a data link whose family slot excludes that descriptor;
elsewhere
membership is an obligation of the emitting module, and what the registry
bounds is which claims a link may name for it.

\emph{(e) The retained-register sector map.} $\Smap_{\Xi}$, the
\emph{retained-register sector map}, is optional: a link may emit a matched
pair of sector structures over one ordered index set, one on Alice's
retained system $A$ and one on the Alice-side descriptor \emph{the
certificate is written against}, which is the \emph{source} descriptor of the
segment's first link, the virtual pair
(Definition~\ref{def:source-segment}\,(S1)), in the sense the block form of
Definition~\ref{def:srcert} is written against. It is
not the descriptor the link produces: a link's target may be a system that
carries no sector structure at all, and for a constructor link whose target
is infinite dimensional it never carries one. The sector map is a
presentation datum: it says how the certificate, the marginal and Alice's
measurement are presented, and it is no statement about any joint state
(Definition~\ref{def:source-spec}, ``What the sector structure is not'').
Where the slot is empty the one-sector structure is meant.

\emph{(f) The transported constraint entries and the licence.}
$\CList_{\Xi}$ is the list of \emph{transported constraint entries}, entries
of a consumer's constraint list carried across the link, each
carrying, \emph{per entry}, the certificate case that discharges it --- the
statement, for that entry, that what the entry asserts of the link's target
is implied by what holds of its source --- and \textsf{licence} is the
link's licence record, the slot every link of a chain has for its transported
entries (\xref{reductions/def:link}), present exactly when
$\CList_{\Xi}\ne\varnothing$. \emph{A non-empty list is refused}: a
consumer's constraint operators live on the measured system, and no claim of
this repository transports an Alice-side entry. An admitted link therefore
carries an empty list and no licence record of its own.

\emph{(g) The status.}
$\textsf{status}\in\{\textsf{exact},\textsf{certified\_outer},
\textsf{exploratory\_numeric}\}$, the link statuses of
\xref{reductions/def:link}, with \textsf{exact} meaning that the
identity of (c1) holds with no error term --- for a constructor link, that
the named claim concludes it with none.

\emph{Alice-side descriptor discipline.} Both descriptors of a source link
are Alice-side: each names a transmitted system and an ensemble on it, and
no condition on a source link mentions the receiver's systems, its outcome
alphabets or its measurements. A source link changes what Alice emits and
nothing else.

\emph{The degenerate link.} A link that supplies its channel as data, whose
declared descriptor family is that one descriptor, whose sector-map slot is
empty and whose entry list is empty gains nothing from its slots (d) and
(e): they carry exactly what the identity of (c1) asserts and the one-sector
structure.
\end{digestclaim}

% claim transmitter/def:source-segment
\digestlabel{def:source-segment}{2.2}\pdfbookmark[2]{Definition 2.2 Source segment and its interface conditions}{digest-transmitter/def:source-segment}
\renewcommand*{\RestatedName}{\hyperref{contract.pdf}{}{claim-transmitter/def:source-segment}{Definition~2.2}}%
\begin{digestclaim}[Source segment and its interface conditions]
The \emph{source segment} of a reduction chain is an ordered list
$\Chsrc=(\Xi_1,\dots,\Xi_{L_{\mathrm s}})$ of source links
(Definition~\ref{def:source-link}), $L_{\mathrm s}\ge0$, whose links act on
the transmitted system and replace Alice's signal ensemble; it is presented
together with a source specification $\Spec_{\mathrm{src}}$
(Definition~\ref{def:source-spec}). The segment satisfies the
\emph{interface conditions} when the following hold.
\begin{enumerate}[label=(S\arabic*),leftmargin=3.0em]
\item \emph{(descriptor match)} the target descriptor of $\Xi_i$ equals the
source descriptor of $\Xi_{i+1}$ for $1\le i<L_{\mathrm s}$; the target
descriptor of $\Xi_{L_{\mathrm s}}$ is the specified pair, the transmitted
system $\Asig$ named by $\Spec_{\mathrm{src}}$ with the specified ensemble
$\{p_x,\rho_x\}$ on it; and the source descriptor of $\Xi_1$ is the virtual
pair, the \emph{virtual transmitted system} $\Avirt$ with the
\emph{virtual ensemble} $\{p_x,\rvirt_x\}_{x\in\Xset}$ on it, on which the
source-replacement certificate is written (Definition~\ref{def:srcert}).
Every descriptor of
this segment is an Alice-side descriptor: no condition of the segment
mentions the receiver's systems, its outcome alphabets or its measurements;
\item \emph{(index preservation)} each link carries the signal index $x$
through unchanged: it maps a family indexed by $\Xset$ to a family indexed
by $\Xset$, with the weights $p_x$ untouched;
\item \emph{(source identity, where checkable)} writing
$\Xi\coloneqq\Xi_{L_{\mathrm s}}\circ\dots\circ\Xi_1$ for the composite
when every link supplies its channel as data,
\begin{equation*}\tag{\getrefnumber{C-eq:srcmap}}\digesteqlabel{eq:srcmap}
  \Xi(\rvirt_x)=\rho_x\qquad\text{for every }x\in\Xset ,
\end{equation*}
an \emph{exact} matrix identity, decided entry by entry. The composite
$\Xi$ is \emph{formed} in exact arithmetic as well, so that the identity
decided is the identity written; that the operands are binary64 is not what
makes the decision exact (the number model of
Definition~\ref{def:source-spec}). The channel runs from the \emph{virtual}
states to the \emph{specified} ones, and the object a consumer receives is
the virtual protocol, the protocol with the virtual ensemble in the place of
the specified one. \emph{Where a link supplies its channel as a theorem
constructor} (Definition~\ref{def:source-link}\,(c2)) the identity of
\emph{that} link is not decided on numbers: it is the conclusion of the claim
the link's obligation record names, at the member of that claim's declared
descriptor family which the link's descriptors are, and the identity of the
composite is the composition of the per-link identities --- each one either
decided, where its link supplies data, or concluded by its named claim, where
its link supplies a constructor. In a segment that mixes the two kinds of
witness, each data link's identity is therefore decided link by link. So
\eqref{eq:srcmap} is what the composite satisfies in every case, and what
changes with the kind of witness is
\emph{by what} it is established and therefore what a consumer may re-decide:
nothing of a constructor link's identity, which is why the obligation
travels;
\item \emph{(slots)} each link carries its constraint entries and its
licence slot, both possibly empty, and its status tag; a link whose channel
is a theorem constructor carries in addition its obligation record, the claim
and the assumptions it names, which are the strings (A2) below decides.
\end{enumerate}

\emph{Which kinds are admitted, and what a consumer decides.} \emph{No
source link kind is admitted unless the submodule that declares it supplies
its soundness claim}: a link of kind $k$ may appear in a segment only where a
submodule of this module declares $k$ and supplies, as a claim of this
repository, the statement that the link's witness satisfies the identity of
Definition~\ref{def:source-link}\,(c1) at every member of its declared
descriptor family --- decided there, for a data link, or proved there, for a
constructor link. A link of a kind for which no such claim exists is
\emph{refused}, and so is a link with a non-empty list of transported
constraint entries, whatever its kind: no claim of this repository transports
an Alice-side entry.

That much is a condition on the \emph{kind}, and a consumer cannot decide it
of a claim it does not read. What a consumer decides is the following, and it
is decided on data.
\begin{enumerate}[label=(A\arabic*),leftmargin=3.0em]
\item \emph{(the soundness registry)} this module emits, with every artefact,
a finite list --- the \emph{soundness registry} --- with one entry per
admitted source-link kind, \emph{closed} in that a kind without an entry is
not admitted (Definition~\ref{def:source-link}\,(a)), each entry the triple
\begin{equation*}\tag{\getrefnumber{C-eq:registry}}\digesteqlabel{eq:registry}
  \big(\textsf{kind},\ \text{the id of its soundness claim},\
  \mathcal{A}(\textsf{kind})\big) ,
\end{equation*}
written out in the module's interface and emitted verbatim. \emph{The third
component $\mathcal{A}(\textsf{kind})$ is the assumptions the obligation
carries}, and it has two sources, both of which a registry entry lists:
\begin{enumerate}[label=(\alph*),leftmargin=2.2em]
\item the \emph{assumption closure of the soundness claim}: the assumptions
in the transitive closure of that claim under its \emph{uses} and
\emph{assumes} edges, the claims and assumptions each claim lists as its
dependencies; it may be empty, and is empty if and only if no claim of that
closure, the soundness claim included, lists an assumption among its
dependencies;
\item the \emph{physical-source assumption the submodule declares for the
source model the claim is stated over}, where the claim is stated over one.
That assumption is not a premise of the soundness claim and must not be made
one: what it attests is that the model the claim is about describes the
physical device, which is what an obligation on a physical source owes and a
mathematical identity does not.
\end{enumerate}
The registry of an artefact is the registry of the version of this module that
emitted it, not a datum a request supplies.
\item \emph{(the triple is decided)} a presented chain's link is
\emph{refused} unless the triple (its \textsf{kind}, the claim id of its
obligation record, the assumption short ids of that record) appears in the
registry of the artefact that carries it --- a comparison of finitely many
declared strings, taken entry by entry --- and unless every assumption of that
entry is in the artefact's assumption ledger, the list of assumptions the
artefact rests on. \emph{Both components of
$\mathcal{A}(\textsf{kind})$ are compared}, so a link that names the soundness
claim's closure and omits the model's attestation is refused exactly as one
that names neither. A link supplying its channel as data carries no obligation
record and is decided by (S3) instead.
\item \emph{(provenance)} a non-empty source segment is accepted only inside
an artefact emitted by the module that admits its links. A segment assembled
from links a consumer received in any other way is refused, whatever its
triples: the registry says which kinds this repository has proved, not that a
particular link was built by the module that proved them.
\end{enumerate}
The admission rule on the kind is applied where a link is admitted, which is
this module and the submodule whose kind it is; it is not among the conditions
(S1) to (S4), which are the conditions a consumer decides on a presented
chain, and (A1) to (A3) are conditions on the artefact and its registry rather
than on the segment's descriptors.

\emph{The empty segment.} When $L_{\mathrm s}=0$ the segment is empty, $\Xi$
is the identity, $\Avirt=\Asig$ and $\rvirt_x=\rho_x$, so \eqref{eq:srcmap}
holds by inspection. That is the case of every specification whose ensemble
is declared as matrices and that carries no link. A specification whose
ensemble is declared by a source model is presented with a non-empty segment
(Definition~\ref{def:source-spec}, case (b)), which is the other case and the
one the conditions above are written for.
\end{digestclaim}

\section*{3\quad The source-replacement certificate}\pdfbookmark[1]{3 The source-replacement certificate}{digest-group-transmitter/fact:sr}

% claim transmitter/fact:sr
\digestlabel{fact:sr}{3.1}\pdfbookmark[2]{Fact 3.1 Source replacement for prepare-and-measure protocols}{digest-transmitter/fact:sr}
\renewcommand*{\RestatedName}{\hyperref{contract.pdf}{}{claim-transmitter/fact:sr}{Fact~3.1}}%
\begin{digestclaim}[Source replacement for prepare-and-measure protocols]
\emph{(\cite[Lemma~8.3, Eqs.~(83)--(84)]{tupkary2026decoy}, with the
protocol framing of its Def.~8.1 and Def.~8.2.)}
Consider a protocol in which Alice first prepares the global state
$\rho_{X_1^n(A')_1^n}=\bigotimes_{j=1}^{n}\rho^{(j)}_{X_jA_j'}$ with
$\rho^{(j)}_{X_jA_j'}=\sum_{k=1}^{d_A}p^{(j)}_k\dyad{k}_{X_j}\otimes
(\rho^{(j)}_k)_{A_j'}$, where $\{\ket{k}\}_{X_j}$ is the classical basis of
$X_j$; she keeps $X_1^n$ and sends the $(A')_1^n$ systems to the adversary.
Then the state $\rho_{X_1^n(A')_1^n}$ can equivalently be obtained by
\emph{(1)} preparing the global state
$\rho_{\Abar_1^n\Ahat_1^n(A')_1^n}=\bigotimes_j\rho^{(j)}_{\Abar_j\Ahat_jA_j'}$
with
$\rho^{(j)}_{\Abar_j\Ahat_jA_j'}=\sum_{k,i=1}^{d_A}
\sqrt{p^{(j)}_ip^{(j)}_k}\,\ketbra{i}{k}_{\Abar}\otimes
\ketbra{\psi^{(j)}_i}{\psi^{(j)}_k}_{\Ahat_jA_j'}$, where
$\ket{\psi^{(j)}_k}_{\Ahat_jA_j'}$ is a purification of
$(\rho^{(j)}_k)_{A_j'}$, $\{\ket{k}\}$ is an orthonormal basis of $\Abar$
and $\Ahat_j$ is the \emph{shield system}; \emph{(2)} sending $(A')_1^n$ to
the adversary; and \emph{(3)} for each $j$, measuring $\Abar_j\Ahat_j$ with
the POVM $\{\dyad{k}_{\Abar_j}\otimes\idon{\Ahat_j}\}_{k\in\{1,\dots,d_A\}}$
and storing the result in $X_j$.

\emph{Symbol map and hypothesis check.} The source's $X_j$ is this module's
outcome register with alphabet $\Xset$ and $k\leftrightarrow x$; its $A_j'$
is the transmitted system carrying the ensemble the lemma is applied to,
and its $\{p^{(j)}_k,\rho^{(j)}_k\}$ is that ensemble. That ensemble is the
virtual pair $(\Avirt,\{p_x,\rvirt_x\})$ on which the source-replacement
certificate is stated, which is the specified pair $(\Asig,\{p_x,\rho_x\})$
of Definition~\ref{def:source-spec} exactly when the source segment is empty.
This module describes one round; that every round is described by the
same data is the consumer's round-homogeneity assumption and is not
asserted here, and it is under that assumption that the superscript $j$ is
dropped in what follows. The source's $\Abar_j\Ahat_j$ is this module's
retained system $A$, and its measurement POVM is $\{M_x\}$ with
$M_x=\dyad{x}_{\Abar}\otimes\idon{\Ahat}$. The $d_A$ of the cited lemma and
the $d_A$ of this module are \emph{not} the same number, and the
identification does not claim they
are: the published index runs $k\in\{1,\dots,d_A\}$ over an orthonormal
basis of $\Abar$, so there $d_A=\dim\Hs{\Abar}=|\Xset|$, whereas throughout
this module $d_A=\dim\Hs{A}=\dim\Hs{\Abar}\cdot\dim\Hs{\Ahat}$; the two
agree exactly when the shield system is trivial. The lemma's one hypothesis
is that Alice prepares the classical-quantum source state of its
Eq.~(83) --- this is Assumption~\ref{ass:ens} for the physical device, and
the declared source data for the specified object. Its only timing element is
that the global state is prepared first; it places no condition on the
sequential structure of later operations, being an identity between the state
of its Eq.~(83) and the measured state of its Eq.~(84). The cited work
applies it within the protocol framing of its Def.~8.1 and Def.~8.2, and this
module applies it to protocols with the sequential structure of
Assumption~\ref{ass:seq}, which supplies that structure: the assumption is the
hypothesis under which this module applies the fact, and not a hypothesis of
the cited lemma. The
published statement quantifies over an
arbitrary purification $\ket{\psi_k}$, so the shield system is trivial
exactly when every state of that ensemble is pure --- the case, for a pure
\emph{specified} ensemble, of \cite[Eqs.~(37)--(39)]{winick2018reliable},
which writes the source-replaced vector as
$\ket{\psi}_{\Abar\Asig}=\sum_i\sqrt{p_i}\ket{i}_{\Abar}\ket{\varphi_i}_{\Asig}$
and records that the resulting marginal
$\rho_{\Abar}=\tr_{\Asig}(\dyad{\psi})$ is fixed independently of the
adversary's attack.
\end{digestclaim}

% claim transmitter/def:srcert
\digestlabel{def:srcert}{3.3}\pdfbookmark[2]{Definition 3.3 Source-replacement certificate}{digest-transmitter/def:srcert}
\renewcommand*{\RestatedName}{\hyperref{contract.pdf}{}{claim-transmitter/def:srcert}{Definition~3.3}}%
\begin{digestclaim}[Source-replacement certificate]
Let $\{p_x,\rvirt_x\}_{x\in\Xset}$ be the \emph{virtual} signal ensemble on
the virtual transmitted system $\Avirt$ --- the pair the source segment
starts from, Definition~\ref{def:source-segment}\,(S1), (S3) --- with
weighted members $\sigma^{\mathrm v}_x\coloneqq p_x\rvirt_x\in\Pos{\Avirt}$,
$p_x=\tr\sigma^{\mathrm v}_x\ge0$ and $\sum_xp_x=1$. These conditions are
decided \emph{exactly}, as every declared condition is, the weight sum in
the declared order and rounding nowhere. A \emph{source-replacement
certificate} for the declared data $(A,\mathbb{T},\tau_A,\{M_x\})$ of a
source specification (Definition~\ref{def:source-spec}) is an operator
$\Psi\in\Lin{A\Avirt}$ such that
\begin{enumerate}[label=(SR\arabic*),leftmargin=3.0em]
\item $\tr_{\Avirt}\Psi=\tau_A$;
\item $\tr_{A}\!\big[(M_x\otimes\idon{\Avirt})\,\Psi\big]=p_x\,\rvirt_x$
for every $x\in\Xset$; and
\item \emph{(retention, an attestation)} $A$ is retained in Alice's
laboratory and never transmitted; no operation of the protocol and no
operation of the adversary acts on $A$, except Alice's declared measurement
$\{M_x\}$.
\end{enumerate}
It is given in one of the three forms below. In the pure form and in the
block form $\Psi$ is declared, together with the declared computational basis
of $A$. In an ensemble presentation neither $\Psi$ nor a basis of $A$ is
declared: $\Psi$ is defined from a declared ensemble by \eqref{eq:enspres}.

\emph{Pure form.} In the \emph{pure form} the sector structure of $A$ is the
one-sector structure, and $\Psi$ is a \emph{rank-one projector} --- $\Psi$
Hermitian, $\Psi^2=\Psi$ and $\tr\Psi=1$, each decided exactly and the trace
summed without rounding, so that $\Psi=\dyad{\psi}$ for a unit vector
$\ket{\psi}\in\Hs{A\Avirt}$.

\emph{Block form.} Let $\Hs{A}=\bigoplus_{t\in\mathbb{T}}\Hs{A_t}$ be the
sector structure of $A$ declared in Definition~\ref{def:source-spec}, along
which $\tau_A=\bigoplus_t\tau^{(t)}_A$ and $M_x=\bigoplus_tM^{(t)}_x$ are
presented as block lists. In the \emph{block form} the virtual transmitted
system carries a declared sector structure \emph{matched} to that of $A$: the
same ordered index set $\mathbb{T}$ and a decomposition
$\Hs{\Avirt}=\bigoplus_{t\in\mathbb{T}}\Hs{S_t}$ into subspaces each spanned
by vectors of the declared basis of $\Avirt$, with projectors
$\Pi^{\mathrm v}_t$. The certificate is a block list
$(\Psi_t)_{t\in\mathbb{T}}$ with $\Psi_t\in\Lin{A_tS_t}$, each block
Hermitian and positive semidefinite, decided exactly, and
$\sum_{t}\tr\Psi_t=1$, the traces summed in the declared order without
rounding. It denotes the operator $\Psi=\bigoplus_t\Psi_t$ on the
\emph{matched subspace} $\bigoplus_t\Hs{A_t}\otimes\Hs{S_t}
\subseteq\Hs{A\Avirt}$, extended by zero to its orthogonal complement. No
purity condition is imposed, on $\Psi$ or on any block. The weighted virtual
members are required to be \emph{block diagonal} along $(S_t)_t$:
$\Pi^{\mathrm v}_t\sigma^{\mathrm v}_x\Pi^{\mathrm v}_{t'}=0$ for all
$t\ne t'$ and every $x$, that is, every entry of $\sigma^{\mathrm v}_x$ whose
row lies in one sector and whose column lies in another is zero, decided entry
by entry on the declared entries of $\sigma^{\mathrm v}_x$ (a virtual ensemble
declared as block lists along $(S_t)_t$ meets it by its presentation). Write
$\sigma^{\mathrm v,(t)}_x\coloneqq\Pi^{\mathrm v}_t\sigma^{\mathrm v}_x
\Pi^{\mathrm v}_t$, read as an operator on $S_t$, for its diagonal blocks.
Against a block list the two identities read block by block,
\begin{equation*}\tag{\getrefnumber{C-eq:srblock}}\digesteqlabel{eq:srblock}
  \tr_{S_t}\Psi_t=\tau^{(t)}_A ,\qquad
  \tr_{A_t}\!\big[(M^{(t)}_x\otimes\idon{S_t})\,\Psi_t\big]
  =\sigma^{\mathrm v,(t)}_x
  \qquad\text{for every }t\in\mathbb{T},\ x\in\Xset ,
\end{equation*}
and \eqref{eq:srblock} together with the block diagonality of every
$\sigma^{\mathrm v}_x$ is equivalent to (SR1) and (SR2) of the operator
$\Psi$. Indeed $\Psi$ vanishes off the matched subspace and $M_x$ is block
diagonal, so the partial traces of $\Psi$ are the direct sums of those of its
blocks: $\tr_{\Avirt}\Psi=\bigoplus_t\tr_{S_t}\Psi_t$, which is
$\tau_A=\bigoplus_t\tau^{(t)}_A$ exactly when the first identity of
\eqref{eq:srblock} holds; and $\tr_A[(M_x\otimes\idon{\Avirt})\Psi]=\bigoplus_t
\tr_{A_t}[(M^{(t)}_x\otimes\idon{S_t})\Psi_t]$, which is block diagonal along
$(S_t)_t$ and so equals $\sigma^{\mathrm v}_x$ exactly when
$\sigma^{\mathrm v}_x$ is block diagonal and its diagonal blocks satisfy the
second identity. Without the block diagonality, \eqref{eq:srblock} is (SR2)
against the pinched members $\bigoplus_t\sigma^{\mathrm v,(t)}_x$ and not
against $\sigma^{\mathrm v}_x$.

\emph{Ensemble presentation.} A block-form certificate need not be declared
at all. In an \emph{ensemble presentation} the declared data are the matched
sector structure $(S_t)_{t\in\mathbb{T}}$ on $\Avirt$ and the weighted
virtual ensemble as block lists,
$\sigma^{\mathrm v}_x=\bigoplus_t\sigma^{\mathrm v,(t)}_x$ with
$\sigma^{\mathrm v,(t)}_x\in\Lin{S_t}$, and \emph{nothing else}: no $\Psi$,
no $\{M_x\}$, no $\tau_A$ and no basis of $A$. The conditions are
\begin{enumerate}[label=(E\arabic*),leftmargin=3.0em]
\item every block $\sigma^{\mathrm v,(t)}_x$ is Hermitian and positive
semidefinite, decided exactly;
\item $p_x=\sum_t\tr\sigma^{\mathrm v,(t)}_x$ for every $x$ and
$\sum_xp_x=1$, the sums formed in the declared orders and rounding nowhere;
and
\item for every $t$ the sum $\sum_{x\in\Xset}\sigma^{\mathrm v,(t)}_x$,
accumulated entry by entry in the declared order of $\Xset$, rounds nowhere,
so that its entries are binary64 numbers.
\end{enumerate}
The system $A$, its sector structure, the marginal $\tau_A$, the POVM
$\{M_x\}$ and the blocks $\Psi_t$ are then \emph{defined} by
\eqref{eq:enspres}: $A_t$ is a copy of $S_t$, $\tau^{(t)}_A$ is the
transpose of the sum in (E3), and $\Psi_t$ and $M^{(t)}_x$ are the canonical
purification and its square-root measurement. By
Lemma~\ref{lem:ensemble-presentation} they are a block-form certificate
with (SR1) and (SR2) holding as identities, with no error term, against the
declared ensemble itself, and with every non-zero block a multiple of a
rank-one projector.
(E1) to (E3) are all that is decided. (SR1), (SR2), positivity of the
blocks, the POVM condition and the purity of the blocks are \emph{theorems}
of the presentation and are decided by nothing, there being no second
declared datum they could be compared with; (SR3) is the attestation it is
in every form. Of the defined objects only $\tau_A$ is ever formed as
numbers, and by (E3) it is admissible declared data. $\Psi_t$ and
$M^{(t)}_x$ carry square roots and quotients of the declared numbers, are
formed by no one, and are never rounded.

\emph{What a consumer reads of a presentation.} What a consumer can read of
Alice's measurement without $\Psi_t$ and $M^{(t)}_x$ is the family
$N^{(t)}_x=(\sigma^{\mathrm v,(t)}_x)^{\transpose}$ of
Lemma~\ref{lem:ensemble-presentation}\,(v). \emph{A consumer whose
declaration carries Alice's measurement as the POVM $\{M_x\}$ cannot receive
an ensemble presentation and refuses it}; the form of declaration that
carries $\{N_x\}$ in its place is that consumer's to define, and the
consuming modules define it: the protocol compiler and the key-rate engine
receive a presentation through the ensemble form of their declarations,
which carries $\{N_x\}$ in the place of $\{M_x\}$
(\xref{protocol-compiler/def:declaration},
\xref{meat-core/def:declaration}), admitted beside a presentation by
\xref{protocol-compiler/def:eligible}. The mode, the alphabet, the source
segment and the attestations of a specification are what they are in the
other forms.

\emph{What is decided.} (SR1) and (SR2) are \emph{exact matrix identities}
on declared data, decided entry by entry. The conditions a consumer decides
--- in pure form the three on $\Psi$, (SR1) and (SR2); in block form
Hermiticity and positive semidefiniteness of every block, the trace, the block
diagonality of every $\sigma^{\mathrm v}_x$ along $(S_t)_t$, and
\eqref{eq:srblock}, which together are (SR1) and (SR2); in an ensemble
presentation (E1) to (E3) and no others
--- are stated in these terms and in no others: none of them is stated on a
vector, which is not a declared datum. (SR3) carries no number and is
decided by nothing; it is the attestation of Assumption~\ref{ass:src}. A
rank-one projector of unit trace is positive semidefinite of unit trace, so
a pure-form certificate satisfies every condition of the block form at the
one-sector structure, and one condition more.

\emph{One reading at every $L_{\mathrm s}$.} The pair
$(\Avirt,\{p_x,\rvirt_x\})$ is the only pair a certificate is ever written
against, at every length of the source segment, because the object a
consumer receives is the virtual protocol
(Definition~\ref{def:source-segment}\,(S3)): the specification's $\tau_A$
and $\{M_x\}$ describe the preparation of $\{p_x,\rvirt_x\}$, and it is the
source identity \eqref{eq:srcmap}, not the certificate, that carries that
description back to the specified ensemble $\{p_x,\rho_x\}$ on $\Asig$. At
$L_{\mathrm s}=0$ the two pairs are literally equal, $\Avirt=\Asig$ and
$\rvirt_x=\rho_x$. So (SR1)--(SR2) have exactly one reading whatever
$L_{\mathrm s}$ is.
\end{digestclaim}

% claim transmitter/lem:ensemble-presentation
\digestlabel{lem:ensemble-presentation}{3.4}\pdfbookmark[2]{Lemma 3.4 Certificate defined by an ensemble presentation}{digest-transmitter/lem:ensemble-presentation}
\begin{mainclaim}
\renewcommand*{\RestatedName}{\hyperref{contract.pdf}{}{claim-transmitter/lem:ensemble-presentation}{Lemma~3.4}}%
\begin{digestclaim}[Certificate defined by an ensemble presentation]
Let the virtual transmitted system carry a sector structure
$\Hs{\Avirt}\supseteq\bigoplus_{t\in\mathbb{T}}\Hs{S_t}$, each $S_t$ spanned
by vectors $(\ket{i}_{S_t})_i$ of the declared basis, and let
$(\sigma^{\mathrm v,(t)}_x)_{t\in\mathbb{T}}$, $x\in\Xset$, be block lists
with $\sigma^{\mathrm v,(t)}_x\in\Pos{S_t}$ and
$\sum_{x}\sum_t\tr\sigma^{\mathrm v,(t)}_x=1$. For every $t$ let $A_t$ be a
copy of $S_t$ with basis $(\ket{i}_{A_t})_i$ in the same order, let
transposition $X\mapsto X^{\transpose}$ be taken in these bases, and put
\begin{equation*}\tag{\getrefnumber{C-eq:enspres}}\digesteqlabel{eq:enspres}
\begin{gathered}
  \tau^{(t)}_{S}\coloneqq\sum_{x\in\Xset}\sigma^{\mathrm v,(t)}_x ,\qquad
  \tau^{(t)}_A\coloneqq\big(\tau^{(t)}_{S}\big)^{\transpose} ,\\
  \ket{\psi_t}\coloneqq\sum_i\ket{i}_{A_t}\otimes
     \big(\tau^{(t)}_{S}\big)^{1/2}\ket{i}_{S_t} ,\qquad
  \Psi_t\coloneqq\dyad{\psi_t} ,\\
  M^{(t)}_x\coloneqq\big(\tau^{(t)}_A\big)^{-1/2}
     \big(\sigma^{\mathrm v,(t)}_x\big)^{\transpose}
     \big(\tau^{(t)}_A\big)^{-1/2}
     +\delta_{x,x_0}\big(\idon{A_t}-P_t\big) ,
\end{gathered}
\end{equation*}
where $P_t$ is the projector onto the support of $\tau^{(t)}_A$, the inverse
square root is taken on that support and is zero off it, and $x_0$ is the
first letter of $\Xset$ in its declared order. Then, with
$A\coloneqq\bigoplus_tA_t$:
\begin{enumerate}[label=(\roman*),leftmargin=2.2em]
\item \emph{(a certificate in block form)} every $\Psi_t$ is positive
semidefinite, $\sum_t\tr\Psi_t=1$, $\{M_x\}_x$ with
$M_x=\bigoplus_tM^{(t)}_x$ is a POVM presented as a block list,
$\tau_A=\bigoplus_t\tau^{(t)}_A$ is a state presented as a block list, and
both identities of \eqref{eq:srblock} hold: $(\Psi_t)_t$ satisfies every
condition Definition~\ref{def:srcert} places on a block-form certificate
for $(A,\mathbb{T},\tau_A,\{M_x\})$ against the virtual ensemble with
weighted members $\sigma^{\mathrm v}_x=\bigoplus_t\sigma^{\mathrm v,(t)}_x$,
(SR3) apart, which is an attestation;
\item \emph{(pure blocks)} every non-zero $\Psi_t$ is
$g_t\coloneqq\tr\tau^{(t)}_S$ times a rank-one projector, so the converse of
Theorem~\ref{thm:certificate-sufficiency}\,(ii) holds of this certificate;
\item \emph{(what is formed as numbers)} the entry $(i,j)$ of
$\tau^{(t)}_A$ is the entry $(j,i)$ of the sum $\tau^{(t)}_S$ of the
declared blocks. No
other object of \eqref{eq:enspres} enters (i) or (ii) through its entries:
$\Psi_t$ and $M^{(t)}_x$ contain square roots and quotients of the declared
numbers and are, in general, not declarable;
\item \emph{(the completion is immaterial)} the letter $x_0$ enters the
POVM condition and nothing else: (SR1), (SR2), (ii) and (v) hold with the
term $\idon{A_t}-P_t$ assigned to any letter, or split among the letters by
any positive semidefinite operators on the kernel of $\tau^{(t)}_A$ that
sum to $\idon{A_t}-P_t$;
\item \emph{(Alice's side without a square root)}
$\big(\tau^{(t)}_A\big)^{1/2}M^{(t)}_x\big(\tau^{(t)}_A\big)^{1/2}
=\big(\sigma^{\mathrm v,(t)}_x\big)^{\transpose}$ for every $x$ and $t$: the
operators $N^{(t)}_x\coloneqq(\sigma^{\mathrm v,(t)}_x)^{\transpose}$ are
entrywise the declared numbers, and
$\sum_xN^{(t)}_x=\tau^{(t)}_A$.
\end{enumerate}
\end{digestclaim}
\end{mainclaim}

\section*{4\quad Sufficiency of the certificate and its scope}\pdfbookmark[1]{4 Sufficiency of the certificate and its scope}{digest-group-transmitter/thm:certificate-sufficiency}

% claim transmitter/thm:certificate-sufficiency
\digestlabel{thm:certificate-sufficiency}{4.1}\pdfbookmark[2]{Theorem 4.1 Checkability and sufficiency of the source-replacement certificate}{digest-transmitter/thm:certificate-sufficiency}
\begin{mainclaim}
\renewcommand*{\RestatedName}{\hyperref{contract.pdf}{}{claim-transmitter/thm:certificate-sufficiency}{Theorem~4.1}}%
\begin{digestclaim}[Checkability and sufficiency of the source-replacement certificate]
Let a source specification (Definition~\ref{def:source-spec}) carry a
source-replacement certificate $\Psi$ for $(A,\mathbb{T},\tau_A,\{M_x\})$
against the virtual ensemble $\{p_x,\rvirt_x\}$, in any form of
Definition~\ref{def:srcert}. In (ii) and (iii) and their proofs an ensemble
presentation is read as the block-form certificate it defines
(Lemma~\ref{lem:ensemble-presentation}\,(i)), so that ``either form'' below
covers it.
\begin{enumerate}[label=(\roman*),leftmargin=2.2em]
\item \emph{(checkability)} the conditions on $\Psi$ of its declared form
--- in pure form Hermiticity, $\Psi^2=\Psi$ and $\tr\Psi=1$; in block form
Hermiticity and positive semidefiniteness of every block and
$\sum_t\tr\Psi_t=1$ --- together with (SR1) and (SR2), are finitely many
exact identities and exact decisions on the declared data, decided by
finitely many arithmetic operations. In block form (SR1) and (SR2) are decided
as the block identities \eqref{eq:srblock} together with the block diagonality
of every weighted virtual member $\sigma^{\mathrm v}_x$ along $(S_t)_t$,
decided entry by entry, which by Definition~\ref{def:srcert} is their
conjunction; \eqref{eq:srblock} alone decides (SR2) only against the pinched
members $\bigoplus_t\Pi^{\mathrm v}_t\sigma^{\mathrm v}_x\Pi^{\mathrm v}_t$.
For a certificate given by an ensemble
presentation the decided conditions are (E1) to (E3) of
Definition~\ref{def:srcert}, equally finite and exact, and the conditions
of the block form hold of the defined certificate by
Lemma~\ref{lem:ensemble-presentation}\,(i) without being decided. (SR3) is
an attestation and is decided by no function of the declared data. (SR1)
moreover implies $\tau_A\in\Seq{A}$, $\tr_{\Avirt}\Psi$ being positive
semidefinite with trace $\tr\Psi=1$. In an ensemble presentation $\tau_A$ is
not declared: it is formed as the transposed sum of (E3), and that it is a
state is Lemma~\ref{lem:ensemble-presentation}\,(i).
\item \emph{(consistency)} more generally than in a source specification,
let $\{M_x\}$ be any family of Hermitian operators on $A$, presented as block
lists against the sector structure, of which $\sum_xM_x=\idon{A}$ is not
assumed; let $\Psi$ meet the conditions on $\Psi$ of its form listed in (i)
and (SR1); and let $p_x\rvirt_x$ stand for
$\tr_A[(M_x\otimes\idon{\Avirt})\Psi]$, so that (SR2) holds. Writing
$U=\operatorname{supp}\tau_A$ with projector $\Pi_U$, the certificate's
signal-side marginal and the virtual average signal state agree,
\begin{equation*}\tag{\getrefnumber{C-eq:srconsist}}\digesteqlabel{eq:srconsist}
  \sum_xp_x\rvirt_x=\tr_A\Psi ,
\end{equation*}
\emph{if} $\Pi_U\big(\sum_xM_x\big)\Pi_U=\Pi_U$, in either form; in
particular they agree for the certificate of a source specification, whose
$\{M_x\}$ is a POVM by Definition~\ref{def:source-spec}. The converse,
\emph{only if}, holds in pure form, and in block form whenever every non-zero
block $\Psi_t$ is a positive multiple of a rank-one projector; it fails for a
general positive block.
\item \emph{(sufficiency at the module boundary)} Assume (SR1), (SR2), the
attestation (SR3) of Assumption~\ref{ass:src}, and
Assumption~\ref{ass:seq}; in block form (SR2) is the identity of operators on
$\Avirt$ against the weighted virtual members $p_x\rvirt_x$ themselves, and
not against their diagonal blocks alone: for a declared block form (i) decides
it with their block diagonality, and for an ensemble presentation it is
proved by Lemma~\ref{lem:ensemble-presentation}\,(i). Let $\pcNmsm$ be Alice's measurement map, from
her retained system $A$ to the classical outcome register $X$,
\begin{equation*}\tag{\getrefnumber{C-eq:Nmsm}}\digesteqlabel{eq:Nmsm}
  \pcNmsm(\xi)\coloneqq\sum_{x\in\Xset}\tr[M_x\,\xi]\;\dyad{x}_{X},
  \qquad \xi\in\Lin{A} ,
\end{equation*}
and let $\mathcal{E}$ be \emph{any} channel from $\Avirt E$ to $E'$, where
$E$ is the adversary's initial register and $E'$ collects her output
registers including whatever is delivered to the receiver. Thus
$\mathcal{E}$ models the adversary's attack and the remainder of the round
\emph{on the transmitted system and the adversary's registers}, and it
acts on neither $A$ nor $X$: it is not permitted to depend on Alice's
outcome. Write
$\rho_{X\Avirt}\coloneqq\sum_{x\in\Xset}p_x\,\dyad{x}_X\otimes\rvirt_x$ for
the virtual prepare-and-measure source state. Then:
\begin{enumerate}[label=(\alph*),leftmargin=2.2em]
\item \emph{(the interchange)}
\begin{equation*}\tag{\getrefnumber{C-eq:srmain}}\digesteqlabel{eq:srmain}
  (\idmap_{X}\otimes\mathcal{E})\circ(\pcNmsm\otimes\idmap_{\Avirt E})
  \;=\;\pcNmsm\otimes\mathcal{E}\;=\;
  (\pcNmsm\otimes\idmap_{E'})\circ(\idmap_{A}\otimes\mathcal{E})
\end{equation*}
as linear maps on $\Lin{A\Avirt E}$, hence in particular on
$\Psi\otimes\sigma$ for every state $\sigma$ on $E$: Alice's measurement may
be performed before the attack or after it, indifferently. This uses no
property of the certificate, only that the two maps act on complementary
registers;
\item \emph{(the source identification)} for every state $\sigma$ on $E$,
\begin{equation*}\tag{\getrefnumber{C-eq:srident}}\digesteqlabel{eq:srident}
  (\pcNmsm\otimes\mathcal{E})\big(\Psi\otimes\sigma\big)
  \;=\;(\idmap_{X}\otimes\mathcal{E})\big(\rho_{X\Avirt}\otimes\sigma\big) ,
\end{equation*}
and both sides are states on $XE'$. Together with (a): on every register
other than $A$, the declared entanglement-based description and the
virtual prepare-and-measure description produce the same state, against
every attack of this form. The states $\rvirt_x$ in $\rho_{X\Avirt}$ are the
declared virtual states, in block form as in every form, and not their
pinching along $(S_t)_t$.
\end{enumerate}
\end{enumerate}
\end{digestclaim}
\end{mainclaim}

% claim transmitter/cor:noncanonical
\digestlabel{cor:noncanonical}{4.3}\pdfbookmark[2]{Corollary 4.3 Certificates beyond the canonical form}{digest-transmitter/cor:noncanonical}
\renewcommand*{\RestatedName}{\hyperref{contract.pdf}{}{claim-transmitter/cor:noncanonical}{Corollary~4.3}}%
\begin{digestclaim}[Certificates beyond the canonical form]
Throughout, the certificate is in the \emph{pure} form of
Definition~\ref{def:srcert}, against the one-sector structure on $A$ that
Definition~\ref{def:source-spec} makes the default.
\begin{enumerate}[label=(\roman*),leftmargin=2.2em]
\item Fact~\ref{fact:sr}'s canonical data --- $A=\Abar\Ahat$ with
$\dim\Hs{\Abar}=|\Xset|$, $M_x=\dyad{x}_{\Abar}\otimes\idon{\Ahat}$, and
$\Psi=\dyad{\psi}$ with
$\ket{\psi}=\sum_x\sqrt{p_x}\ket{x}_{\Abar}\ket{\psi_x}_{\Ahat\Avirt}$ and
$\ket{\psi_x}$ a purification of $\rvirt_x$ --- satisfy (SR1) and (SR2).
\item They are not the only such data: Theorem~\ref{thm:certificate-sufficiency}\,(iii)
applies to every triple $(\Psi,\tau_A,\{M_x\})$ satisfying (SR1)--(SR3), and
there are source
specifications with $d_A<|\Xset|$ carrying a certificate. In particular the
four-state BB84 ensemble on a qubit, $\Xset=\{(b,k):b,k\in\{0,1\}\}$ with
$p_{(b,k)}=\tfrac14$ and $\rvirt_{(b,k)}=\dyad{k_b}$ for the computational
basis ($b=0$) and the Hadamard basis ($b=1$), admits the certificate
$A=\Avirt=\mathbb{C}^2$, $M_{(b,k)}=\tfrac12\dyad{k_b}$,
$\tau_A=\tfrac12\idon{A}$, $\Psi=\dyad{\Phip}$ with
$\ket{\Phip}=\tfrac1{\sqrt2}(\ket{00}+\ket{11})$: $d_A=2$ where the
canonical form has $d_A=4$.
\end{enumerate}
\end{digestclaim}

% claim transmitter/lem:matched-pinching
\digestlabel{lem:matched-pinching}{4.5}\pdfbookmark[2]{Lemma 4.5 Matched pinching of a block certificate}{digest-transmitter/lem:matched-pinching}
\begin{mainclaim}
\renewcommand*{\RestatedName}{\hyperref{contract.pdf}{}{claim-transmitter/lem:matched-pinching}{Lemma~4.5}}%
\begin{digestclaim}[Matched pinching of a block certificate]
Let a source specification carry a certificate in block form
(Definition~\ref{def:srcert}), with sector projectors $\Pi^{A}_t$ on $A$ and
$\Pi^{\mathrm v}_t$ on $\Avirt$, $t\in\mathbb{T}$, which sum to the identity
on $A$ and on $\Avirt$ respectively (Definitions~\ref{def:source-spec}
and~\ref{def:srcert}), and write
\begin{equation*}
  \mathrm{Pin}_{A}(\xi)\coloneqq\sum_{t\in\mathbb{T}}\Pi^{A}_t\,\xi\,\Pi^{A}_t ,
  \qquad
  \mathrm{Pin}_{\mathrm v}(\zeta)\coloneqq
  \sum_{t\in\mathbb{T}}\Pi^{\mathrm v}_t\,\zeta\,\Pi^{\mathrm v}_t
\end{equation*}
for the two pinching channels.
\begin{enumerate}[label=(\roman*),leftmargin=2.2em]
\item \emph{(matched structure)} for every $t\in\mathbb{T}$,
\begin{equation*}\tag{\getrefnumber{C-eq:matched}}\digesteqlabel{eq:matched}
  (\Pi^{A}_t\otimes\idon{\Avirt})\,\Psi
  =\Psi\,(\Pi^{A}_t\otimes\idon{\Avirt})
  =(\idon{A}\otimes\Pi^{\mathrm v}_t)\,\Psi
  =\Psi\,(\idon{A}\otimes\Pi^{\mathrm v}_t)
  =\Psi_t ,
\end{equation*}
the block $\Psi_t$ read as an operator on $\Hs{A\Avirt}$ extended by zero.
Hence
$(\mathrm{Pin}_{A}\otimes\idmap_{\Avirt})(\Psi)=\Psi
=(\idmap_{A}\otimes\mathrm{Pin}_{\mathrm v})(\Psi)$.
\item \emph{(reached states)} assume (SR3), and let $\mathcal{E}$ be any
channel from $\Avirt E$ to $E'$ and $\sigma$ any state on $E$, as in
Theorem~\ref{thm:certificate-sufficiency}\,(iii). The state
$\omega\coloneqq(\idmap_{A}\otimes\mathcal{E})(\Psi\otimes\sigma)$ on $AE'$
satisfies
\begin{equation*}\tag{\getrefnumber{C-eq:reached}}\digesteqlabel{eq:reached}
  \tr_{E'}\omega=\tau_A ,\qquad
  (\mathrm{Pin}_{A}\otimes\idmap_{E'})(\omega)=\omega ,\qquad
  \omega=(\idmap_{A}\otimes\mathcal{E}')(\Psi\otimes\sigma)
  \ \text{ with }\
  \mathcal{E}'\coloneqq\mathcal{E}\circ(\mathrm{Pin}_{\mathrm v}\otimes\idmap_{E}) .
\end{equation*}
So $\omega=\bigoplus_t\omega_t$ with $\omega_t\in\Pos{A_tE'}$ and
$\tr_{E'}\omega_t=\tau^{(t)}_A$; and the attack $\mathcal{E}$ and the attack
$\mathcal{E}'$ that first measures the sector of the virtual signal reach
the same state.
\item \emph{(operators that are not reached)} the block diagonality of (ii)
does not extend to operators merely constrained to have the marginal: in
general a positive operator on $AE'$ with marginal
$\tau_A=\bigoplus_t\tau^{(t)}_A$ need not be block diagonal along
$\mathbb{T}$, and this happens already for two one-dimensional sectors,
$\tau_A=\tfrac12\idon{A}$ and $E'$ a qubit.
\end{enumerate}
\end{digestclaim}
\end{mainclaim}

\section*{5\quad Assumptions}\pdfbookmark[1]{5 Assumptions}{digest-group-transmitter/ass:ens}

% claim transmitter/ass:ens
\digestlabel{ass:ens}{5.1}\pdfbookmark[2]{Assumption 5.1 Emission of the specified ensemble}{digest-transmitter/ass:ens}
\SetShortId{A-ENS}
\renewcommand*{\RestatedName}{\hyperref{contract.pdf}{}{claim-transmitter/ass:ens}{Assumption~5.1}}%
\begin{digestassumption}[Emission of the specified ensemble]
For the declared signal ensemble $\{p_x,\rho_x\}$ of
Definition~\ref{def:source-spec} on the transmitted system, the
letters the physical source draws and the systems it emits satisfy the
following; a letter that Alice obtains in untrusted-source mode by measuring a
received system is not drawn by the source, and nothing below concerns it:
(i) the letters $x_1,\dots,x_n$ it draws in the $n$ rounds are drawn
independently, each with the weights $p_x$;
(ii) given those letters, the systems it emits in the $n$ rounds are in the
state $\rho_{x_1}\otimes\dots\otimes\rho_{x_n}$;
(iii) neither the letter it draws in a round nor the state it emits in it
depends on what happens in the other rounds, their outcomes and public
messages included.
\end{digestassumption}

% claim transmitter/ass:seq
\digestlabel{ass:seq}{5.2}\pdfbookmark[2]{Assumption 5.2 Sequential round structure}{digest-transmitter/ass:seq}
\SetShortId{A-SEQ}
\renewcommand*{\RestatedName}{\hyperref{contract.pdf}{}{claim-transmitter/ass:seq}{Assumption~5.2}}%
\begin{digestassumption}[Sequential round structure]
The physical protocol has the sequential structure the declared protocol
describes:
(i) each round runs, in this order, state preparation and transmission; the
measurements of the round, by every party that measures in it; the public
messages of the round, in their declared order, \emph{each of them sent after
every measurement of the round is complete}, whichever party sends it and
whichever party measured; and announcement-determined sifting;
(ii) the classical post-processing is applied after the $n$ rounds;
(iii) every message is public;
(iv) no operation on a quantum system of round $j$ depends on a public
message of round $j$;
(v) every party that measures makes its measurements in round order: its
measurement of round $j$ is complete before its measurement of round $j+1$
is.
\end{digestassumption}

% claim transmitter/ass:src
\digestlabel{ass:src}{5.3}\pdfbookmark[2]{Assumption 5.3 Retention of the system A in trusted-source mode}{digest-transmitter/ass:src}
\SetShortId{A-SRC}
\renewcommand*{\RestatedName}{\hyperref{contract.pdf}{}{claim-transmitter/ass:src}{Assumption~5.3}}%
\begin{digestassumption}[Retention of the system A in trusted-source mode]
For a source specification (Definition~\ref{def:source-spec}) in
trusted-source mode, which fills the marginal slot with
$\tau_A\in\Seq{A}$, condition~(SR3) of Definition~\ref{def:srcert} holds of
the apparatus:
(i) the system $A$ is retained in the laboratory of the party whose label
the party partition of the consuming declaration
(\xref{protocol-compiler/def:declaration}) gives $A$, and is never
transmitted;
(ii) no operation of the protocol or of the adversary acts on $A$ except the
declared measurement $\{M_x\}$ of that party.
\end{digestassumption}

\part*{Engine requirements}\pdfbookmark[0]{Engine requirements}{digest-part-2}

\section*{6\quad Certificate and licence requirements}\pdfbookmark[1]{6 Certificate and licence requirements}{digest-group-transmitter/req:licence-scope}

% claim transmitter/req:licence-scope
\digestlabel{req:licence-scope}{6.1}\pdfbookmark[2]{Requirement 6.1 Certificate and attestation demanded only at d\_A≥2}{digest-transmitter/req:licence-scope}
\renewcommand*{\RestatedName}{\hyperref{contract.pdf}{}{claim-transmitter/req:licence-scope}{Requirement~6.1}}%
\begin{digestrequirement}[Certificate and attestation demanded only at $d_A\ge2$]
The engine demands the source-replacement certificate of
Definition~\ref{def:srcert} of a request in trusted-source mode only at
$d_A\ge2$, and there refuses such a request that carries none
(\iface{trusted_mode_without_certificate}). At $d_A=1$, where the licensed
condition $\omega_A=\tau_A$ excludes no state (Remark~\ref{C-rem:da1}),
neither the engine nor a consumer of its artefact demands the certificate or
the attestation of Assumption~\ref{ass:src}.
\end{digestrequirement}

% claim transmitter/req:marginal-fill
\digestlabel{req:marginal-fill}{6.2}\pdfbookmark[2]{Requirement 6.2 The unique marginal served at d\_A=1}{digest-transmitter/req:marginal-fill}
\renewcommand*{\RestatedName}{\hyperref{contract.pdf}{}{claim-transmitter/req:marginal-fill}{Requirement~6.2}}%
\begin{digestrequirement}[The unique marginal served at $d_A=1$]
At $d_A=1$ a request in trusted-source mode may leave the marginal slot
empty. The engine then fills the slot with $\tau_A=[1]$, the unique element
of $\Seq{A}$ (Remark~\ref{C-rem:da1}); the emitted specification is in
trusted-source mode, trusted-source and untrusted-source mode imposing the
same condition there; and the emitted artefact records that the slot was
filled. This is the only slot the engine ever fills.
\end{digestrequirement}

% claim transmitter/req:failing-certificate
\digestlabel{req:failing-certificate}{6.3}\pdfbookmark[2]{Requirement 6.3 Refusal of a certificate that fails a decided condition}{digest-transmitter/req:failing-certificate}
\renewcommand*{\RestatedName}{\hyperref{contract.pdf}{}{claim-transmitter/req:failing-certificate}{Requirement~6.3}}%
\begin{digestrequirement}[Refusal of a certificate that fails a decided condition]
The engine refuses a request whose certificate is present and fails one of
the conditions Definition~\ref{def:srcert} decides of its form --- in pure
form the three on $\Psi$, (SR1) and (SR2); in block form Hermiticity and
positive semidefiniteness of every block, the trace, the block diagonality of
every weighted virtual member along the matched sectors, and
\eqref{eq:srblock}; in an ensemble presentation (E1) to (E3) --- at every
$d_A$, including $d_A=1$, where no certificate is demanded
(Requirement~\ref{req:licence-scope}): a certificate that does not verify is
a defect in the supplied data, not a weaker claim.
\end{digestrequirement}

% claim transmitter/req:marginal-decided
\digestlabel{req:marginal-decided}{6.4}\pdfbookmark[2]{Requirement 6.4 Direct decision of the declared marginal}{digest-transmitter/req:marginal-decided}
\renewcommand*{\RestatedName}{\hyperref{contract.pdf}{}{claim-transmitter/req:marginal-decided}{Requirement~6.4}}%
\begin{digestrequirement}[Direct decision of the declared marginal]
Where $\tau_A$ is declared, the engine decides $\tau_A\in\Seq{A}$ on the
declared matrix directly, also where a certificate is present and satisfies
(SR1), and never infers it from the implication
(SR1)\,$\Rightarrow\tau_A\in\Seq{A}$ of Theorem~\ref{thm:certificate-sufficiency}\,(i); it
refuses a request whose declared $\tau_A$ is not a state
(\iface{alice_marginal_not_admissible}). Consumers of the declared data
decide it in the same way under their own rules; this requirement binds the
engine.
\end{digestrequirement}

% claim transmitter/req:attestation-record
\digestlabel{req:attestation-record}{6.5}\pdfbookmark[2]{Requirement 6.5 Pass-through of the retention attestation}{digest-transmitter/req:attestation-record}
\renewcommand*{\RestatedName}{\hyperref{contract.pdf}{}{claim-transmitter/req:attestation-record}{Requirement~6.5}}%
\begin{digestrequirement}[Pass-through of the retention attestation]
Where the engine emits the licence record \iface{source_replacement_license},
which it does only in trusted-source mode at $d_A\ge2$, the scope in
which the certificate is demanded (Requirement~\ref{req:licence-scope}), it
records the (SR3) attestation it receives (Definition~\ref{def:srcert}) in that
record, and passes it through: it never
verifies it and never asserts it of any physical apparatus
(Assumption~\ref{ass:src}; Remark~\ref{C-rem:boundary}). No consumer can
verify it either, since no function of the declared data decides it
(Theorem~\ref{thm:certificate-sufficiency}\,(i)).
\end{digestrequirement}

\part*{Decoy-state sources and the tagging channel}\pdfbookmark[0]{Decoy-state sources and the tagging channel}{digest-part-3}

\section*{7\quad The source model and the tagged declaration}\pdfbookmark[1]{7 The source model and the tagged declaration}{digest-group-transmitter/decoy/def:wcp-model}

% claim transmitter/decoy/def:wcp-model
\digestlabel{def:wcp-model}{7.1}\pdfbookmark[2]{Definition 7.1 Phase-randomised weak-coherent-pulse source model}{digest-transmitter/decoy/def:wcp-model}
\renewcommand*{\RestatedName}{\hyperref{contract.pdf}{}{claim-transmitter/decoy/def:wcp-model}{Definition~7.1}}%
\begin{digestclaim}[Phase-randomised weak-coherent-pulse source model]
A \emph{weak-coherent-pulse model}
is a source model, the case (b) of
Definition~\ref{def:source-spec} in which the specified ensemble
$\{p_x,\rho_x\}$ is fixed as mathematical objects rather than as matrices, on
a transmitted system whose Hilbert space is separable. It consists of the
following declared data.
\begin{enumerate}[label=(W\arabic*),leftmargin=3.0em]
\item \emph{(modes and polarisations)} two declared optical modes with a fixed
orthonormal basis of the one-photon space $\mathbb{C}^2$, and a finite
non-empty set $\mathrm{Pol}$ of
\emph{polarisations},
declared labels each standing for a unit vector $\ket{a}\in\mathbb{C}^2$,
$a\in\mathrm{Pol}$, with creation operator $c^{\dagger}_a$. The vector
$\ket{a}$ is a mathematical object of the model, as the states of (W4) are,
and not an array of declared numbers: a declaration names the vector a label
stands for and carries none of its coordinates, which need not be binary64
numbers ($(\ket{H}+\ket{V})/\sqrt2$ has the coordinates $1/\sqrt2$). The
transmitted system $\Asig$ is the Fock space over the two modes, with vacuum
$\ket{\vac}$; its Hilbert space is separable and infinite dimensional.
\item \emph{(intensities and labels)} a finite non-empty set
$\mathrm{Int}\subset[0,\infty)$ of
\emph{intensities};
the signal alphabet
$\Xset\subseteq\mathrm{Pol}\times\mathrm{Int}\times\{\tst,\gen\}$ with its
declared order, an element of
$\mathrm{Pol}\times\mathrm{Int}\times\{\tst,\gen\}$ being a label and a label in
$\Xset$ a letter, written $x=(a_x,\mu_x,z_x)$, whose third component is the
\emph{round type};
and weights $p_x>0$ for every $x\in\Xset$ with $\sum_{x\in\Xset}p_x=1$ in the
declared order, rounding nowhere. A label
of weight zero is not a letter of $\Xset$.
\item \emph{(number states)} for $a\in\mathrm{Pol}$ and an integer
$\mathsf{n}\ge0$, the unit vector
\begin{equation*}\tag{\getrefnumber{C-eq:numberstate}}\digesteqlabel{eq:numberstate}
  \ket{\mathsf{n},a}\;\coloneqq\;
  \frac{(c^{\dagger}_a)^{\mathsf{n}}}{\sqrt{\mathsf{n}!}}\,\ket{\vac} ,
\end{equation*}
and the subspace
$S_{\mathsf{n}}\coloneqq\operatorname{span}\{\ket{\mathsf{n},a}:a\in\mathrm{Pol}\}$
of the $\mathsf{n}$-photon subspace. Since the $\mathsf{n}$-photon subspace of
a two-mode Fock space has dimension $\mathsf{n}+1$ and $S_{\mathsf{n}}$ is
spanned by $|\mathrm{Pol}|$ vectors,
$\dim S_{\mathsf{n}}\le\min(\mathsf{n}+1,|\mathrm{Pol}|)$; and
$S_{\mathsf{n}}\perp S_{\mathsf{n}'}$ for $\mathsf{n}\ne\mathsf{n}'$, the two
subspaces lying in different photon-number eigenspaces.
$S_0=\mathbb{C}\ket{\vac}$ is one dimensional, every polarisation sharing the
vacuum.
\item \emph{(the specified states)} with the
\emph{Poisson weight}
$P_{\mu}(\mathsf{n})\coloneqq e^{-\mu}\mu^{\mathsf{n}}/\mathsf{n}!$, the state
prepared at the letter $x=(a,\mu,z)$ is
\begin{equation*}\tag{\getrefnumber{C-eq:wcpstate}}\digesteqlabel{eq:wcpstate}
  \rho_x\;=\;\sum_{\mathsf{n}\ge0}P_{\mu}(\mathsf{n})\,
  \dyad{\mathsf{n},a} ,
\end{equation*}
which does not depend on the round type $z$. It is a state:
$\sum_{\mathsf{n}\ge0}P_{\mu}(\mathsf{n})=1$ and each $\dyad{\mathsf{n},a}$
is a rank-one projector.
\end{enumerate}
The physical-source assumption of the model, which asserts that the device
emits \eqref{eq:wcpstate}, is Assumption~\ref{ass:wcp}; the model declares it
in the place of Assumption~\ref{ass:ens}.
\end{digestclaim}

% claim transmitter/decoy/def:tagged-source
\digestlabel{def:tagged-source}{7.2}\pdfbookmark[2]{Definition 7.2 Tagged virtual source: cutoffs, flags and sector map}{digest-transmitter/decoy/def:tagged-source}
\renewcommand*{\RestatedName}{\hyperref{contract.pdf}{}{claim-transmitter/decoy/def:tagged-source}{Definition~7.2}}%
\begin{digestclaim}[Tagged virtual source: cutoffs, flags and sector map]
Let a weak-coherent-pulse model be declared (Definition~\ref{def:wcp-model}). A
\emph{tagged source declaration}
over it consists of the following.
\begin{enumerate}[label=(T\arabic*),leftmargin=3.0em]
\item \emph{(the two cutoffs)} integers $N_{\mathrm{ent}}$ and
$N_{\mathrm{src}}$ with $0\le N_{\mathrm{ent}}\le N_{\mathrm{src}}$: the
\emph{entropy cutoff},
the largest photon number whose sector stays quantum, and the
\emph{source cutoff},
the largest photon number that keeps a sector of its own.
\item \emph{(the sector weights)} declared numbers $\tilde p(\mathsf{n}\mid\mu)$
for $0\le\mathsf{n}\le N_{\mathrm{src}}$ and $\mu\in\mathrm{Int}$, depending on
the intensity alone and not on the polarisation or the round type, satisfying
\begin{equation*}\tag{\getrefnumber{C-eq:weightineq}}\digesteqlabel{eq:weightineq}
  0\;\le\;\tilde p(\mathsf{n}\mid\mu)\;\le\;P_{\mu}(\mathsf{n})
  \qquad\text{for every }\mathsf{n}\le N_{\mathrm{src}},\ \mu\in\mathrm{Int} ,
\end{equation*}
and the
\emph{tail weight}
\begin{equation*}\tag{\getrefnumber{C-eq:tailweight}}\digesteqlabel{eq:tailweight}
  \tilde w_{\mu}\;\coloneqq\;
  1-\sum_{\mathsf{n}=0}^{N_{\mathrm{src}}}\tilde p(\mathsf{n}\mid\mu) ,
\end{equation*}
the sum formed in increasing $\mathsf{n}$ and rounding nowhere. Then
$\tilde w_{\mu}\ge0$, because
$\sum_{\mathsf{n}\le N_{\mathrm{src}}}\tilde p(\mathsf{n}\mid\mu)
\le\sum_{\mathsf{n}\le N_{\mathrm{src}}}P_{\mu}(\mathsf{n})\le1$ by
\eqref{eq:weightineq}.
\item \emph{(the sector index set and the virtual transmitted system)} a
sector that no letter of $\Xset$ puts positive weight on is not declared: the
\emph{sector index set}
carries the sectors the declaration uses and is read off the weights,
\begin{equation*}\tag{\getrefnumber{C-eq:tset}}\digesteqlabel{eq:tset}
\begin{aligned}
  \mathbb{T}\;=\;
  &\{\mathsf{n}\le N_{\mathrm{ent}}:\
     \tilde p(\mathsf{n}\mid\mu_x)>0\text{ for some }x\in\Xset\}\\
  \cup\;&\{(\mathsf{n},a):N_{\mathrm{ent}}<\mathsf{n}\le N_{\mathrm{src}},\
     \tilde p(\mathsf{n}\mid\mu_x)>0\text{ for some }x\in\Xset
     \text{ with }a_x=a\}\\
  \cup\;&\{(a,\mu):\tilde w_{\mu}>0\text{ and }(a,\mu)=(a_x,\mu_x)
     \text{ for some }x\in\Xset\} ,
\end{aligned}
\end{equation*}
in that declared order: the quantum sectors in increasing photon number, then
the polarisation flags in the declared order of $(\mathsf{n},a)$, then the tail
flags in the declared order of $(a,\mu)$. The finite-dimensional virtual
transmitted system is the orthogonal direct sum over it,
\begin{equation*}\tag{\getrefnumber{C-eq:avirt}}\digesteqlabel{eq:avirt}
  \Avirt\;=\;\bigoplus_{t\in\mathbb{T}}S_t ,\qquad
  S_t=\begin{cases}
    S_{\mathsf{n}} & t=\mathsf{n}\le N_{\mathrm{ent}},\\
    \mathbb{C}f_{\mathsf{n},a} & t=(\mathsf{n},a),\\
    \mathbb{C}f_{a,\mu} & t=(a,\mu),
  \end{cases}
\end{equation*}
with the basis of each summand described next; it is the declared basis of
$\Avirt$ in the sense of Definition~\ref{def:srcert}, the basis in which
the blocks of the declaration are written. At a quantum sector
$t=\mathsf{n}$ the summand is the space $S_{\mathsf{n}}$ of
Definition~\ref{def:wcp-model}\,(W3), taken into the direct sum with its
vectors $\ket{\mathsf{n},a}$ and an orthonormal basis of $S_{\mathsf{n}}$;
like those vectors, that basis is a mathematical object of the model and not
an array of declared numbers (Proposition~\ref{prop:tagging} returns
$S_{\mathsf{n}}$ to the Fock space by an isometry $V_{\mathsf{n}}$). There is
one
\emph{polarisation flag}
$f_{\mathsf{n},a}$ per photon number above the entropy cutoff and per
polarisation of a letter that weights it, shared across intensities and round
types; and one
\emph{tail flag}
$f_{a,\mu}$ per pair a letter carries with
$\tilde w_{\mu}>0$, shared across round types. The flag vectors are declared
orthonormal, the summands of \eqref{eq:avirt} are mutually orthogonal in
$\Avirt$, and $\Pi^{\mathrm v}_t$ is the projector onto $S_t$. Every flag
sector is one dimensional, and so is the quantum sector $\mathsf{n}=0$. The
index set is a function of the declared weights.
\item \emph{(the virtual ensemble)} for $x=(a,\mu,z)\in\Xset$,
\begin{equation*}\tag{\getrefnumber{C-eq:rvirt}}\digesteqlabel{eq:rvirt}
  \rvirt_x\;=\;
  \sum_{\mathsf{n}\le N_{\mathrm{ent}}}
    \tilde p(\mathsf{n}\mid\mu)\,\dyad{\mathsf{n},a}
  \;+\;
  \sum_{N_{\mathrm{ent}}<\mathsf{n}\le N_{\mathrm{src}}}
    \tilde p(\mathsf{n}\mid\mu)\,\dyad{f_{\mathsf{n},a}}
  \;+\;
  \tilde w_{\mu}\,\dyad{f_{a,\mu}} ,
\end{equation*}
every sum running over the sectors of the index set \eqref{eq:tset} only: a
quantum sector or a polarisation flag that no letter puts positive weight on is
not a sector of $\Avirt$ by (T3) and contributes no term, as the tail flag
contributes none where $\tilde w_{\mu}=0$. Each $\rvirt_x$ is a state on
$\Avirt$ that is block diagonal along $\mathbb{T}$. Its component in a quantum
sector $\mathsf{n}\le N_{\mathrm{ent}}$ is
$\tilde p(\mathsf{n}\mid\mu)\dyad{\mathsf{n},a}$; in the polarisation flag
sector $(\mathsf{n},a')$ it is $\tilde p(\mathsf{n}\mid\mu)$ if $a'=a$ and zero
otherwise; in the tail flag sector $(a',\mu')$ it is $\tilde w_{\mu}$ if
$(a',\mu')=(a,\mu)$ and zero otherwise. It is positive semidefinite, being a
sum of non-negative multiples of rank-one projectors, and its trace is
$\sum_{\mathsf{n}\le N_{\mathrm{src}}}\tilde p(\mathsf{n}\mid\mu)
+\tilde w_{\mu}=1$ by \eqref{eq:tailweight}, the absent terms contributing
nothing to it: a sector that \eqref{eq:tset} leaves out is one on which
$\tilde p(\mathsf{n}\mid\mu)$ vanishes for every letter, this one included, or
a tail flag with $\tilde w_{\mu}=0$. Two labels of the same polarisation and
different intensities put weight on the same polarisation flag, and two labels
of the same polarisation and intensity and different round types put weight on
the same tail flag.
\item \emph{(the matched sector map)} the retained system and its
sector structure are not declared: a tagged source is presented as an
ensemble presentation (Definition~\ref{def:srcert}: a certificate defined from declared
weighted blocks rather than declared), in which $A$, its structure, $\tau_A$
and $\{M_x\}$ are the objects \eqref{eq:enspres} defines, $A_t$ being a copy of
$S_t$ sector by sector and in the same order. In particular
$\dim\Hs{A_t}=\dim\Hs{S_t}$ whatever the rank of the sector's marginal
(Proposition~\ref{C-prop:decoy-cert}\,(iii)). The matched pair
$(\{\Pi^{A}_t\},\{\Pi^{\mathrm v}_t\})$, with $\Pi^{A}_t$ the projectors onto
Alice's sectors (Definition~\ref{def:source-spec}), is the structure of (T3)
read on both sides. It is the retained-register sector map of the link of
(T6), the slot (e) of Definition~\ref{def:source-link}, which fixes how the
certificate, $\tau_A$ and $\{M_x\}$ are presented, sector by sector. It is a
statement about those three objects only: it does not assert that any joint
state of $A$ and another system is block diagonal
(Definition~\ref{def:source-spec}).
\item \emph{(the link)} the channel $\Xi_{\mathrm{tag}}$ that
Proposition~\ref{prop:tagging} constructs and proves to carry each $\rvirt_x$
to $\rho_x$, recorded as one source link
(Definition~\ref{def:source-link}) with the following slots: the kind
\textsf{tagging} (slot (a)); the source descriptor
$(\Avirt,\{p_x,\rvirt_x\})$ of (T3) and (T4) and the target descriptor
$(\Asig,\{p_x,\rho_x\})$ of the model (slot (b)); the channel witness supplied
as a theorem constructor, namely Proposition~\ref{prop:tagging}, with an
obligation record that names Assumption~\ref{ass:wcp} (slot (c2)); the
declared descriptor family, every tagged source declaration over every
weak-coherent-pulse model (slot (d)); the retained-register sector map, the
matched pair of (T5) (slot (e)); an empty list of
transported constraint entries, so that no constraint is carried
across the link, and hence no licence record, a licence being present only
beside a non-empty list (slot (f)); and the status \textsf{exact}, the identity
holding with no error term (slot (g)).
\end{enumerate}
The cutoffs of (T1), the weights $\tilde p(\mathsf{n}\mid\mu)$ of (T2), and the
model's alphabet, polarisation labels, intensities and label weights are
\emph{declared}. The polarisation vectors $\ket{a}$, the number states
$\ket{\mathsf{n},a}$ and the orthonormal basis of each quantum sector
$S_{\mathsf{n}}$ of (T3) are objects of the model and are not declared
numbers. The tail weight \eqref{eq:tailweight}, the index set \eqref{eq:tset}
and the weighted members $p_x\rvirt_x$ of \eqref{eq:rvirt} are
\emph{formed}: the tail weight and the index set from the declared numbers in
exact arithmetic, and each entry of a weighted member, in the basis of its
sector, as the product of declared numbers with an entry of the dyad of the
sector's signal vector in that basis, a number fixed by the model and the
basis. They are the declaration's own data: read sector by sector, the
weighted members are the blocks that the ensemble presentation of (T5) takes as
its data. Everything on Alice's side --- $A$, $\tau_A$, $\{M_x\}$ and the
certificate --- is defined from those blocks and declared by no one.
\end{digestclaim}

\section*{8\quad The tagging channel}\pdfbookmark[1]{8 The tagging channel}{digest-group-transmitter/decoy/prop:tagging}

% claim transmitter/decoy/prop:tagging
\digestlabel{prop:tagging}{8.1}\pdfbookmark[2]{Proposition 8.1 Tagging channel for weights at or below the Poisson weights}{digest-transmitter/decoy/prop:tagging}
\begin{mainclaim}
\renewcommand*{\RestatedName}{\hyperref{contract.pdf}{}{claim-transmitter/decoy/prop:tagging}{Proposition~8.1}}%
\begin{digestclaim}[Tagging channel for weights at or below the Poisson weights]
Let a weak-coherent-pulse model be given
(Definition~\ref{def:wcp-model}): labels $x=(a,\mu,z)$ with $z$ the
round type, unit polarisation vectors
$\ket{a}$, Poisson weights
$P_\mu(\mathsf{n})=e^{-\mu}\mu^{\mathsf{n}}/\mathsf{n}!$ and model states
$\rho_x=\sum_{\mathsf{n}\ge0}P_\mu(\mathsf{n})\,\dyad{\mathsf{n},a}$ on the
two-mode Fock space. Let a tagged source be declared on it
(Definition~\ref{def:tagged-source}) with cutoffs
$0\le N_{\mathrm{ent}}\le N_{\mathrm{src}}$, virtual states $\rvirt_x$ of
\eqref{eq:rvirt} and declared weights $\tilde p(\mathsf{n}\,|\,\mu)$,
$\mathsf{n}\le N_{\mathrm{src}}$, satisfying the inequality
\eqref{eq:weightineq} of Definition~\ref{def:tagged-source}\,(T2), which
reads
\begin{equation*}\tag{\getrefnumber{C-eq:tag-ineq}}\digesteqlabel{eq:tag-ineq}
  0\le\tilde p(\mathsf{n}\,|\,\mu)\le P_\mu(\mathsf{n})
  \qquad\text{for every intensity }\mu\text{ and every }
  \mathsf{n}\le N_{\mathrm{src}} ,
\end{equation*}
with tail weights
$\tilde w_\mu=1-\sum_{\mathsf{n}\le N_{\mathrm{src}}}\tilde p(\mathsf{n}\,|\,\mu)$.
Then:
\begin{enumerate}[label=(\roman*),leftmargin=2.2em]
\item \emph{(the tail states)} for every polarisation $a$ and every intensity
$\mu$ with $\tilde w_\mu>0$,
\begin{equation*}\tag{\getrefnumber{C-eq:tailstate}}\digesteqlabel{eq:tailstate}
  \vartheta_{a,\mu}\coloneqq\frac{1}{\tilde w_\mu}\Big[
  \sum_{\mathsf{n}\le N_{\mathrm{src}}}
  \big(P_\mu(\mathsf{n})-\tilde p(\mathsf{n}\,|\,\mu)\big)\dyad{\mathsf{n},a}
  +\sum_{\mathsf{n}>N_{\mathrm{src}}}P_\mu(\mathsf{n})\,\dyad{\mathsf{n},a}
  \Big]
\end{equation*}
is a state on the Fock space. Where $\tilde w_\mu=0$ the declaration carries no
tail flag for $(a,\mu)$ (Definition~\ref{def:tagged-source}) and no tail state
is defined.
\item \emph{(one channel)} the map
\begin{equation*}\tag{\getrefnumber{C-eq:xitag}}\digesteqlabel{eq:xitag}
\begin{aligned}
  \Xi_{\mathrm{tag}}(\zeta)\coloneqq{}&
  \sum_{\mathsf{n}\le N_{\mathrm{ent}}}
    V_{\mathsf{n}}\,\Pi^{\mathrm v}_{\mathsf{n}}\,\zeta\,
    \Pi^{\mathrm v}_{\mathsf{n}}V_{\mathsf{n}}^{\dagger}
  +\sum_{N_{\mathrm{ent}}<\mathsf{n}\le N_{\mathrm{src}}}\sum_{a}
    \bra{f_{\mathsf{n},a}}\zeta\ket{f_{\mathsf{n},a}}\,\dyad{\mathsf{n},a}\\
  &+\sum_{(a,\mu):\,\tilde w_\mu>0}
    \bra{f_{a,\mu}}\zeta\ket{f_{a,\mu}}\,\vartheta_{a,\mu} ,
\end{aligned}
\end{equation*}
in which $\Pi^{\mathrm v}_{\mathsf{n}}$ projects $\Hs{\Avirt}$ onto the quantum
sector $S_{\mathsf{n}}$, $V_{\mathsf{n}}$ is the isometry that embeds that
sector into the Fock space as $\operatorname{span}\{\ket{\mathsf{n},a}\}_a$,
and every sum runs over the sectors of the declared index set $\mathbb{T}$
only, is a channel from $\Avirt$ to the Fock space, and
\begin{equation*}\tag{\getrefnumber{C-eq:tagid}}\digesteqlabel{eq:tagid}
  \Xi_{\mathrm{tag}}(\rvirt_x)=\rho_x\qquad\text{for every label }x .
\end{equation*}
A quantum sector or a polarisation flag that no letter puts weight on is not a
sector of $\Avirt$ by Definition~\ref{def:tagged-source}\,(T3) and contributes
no term to \eqref{eq:xitag}, exactly as for the tail flags.
The channel $\Xi_{\mathrm{tag}}$ depends on the declared weights and the
cutoffs and on no label.
\item \emph{(the role of the inequality)} \eqref{eq:tagid} holds for every
family of weights satisfying \eqref{eq:tag-ineq}. Where \eqref{eq:tag-ineq}
fails at an intensity $\mu$, its two halves fail differently.
\begin{enumerate}[label=(\alph*),leftmargin=2.2em]
\item If $\tilde p(\mathsf{n}\,|\,\mu)>P_\mu(\mathsf{n})$ for some
$\mathsf{n}\le N_{\mathrm{src}}$ and $\tilde w_\mu>0$, the operator
\eqref{eq:tailstate} is not positive semidefinite, for every polarisation $a$,
and the construction is not available: where a letter $x$ has intensity $\mu$,
the map \eqref{eq:xitag} sends the state $\dyad{f_{a_x,\mu}}$ to
$\vartheta_{a_x,\mu}$ and is not positive.
\item If $\tilde p(\mathsf{n}\,|\,\mu)<0$ for some
$\mathsf{n}\le N_{\mathrm{src}}$ while
$\tilde p(\mathsf{n}'\,|\,\mu)\le P_\mu(\mathsf{n}')$ for every
$\mathsf{n}'\le N_{\mathrm{src}}$, then $\tilde w_\mu>0$ and the operator
\eqref{eq:tailstate} is a state for every polarisation $a$. What fails is the
virtual state: for every letter $x$ of intensity $\mu$ the operator $\rvirt_x$
of \eqref{eq:rvirt} is not positive semidefinite or has trace greater than
one, so it is not a state.
\end{enumerate}
\item \emph{(the published case)} for
$\tilde p(\mathsf{n}\,|\,\mu)=P_\mu(\mathsf{n})$ and
$N_{\mathrm{ent}}=N_{\mathrm{src}}$ the virtual states are those of
\cite[Eq.~(145)]{tupkary2026decoy} and \eqref{eq:tagid} is
\cite[Lemma~9.2]{tupkary2026decoy}.
\end{enumerate}
\end{digestclaim}
\end{mainclaim}

\section*{9\quad The certificate and the emitted declaration}\pdfbookmark[1]{9 The certificate and the emitted declaration}{digest-group-transmitter/decoy/prop:decoy-cert}

% claim transmitter/decoy/cor:emitted-decoy
\digestlabel{cor:emitted-decoy}{9.2}\pdfbookmark[2]{Corollary 9.2 The emitted decoy declaration}{digest-transmitter/decoy/cor:emitted-decoy}
\begin{mainclaim}
\renewcommand*{\RestatedName}{\hyperref{contract.pdf}{}{claim-transmitter/decoy/cor:emitted-decoy}{Corollary~9.2}}%
\begin{digestclaim}[The emitted decoy declaration]
Let a weak-coherent-pulse model (Definition~\ref{def:wcp-model}) and a tagged
source declaration over it (Definition~\ref{def:tagged-source}) be presented,
with the certified inequality \eqref{eq:weightineq},
$0\le\tilde p(\mathsf{n}\mid\mu)\le P_\mu(\mathsf{n})$, discharged for every
$\mathsf{n}\le N_{\mathrm{src}}$ and every intensity $\mu\in\mathrm{Int}$. Its
\emph{decoy declaration} is the ensemble presentation, in the sense of
Definition~\ref{def:srcert}, whose declared data are exactly
\begin{itemize}[leftmargin=1.6em]
\item the alphabet $\Xset$ with its declared order and the weights $p_x$;
\item the matched sector structure, matched meaning that Alice's system and
  $\Avirt$ carry one ordered index set: the sectors $S_t$ of
  Definition~\ref{def:tagged-source}\,(T3) that some letter puts positive
  weight on, in the declared order of $\mathbb{T}$, each with its orthonormal
  basis of (T3);
\item the weighted blocks $\sigma^{\mathrm v,(t)}_x=p_xw_{x,t}
  \dyad{\varphi_{x,t}}$ of \eqref{C-eq:decoy-blocks}, as block lists;
\item the specified ensemble $\{p_x,\rho_x\}$ of the model as the source model,
  the source segment with the one tagging link of
  Definition~\ref{def:tagged-source}\,(T6) and its obligation record, and the
  assumption ledger;
\end{itemize}
so that no retained system, no measurement, no marginal and no certificate is
among them. Below, \emph{emitted} refers to this decoy declaration: the emitted
source descriptor, segment and assumption ledger are the ones it
declares. Then
the following hold.
\begin{enumerate}[label=(\roman*),leftmargin=3.0em]
\item \emph{(a certificate, and an admissible specification)} where (E1) to
(E3) of Definition~\ref{def:srcert} hold of the declared blocks, the objects
\eqref{eq:enspres} --- the retained system $A=\bigoplus_tA_t$ with $A_t$ a copy
of $S_t$, the marginal $\tau_A$, the POVM $\{M_x\}$ and the blocks
$(\Psi_t)_t$ --- are a source-replacement certificate in block form
with every non-zero block a multiple of a rank-one projector, satisfying (SR1)
and (SR2) exactly against the tagged virtual ensemble of
Definition~\ref{def:tagged-source}\,(T4); the data are a source specification
in the sense of Definition~\ref{def:source-spec} in
trusted-source mode; and
Theorem~\ref{thm:certificate-sufficiency} holds of it, its part (i) reading (E1) to (E3) as the
decided conditions.
\item \emph{(the segment, and the source identity at the emitted descriptors)}
\begin{enumerate}[label=(\alph*),leftmargin=2.2em]
\item \emph{(the source identity)} The emitted source descriptor is
$(\Avirt,\{\sigma^{\mathrm v}_x\})$ with $\sigma^{\mathrm v}_x=p_x\rvirt_x$ for
the tagged ensemble of Definition~\ref{def:tagged-source}\,(T4) with the
declared weights, hence a member of the declared descriptor family of
Proposition~\ref{prop:tagging}, every tagged source declaration over every
weak-coherent-pulse model; so the source identity \eqref{eq:srcmap}, here
$\Xi_{\mathrm{tag}}(\rvirt_x)=\rho_x$, holds exactly at the emitted
descriptors, as that proposition's conclusion.
\item \emph{(the segment record)} The emitted segment satisfies the
interface conditions (S1) (descriptor match), (S2) (index preservation)
and (S4) (slots) of Definition~\ref{def:source-segment}, and its (S3) (source
identity) is the conclusion of (a). Its list of
transported constraint entries is empty and it carries no licence
record. Its obligation record, the triple a consumer compares with the
soundness registry to admit the link, is
$\big(\textsf{tagging},\ \text{the id of Proposition~\ref{prop:tagging}},\
\{\mathrm{A\textnormal{-}WCP}\}\big)$, together with the member of the declared
descriptor family that its descriptors are
(Definition~\ref{def:source-link}\,(c2)). The third component is the union of
the assumption closure of Proposition~\ref{prop:tagging}, which is empty, and
the physical-source assumption A-WCP that this submodule declares for the model
that proposition is stated over (Definition~\ref{def:source-segment}\,(A1)).
Inside an artefact that this submodule emits and whose assumption ledger lists
A-WCP, the segment meets the admission conditions (A1) to (A3) of
Definition~\ref{def:source-segment}: the registry entry, the decided triple
and the provenance.
\end{enumerate}
\item \emph{(what the certificate reaches)} every state reached from the
certificate of (i) by an attack on the
transmitted system and the adversary's
registers is block diagonal along $\mathbb{T}$ with $A$-marginal $\tau_A$, and
pinching $A$ along $\mathbb{T}$ after the attack is the same operation as
pinching the virtual signal before it (Lemma~\ref{lem:matched-pinching});
Theorem~\ref{thm:certificate-sufficiency}\,(iii) holds against the tagged ensemble of
Definition~\ref{def:tagged-source}\,(T4) itself;
and a positive operator constrained only to have the marginal $\tau_A$ need not
be block diagonal.
\item \emph{(the entries of the declared blocks)} (E1) to (E3) are conditions on
the declared blocks, and every entry of a block is a product
$p_x\,w_{x,t}\,\big(\dyad{\varphi_{x,t}}\big)_{ij}$ of the declared weight
$p_x$, the weight $w_{x,t}$, which is a declared sector weight or a
tail weight formed from them, and an entry of the dyad of $\ket{\varphi_{x,t}}$ in the
orthonormal basis of $S_t$ of Definition~\ref{def:tagged-source}\,(T3), a
number fixed by the model and that basis and not a declared number.
\begin{enumerate}[label=(\alph*),leftmargin=2.2em]
\item \emph{(the weights)} The products $p_xw_{x,t}$ are formed from the
declared $p_x$, the declared $\tilde p(\mathsf{n}\mid\mu)$ and the
tail weight \eqref{eq:tailweight}, which is itself formed exactly and rounds nowhere.
Choosing the $p_x$ and the $\tilde p(\mathsf{n}\mid\mu)$ with few enough
significant bits makes every product exact, and lowering a $\tilde p$ to a
value that is still non-negative is permitted by \eqref{eq:weightineq} and
yields again a tagged source declaration by (v).
\item \emph{(the signal dyads)} The entries of $\dyad{\varphi_{x,t}}$ are the
number $1$ on every flag sector and on the vacuum, and on a quantum sector
$\mathsf{n}\ge1$ the entries of $\dyad{\mathsf{n},a}$ in the basis of
$S_{\mathsf{n}}$. For every unit vector $\ket{b}\in\mathbb{C}^2$, in
$\mathrm{Pol}$ or not, write
$\ket{\mathsf{n},b}\coloneqq(c^{\dagger}_b)^{\mathsf{n}}\ket{\vac}/
\sqrt{\mathsf{n}!}$, which for $b\in\mathrm{Pol}$ is the vector of
\eqref{eq:numberstate}; then
$\big|\braket{\mathsf{n},a}{\mathsf{n},b}\big|^2=|\braket{a}{b}|^{2\mathsf{n}}$
for all unit vectors $a,b$. Let every polarisation be one of the four
vectors of $\Dset$, with $\ket{H},\ket{V}$ the fixed orthonormal basis of the
one-photon space of Definition~\ref{def:wcp-model}\,(W1) and
$\ket{\pm}=(\ket{H}\pm\ket{V})/\sqrt2$, distinct polarisations being distinct
vectors, and put $v_H=(1,1,0)$, $v_V=(1,-1,0)$ and $v_{\pm}=(1,0,\pm1)$. If
$|\mathrm{Pol}|\ge2$, then $S_1$ is the one-photon space,
$\{\ket{1,H},\ket{1,V}\}$ is an orthonormal basis of it, and in that basis
$\dyad{1,H}$, $\dyad{1,V}$ and $\dyad{1,\pm}$ are
$\big(\begin{smallmatrix}1&0\\0&0\end{smallmatrix}\big)$,
$\big(\begin{smallmatrix}0&0\\0&1\end{smallmatrix}\big)$ and
$\tfrac12\big(\begin{smallmatrix}1&\pm1\\\pm1&1\end{smallmatrix}\big)$. If
$|\mathrm{Pol}|\ge3$, then $S_2$ is the two-photon space, and it has an
orthonormal basis in which $\dyad{2,a}=\tfrac12\,v_av_a^{\transpose}$ for every
$a\in\Dset$. So every entry of these dyads is $0$, $1$ or $\pm\tfrac12$ at one
photon and $0$ or $\pm\tfrac12$ at two.
\item \emph{(the sector sums)} With the entries of (a) and (b) representable,
(E3) is a condition on the sums $\sum_x\sigma^{\mathrm v,(t)}_x$, accumulated
entry by entry in the declared order, and on nothing else.
\end{enumerate}
\item \emph{(lowered weights)} if
$0\le\tilde p'(\mathsf{n}\mid\mu)\le\tilde p(\mathsf{n}\mid\mu)$ for every
$\mathsf{n}\le N_{\mathrm{src}}$ and every $\mu$, then the weights $\tilde p'$
are again a tagged source declaration over the same model and
Proposition~\ref{prop:tagging} holds of it, the difference passing into the
tail weight. Its weights, its virtual ensemble, its tail states and its sector
list are those formed from $\tilde p'$: the sector list may gain a tail flag,
where $\tilde w_{\mu}$ was zero and $\tilde w'_{\mu}$ is not, and it may lose a
polarisation flag or a quantum sector, where a weight is lowered to zero at
every letter that carried it, so that neither sector list need contain the
other.
\end{enumerate}
\end{digestclaim}
\end{mainclaim}

\section*{10\quad The physical-source assumption}\pdfbookmark[1]{10 The physical-source assumption}{digest-group-transmitter/decoy/ass:wcp}

% claim transmitter/decoy/ass:wcp
\digestlabel{ass:wcp}{10.1}\pdfbookmark[2]{Assumption 10.1 Emission of the declared weak-coherent-pulse states}{digest-transmitter/decoy/ass:wcp}
\SetShortId{A-WCP}
\renewcommand*{\RestatedName}{\hyperref{contract.pdf}{}{claim-transmitter/decoy/ass:wcp}{Assumption~10.1}}%
\begin{digestassumption}[Emission of the declared weak-coherent-pulse states]
For the tagged source declaration of Definition~\ref{def:tagged-source},
with the cutoffs, the weights and the polarisations it declares, the physical
source satisfies:
(i) in every round it emits the specified state $\rho_x$ of
Definition~\ref{def:wcp-model} of the letter $x=(a,\mu,z)$ drawn in that
round, and nothing else;
(ii) the intensity is the declared $\mu$ exactly, with no fluctuation and no
uncertainty interval;
(iii) the phase randomisation is full, so that the emitted state is the phase
average \eqref{eq:wcpstate};
(iv) the emission is in the two optical modes declared in
Definition~\ref{def:wcp-model}\,(W1), and the polarisation is the declared
$\ket{a}$;
(v) no degree of freedom other than those two modes carries any dependence on
$x$;
(vi) the round type $z$ changes nothing the source emits;
(vii) the letters $x_1,\dots,x_n$ of the $n$ rounds are drawn independently,
each with the weights $p_x$ of Definition~\ref{def:wcp-model}\,(W2);
(viii) given those letters, the systems the source emits in the $n$ rounds
are in the product state $\rho_{x_1}\otimes\dots\otimes\rho_{x_n}$ of the
per-round states;
(ix) neither the letter of a round nor the state emitted in it depends on
what happens in the other rounds, their outcomes and public messages
included.
\end{digestassumption}

\section*{11\quad Emission and refusals}\pdfbookmark[1]{11 Emission and refusals}{digest-group-transmitter/decoy/req:decoy-emission}

% claim transmitter/decoy/req:decoy-emission
\digestlabel{req:decoy-emission}{11.1}\pdfbookmark[2]{Requirement 11.1 Emission of the tagged ensemble as declared}{digest-transmitter/decoy/req:decoy-emission}
\renewcommand*{\RestatedName}{\hyperref{contract.pdf}{}{claim-transmitter/decoy/req:decoy-emission}{Requirement~11.1}}%
\begin{digestrequirement}[Emission of the tagged ensemble as declared]
When the engine emits an artefact from a tagged source declaration over
a weak-coherent-pulse model, the artefact is that declaration's decoy
declaration (Corollary~\ref{cor:emitted-decoy}) and nothing else: the engine
declares no retained system, no measurement, no marginal and no certificate,
and it rounds or otherwise alters no declared number. The matched
sector structure among the declared data of that corollary is emitted as
the retained-register sector map of
Definition~\ref{def:tagged-source}\,(T5), with
the block kind \iface{proof_sector}, the tag that records a decomposition
licensed by a proof, here source replacement, rather than by the device's own
mode structure;
it names Alice's sectors only as copies of the declared sectors $S_t$, so it
declares no retained system. An artefact emitted from a tagged declaration
carries A-WCP and not A-ENS in its assumption ledger
(Assumption~\ref{ass:wcp}).
\end{digestrequirement}

% claim transmitter/decoy/req:cutoff-order
\digestlabel{req:cutoff-order}{11.2}\pdfbookmark[2]{Requirement 11.2 Order of the entropy and source cutoffs}{digest-transmitter/decoy/req:cutoff-order}
\renewcommand*{\RestatedName}{\hyperref{contract.pdf}{}{claim-transmitter/decoy/req:cutoff-order}{Requirement~11.2}}%
\begin{digestrequirement}[Order of the entropy and source cutoffs]
The engine refuses, with no artefact, a declaration whose cutoffs are not
integers with $0\le N_{\mathrm{ent}}\le N_{\mathrm{src}}$, which is (T1) of
Definition~\ref{def:tagged-source} (\iface{cutoffs_out_of_order}, the one
refusal for a cutoff that is not an integer and for cutoffs out of order).
\end{digestrequirement}

% claim transmitter/decoy/req:entropy-cutoff
\digestlabel{req:entropy-cutoff}{11.3}\pdfbookmark[2]{Requirement 11.3 Entropy cutoff of at most two photons}{digest-transmitter/decoy/req:entropy-cutoff}
\renewcommand*{\RestatedName}{\hyperref{contract.pdf}{}{claim-transmitter/decoy/req:entropy-cutoff}{Requirement~11.3}}%
\begin{digestrequirement}[Entropy cutoff of at most two photons]
The engine refuses, with no artefact, a declaration with
$N_{\mathrm{ent}}\ge3$ (\iface{entropy_cutoff_above_two}). The refusal is a
scope decision and not a consequence of anything the claims state: every claim
of this submodule is stated for arbitrary $N_{\mathrm{ent}}\le N_{\mathrm{src}}$
and nothing in them fails at $N_{\mathrm{ent}}\ge3$. What is missing is a
construction of the kind Corollary~\ref{cor:emitted-decoy}\,(iv)(b) supplies
for the four polarisations $\Dset$ at one and at two photons: a declared
basis in which the signal dyads have exactly declarable entries. This
repository has no such construction at three photons and above. Admitting such
a sector would mean deciding exact conditions on dyad entries that differ from
the entries of the signal states, so until the construction exists the sector
is refused rather than declared with entries that are not the numbers written.
\end{digestrequirement}

% claim transmitter/decoy/req:signal-dyads
\digestlabel{req:signal-dyads}{11.4}\pdfbookmark[2]{Requirement 11.4 Signal dyads formed by the exhibited construction}{digest-transmitter/decoy/req:signal-dyads}
\renewcommand*{\RestatedName}{\hyperref{contract.pdf}{}{claim-transmitter/decoy/req:signal-dyads}{Requirement~11.4}}%
\begin{digestrequirement}[Signal dyads formed by the exhibited construction]
The engine forms the entries of the signal dyads of a decoy declaration
itself, and takes neither a dyad nor a basis of a sector from a request. On
the vacuum sector and on every flag sector it forms the number $1$. On a
quantum sector $\mathsf{n}\in\{1,2\}$ of the index set it forms them only where
every polarisation of $\mathrm{Pol}$ is declared as one of the four vectors of
$\Dset$, distinct polarisations as distinct vectors, and
$|\mathrm{Pol}|\ge\mathsf{n}+1$; it then forms the dyad of $\ket{\mathsf{n},a}$
with the entries Corollary~\ref{cor:emitted-decoy}\,(iv)(b) states for it. It
decides of the formed dyads, exactly, that each is a rank-one projector on
$S_{\mathsf{n}}$ (Hermitian, idempotent and of unit trace) and that
$\tr\big[\dyad{\mathsf{n},a}\,\dyad{\mathsf{n},b}\big]
=|\braket{a}{b}|^{2\mathsf{n}}$ for every pair of polarisations of
$\mathrm{Pol}$, the value that Definition~\ref{def:wcp-model}\,(W3) gives by
that corollary and that is $0$, $2^{-\mathsf{n}}$ or $1$ for two vectors of
$\Dset$. It refuses, with no artefact, a declaration whose index set carries
a quantum sector $\mathsf{n}\ge1$ outside this construction, or at which one of
these decisions fails (\iface{quantum_sector_outside_construction}). A
declaration with $N_{\mathrm{ent}}\ge3$ is refused under its own reason
(Requirement~\ref{req:entropy-cutoff}).
\end{digestrequirement}

% claim transmitter/decoy/req:sector-weight-bound
\digestlabel{req:sector-weight-bound}{11.5}\pdfbookmark[2]{Requirement 11.5 Certified bound on the sector weights}{digest-transmitter/decoy/req:sector-weight-bound}
\renewcommand*{\RestatedName}{\hyperref{contract.pdf}{}{claim-transmitter/decoy/req:sector-weight-bound}{Requirement~11.5}}%
\begin{digestrequirement}[Certified bound on the sector weights]
The obligation \eqref{eq:weightineq} of Definition~\ref{def:tagged-source}\,(T2),
$0\le\tilde p(\mathsf{n}\mid\mu)\le P_\mu(\mathsf{n})$ for every
$\mathsf{n}\le N_{\mathrm{src}}$ and every intensity $\mu\in\mathrm{Int}$ of
Definition~\ref{def:wcp-model}\,(W2), is an inequality and never an equality:
its right-hand side is not a declared datum, and at a positive intensity,
declared as a binary64 number, it is a transcendental number (at $\mu=0$ it is
$1$ for $\mathsf{n}=0$ and $0$ for $\mathsf{n}\ge1$). For every such
$(\mathsf{n},\mu)$ the engine forms, from the declared intensity $\mu$ and in
exact rational arithmetic, the lower end
$\underline{P}_\mu(\mathsf{n})\coloneqq\max(0,S_{K+1})\,\mu^{\mathsf{n}}/\mathsf{n}!$
of the enclosure of Lemma~\ref{C-lem:poisson-enclosure}, with $S_{K+1}$ the
partial sum \eqref{C-eq:expsum} at $K=K_\mu$, the least even integer
$K\ge\mu$ with $t_{K+1}=\mu^{K+1}/(K+1)!\le2^{-90}$, which exists because
$t_{K+1}\to0$. It records $K_\mu$ for every intensity in the artefact, and it
reads no enclosure of $P_\mu(\mathsf{n})$ from the request. By part (iv) of
that lemma no even $K$ with $\mu\le K<K_\mu$ gives a larger lower end, so the
rule refuses no weight that a smaller admissible index would admit. It
refuses, with no artefact, a declaration in which
$\tilde p(\mathsf{n}\mid\mu)<0$ or
$\tilde p(\mathsf{n}\mid\mu)>\underline{P}_\mu(\mathsf{n})$ at some
$(\mathsf{n},\mu)$ (\iface{sector_weight_above_poisson}). Since
$\underline{P}_\mu(\mathsf{n})\le P_\mu(\mathsf{n})$ by that lemma, every
declaration violating \eqref{eq:weightineq} is refused, and every declaration
it admits satisfies \eqref{eq:weightineq}. It refuses, with no artefact, a
declaration presenting the obligation as an equality
$\tilde p(\mathsf{n}\mid\mu)=P_\mu(\mathsf{n})$ to be decided
(\iface{sector_weight_as_equality}). No condition of this submodule compares
$\tilde p(\mathsf{n}\mid\mu)$ with $P_{\mu}(\mathsf{n})$ for equality and none
is decided at a tolerance.
\end{digestrequirement}

% claim transmitter/decoy/req:shared-flags
\digestlabel{req:shared-flags}{11.6}\pdfbookmark[2]{Requirement 11.6 Shared polarisation and tail flags}{digest-transmitter/decoy/req:shared-flags}
\renewcommand*{\RestatedName}{\hyperref{contract.pdf}{}{claim-transmitter/decoy/req:shared-flags}{Requirement~11.6}}%
\begin{digestrequirement}[Shared polarisation and tail flags]
The engine refuses, with no artefact, a declaration carrying a flag family finer
than the one Definition~\ref{def:tagged-source}\,(T3) declares
(\iface{flag_family_split_by_intensity}). That family is one
polarisation flag $f_{\mathsf{n},a}$ for each sector $(\mathsf{n},a)$ of
the index set of (T3), $N_{\mathrm{ent}}<\mathsf{n}\le N_{\mathrm{src}}$, shared
by every intensity and both round types, and one
tail flag $f_{a,\mu}$ for each sector $(a,\mu)$ of
that index set, shared by both round types; a finer family is one with separate
flags where (T3) declares one. In particular the engine refuses a family whose
polarisation flags are split by intensity at a photon number
$N_{\mathrm{ent}}<\mathsf{n}\le N_{\mathrm{src}}$. A family split by intensity
still satisfies the source identity but gives up the relation across
intensities that is the decoy content of the declaration, one photon number's
weight at every intensity sitting on one flag
(Remark~\ref{C-rem:no-intensity-split}), and without the refusal that loss would
be silent.
\end{digestrequirement}

% claim transmitter/decoy/req:exact-decoy-blocks
\digestlabel{req:exact-decoy-blocks}{11.7}\pdfbookmark[2]{Requirement 11.7 Binary64 entries of the emitted blocks and sector sums}{digest-transmitter/decoy/req:exact-decoy-blocks}
\renewcommand*{\RestatedName}{\hyperref{contract.pdf}{}{claim-transmitter/decoy/req:exact-decoy-blocks}{Requirement~11.7}}%
\begin{digestrequirement}[Binary64 entries of the emitted blocks and sector sums]
The engine decides (E3) of Definition~\ref{def:srcert} on the sums of the
declared blocks of a decoy declaration, accumulated entry by entry in the
declared order. Each entry of a block, a product of declared numbers, the
tail weight formed from them and a formed entry of a signal dyad
(Corollary~\ref{cor:emitted-decoy}\,(iv)), is
required to be a binary64 number and is refused, with no artefact, when it is
not (\iface{formed_entry_not_declarable}). Those sums are required to round
nowhere, so that their entries are binary64 numbers, and a declaration at a
sector of which the accumulation rounds fails (E3) and is refused, with no
artefact (\iface{ensemble_presentation_not_admissible}). No entry is ever
rounded, because rounding an entry of a declared block would change which
ensemble is declared and so which ensemble the source identity of
Corollary~\ref{cor:emitted-decoy}\,(ii) is asserted of.
\end{digestrequirement}

% claim transmitter/decoy/req:decoy-ensemble-conditions
\digestlabel{req:decoy-ensemble-conditions}{11.8}\pdfbookmark[2]{Requirement 11.8 Positivity and weight sums of the emitted blocks}{digest-transmitter/decoy/req:decoy-ensemble-conditions}
\renewcommand*{\RestatedName}{\hyperref{contract.pdf}{}{claim-transmitter/decoy/req:decoy-ensemble-conditions}{Requirement~11.8}}%
\begin{digestrequirement}[Positivity and weight sums of the emitted blocks]
The engine decides (E1) and (E2) of Definition~\ref{def:srcert} on the declared
blocks of a decoy declaration, as it decides (E3)
(Requirement~\ref{req:exact-decoy-blocks}). It refuses, with no artefact, a
declared block failing (E1) or a weight sum failing (E2)
(\iface{ensemble_presentation_not_admissible}).
\end{digestrequirement}

\bibliographystyle{plain}
\bibliography{../../../refs/references}

\end{document}
